% =============================================================================================================
% Thesis: an explainable, hallucination-aware Swin framework for chest CT and X-ray (CT-supervised bone suppression, lesion absorption and
% the twin regulariser; bone-suppressed classifier inputs; CT-to-X-ray distillation, Chapter ch:crossmodal).
% It shares tables, figures, section files and the bibliography with ../paper.
%
% Build (inside thesis/):   pdflatex main && bibtex main && pdflatex main && pdflatex main
% Rules (same as the paper)
%   1. No result is typed by hand. Result tables are \input from ../paper/tables (tools/final_stats.py, tools/paper_extras.py,
%      tools/bdcxr_stats.py), figures come from ../paper/figures; numbers quoted in the text are copied from those files or from
%      results/final/stats_held.json and results/bdcxr/stats.json.
%   2. Chapters \input paper sections where the text overlaps; each paper section is input at most once (labels stay unique). The paper's
%      results.tex is NOT input: the result chapters define the same labels with thesis-length text.
%   3. Red [TODO: ...] markers are the places where text or a decision is still missing.
% =============================================================================================================
\documentclass[11pt,a4paper,openany]{report}

\usepackage{amsmath}
\usepackage{graphicx}
\usepackage{booktabs}
\usepackage{longtable}
\usepackage{xcolor}
\usepackage{siunitx}
\usepackage{setspace}
\usepackage{adjustbox}
\usepackage{microtype}
\usepackage[section]{placeins}
\usepackage[authoryear,round]{natbib}
\usepackage[colorlinks=true,linkcolor=blue!50!black,citecolor=blue!50!black,urlcolor=blue!50!black]{hyperref}
\onehalfspacing

% ---- layout (same as paper/main.tex) ---------------------------------------------------------------------------
\setlength{\emergencystretch}{3em}   % justified lines stretch instead of running into the margin
\def\UrlBreaks{\do\/\do\.\do\_\do\-}  % file paths in \path{...} may break after / . _ -
\renewcommand{\topfraction}{0.9}
\renewcommand{\bottomfraction}{0.8}
\renewcommand{\textfraction}{0.07}
\renewcommand{\floatpagefraction}{0.75}
\setcounter{topnumber}{3}
\setcounter{bottomnumber}{2}
\setcounter{totalnumber}{5}

% Look for \input files, tables and figures in the paper folder as well as in thesis/
\makeatletter
\def\input@path{{../paper/}}
\makeatother
\graphicspath{{../paper/}}

% ---- helpers (keep identical to paper/main.tex) ---------------------------------------------------------------
% Red to-do marker. The argument is detokenized, so file names with underscores are safe; use plain ASCII and no
% macros, no percent sign and no empty line inside \todo{...}.
\newcommand{\todo}[1]{\textcolor{red}{[TODO: \detokenize{#1}]}}
% Generated table, shrunk to the text width or 80 % of the page height when larger (never enlarged): \fittable{tables/name.tex}
\newcommand{\fittable}[1]{\begin{adjustbox}{max width=\linewidth,max totalheight=0.8\textheight}\input{#1}\end{adjustbox}}
% Figure that falls back to a red box when the PDF has not been generated yet: \safefig[width]{figures/name.pdf}
\newcommand{\safefig}[2][\linewidth]{%
  \IfFileExists{#2}%
    {\includegraphics[width=#1]{#2}}%
    {\fbox{\parbox[c][2.4cm][c]{0.92\linewidth}{\centering\todo{missing figure #2}}}}}
% some bibliography styles print DOI fields through \doi{...}
\providecommand{\doi}[1]{\url{https://doi.org/#1}}
% Paper sections shared by the two papers and the thesis name the other paper; here they point to chapters instead
% (paper/main.tex and paper/main_cls.tex define the paper wording).
\newcommand{\companion}[1]{Chapter~\ref{#1}}
\newcommand{\companioncite}{}
\newcommand{\clsstudy}{the pre-specified study of Chapter~\ref{ch:classifier}}

\begin{document}

\begin{titlepage}
\centering
\vspace*{2cm}
{\LARGE\bfseries Development of an Explainable and Hallucination-Aware Swin Transformer Based Framework for Cross Modality Chest CT and
X-Ray Disease Classification\par}
\vspace{0.8em}
{\Large Lesion-faithful CT-supervised bone suppression, a counterfactual lesion audit and CT-to-X-ray knowledge transfer\par}
\vspace{3em}
{\large A thesis submitted in partial fulfilment of the requirements for the degree of\par}
\vspace{0.5em}
{\large \todo{degree and programme}\par}
\vspace{3em}
{\large by\par}
\vspace{0.5em}
{\Large Tanvir Hossin\par}
{\large \todo{student ID}\par}
\vspace{3em}
{\large Supervisor: \todo{name, title}\par}
\vfill
{\large \todo{department}\par}
{\large \todo{university}\par}
\vspace{1em}
{\large \todo{month and year of submission}\par}
\end{titlepage}

\pagenumbering{roman}

\chapter*{Declaration}
\addcontentsline{toc}{chapter}{Declaration}
I declare that this thesis is my own work and has not been submitted for any other degree or qualification. Work by others is acknowledged
by citation, and the data sets are used under the terms stated in Chapter~\ref{ch:ethics}.
\todo{department wording of the declaration, signature and date; AI-assistance statement: an AI assistant (Claude, Anthropic) was used for code, analysis automation and drafting, and the author reviewed its output -- adapt to the department rules}

\chapter*{Abstract}
\addcontentsline{toc}{chapter}{Abstract}
Chest CT shows in three dimensions what a chest radiograph superimposes on one plane. This thesis develops two kinds of model in which CT
supervises the reading of radiographs, bone suppressors and Swin transformer classifiers, and asks of each whether it can be trusted:
whether it keeps what it should keep, whether it reports findings that are not there, and whether its explanations point at the disease.

The first part concerns bone suppressors trained on projections rendered from CT, where the bone image is known exactly. Such a network may
move part of a lesion into the bone image and so delete it from the suppressed radiograph, which fidelity metrics do not show. We introduce
a counterfactual lesion audit, which adds nodules from CT scans unseen in training to real radiographs at fixed sites and
measures the absorption $A$, the fraction of the added signal that moves into the predicted bone image, and a twin lesion-invariance
regulariser, which ties the bone predictions of a training image and of a copy with a pasted lesion. Two studies with pre-specified plans
trained suppressors on LIDC-IDRI renders and on the public Qure.ai projections, and every change to the plans is reported. The regulariser
lowered absorption in a dose-dependent way. On VinDr-RibCXR, unseen during development, the final three-model ensemble
absorbed about a third of the lesion signal that a released CT-supervised model absorbed and less than a third of what a matched control
without the regulariser absorbed, with non-inferior bone-region fidelity on held-out dual-energy pairs; the effect replicated with three
encoders and on 3,257 paediatric radiographs from Bangladesh. Not everything held: structural similarity to the dual-energy image missed its
margin, real-nodule contrast rose more than with the reference model on one data set and less on another, the suppressed image alone did
not improve downstream classification or detection, and intervals calibrated on synthetic data did not cover real errors until a post hoc
recalibration on 24 real pairs. A third pre-specified study, 175 classifier runs on three data sets, found no gain in overall accuracy from
the soft-tissue and bone images against a control of the same capacity, only a small gain for nodules and masses, significant on NIH
ChestX-ray14 alone.

The second part builds an explainable, hallucination-aware Swin framework for cross-modality chest CT and X-ray classification. A 2.5D
SwinV2 teacher learns six findings from CT-RATE volumes, an X-ray model learns them from projections of the same volumes with and without
distillation from the teacher, and both initialise NIH ChestX-ray14 classifiers. Under a fourth pre-specified plan, distillation did not
improve the classification of real radiographs (difference in mean AUROC $-0.0001$, 95\,\% interval $-0.0016$ to $0.0014$), and the CT and
X-ray heat maps of the same patient did not agree. Grad-CAM++ maps of the Swin classifiers located expert boxes less than half as often as
those of a DenseNet-121, whereas Integrated Gradients and decoder attention localised about as well. The Swin classifiers were far less
sensitive to image noise than DenseNet-121 (27--29\,\% against 78\,\% of finding-free images flipped), and rib emphasis
flipped slightly fewer images after distillation (12.9\,\% against 13.7\,\%, $p=0.044$, uncorrected), as it did after pretraining on the
projections without the teacher.

Across both parts, models of equal accuracy or fidelity differed widely in what they erased, invented or highlighted. Lesion absorption,
hallucination and localisation should therefore be measured before CT-supervised models are used on radiographs.

\chapter*{Acknowledgements}
\addcontentsline{toc}{chapter}{Acknowledgements}
\todo{personal acknowledgements: supervisor, lab, family, funding}
This work used public and licensed data provided by Qure.ai, the CT-RATE team, The Cancer Imaging Archive and the LIDC-IDRI investigators, the
authors of SZCH-X-Rays, VinDr (VinDr-RibCXR and VinDr-CXR), the Japanese Society of Radiological Technology, the Node21, TBX11K and NIH
ChestX-ray14 teams, and Dr.~M~R~Khan Shishu Hospital \& Institute of Child Health, Dhaka, which provided BDCXR-3257 under the ethical approval
REC-FSIT-2025/No:~11928. Training ran on Kaggle GPU quota.

\tableofcontents
\listoffigures
\listoftables

\chapter*{Abbreviations and symbols}
\addcontentsline{toc}{chapter}{Abbreviations and symbols}
\begin{longtable}{@{}lp{0.74\linewidth}@{}}
$A$ & lesion absorption (Chapter~\ref{ch:audit}) \\
AUPRC & area under the precision-recall curve \\
AURC & area under the risk-coverage curve (selective classification) \\
AUROC & area under the receiver operating characteristic curve \\
CBAM & convolutional block attention module \\
CLIA & counterfactual lesion insertion augmentation (lesions pasted into training images) \\
CNR & contrast-to-noise ratio \\
CPM & competition performance metric (mean FROC sensitivity) \\
DES & dual-energy subtraction \\
EMA & exponential moving average of the weights \\
FNR & false-negative rate \\
F$_\mathrm{in}$, F$_\mathrm{p0}$, F$_\mathrm{p1}$ & NIH classifiers initialised from ImageNet, P0 and P1 (Chapter~\ref{ch:crossmodal}) \\
FROC & free-response receiver operating characteristic \\
H1$'$--H5$'$ & confirmatory hypotheses of the second analysis plan \\
H6--H10 & hypotheses of the third analysis plan (Chapter~\ref{ch:classifier}) \\
H11--H15 & hypotheses of the fourth analysis plan (Chapter~\ref{ch:crossmodal}) \\
IG & Integrated Gradients (attribution method) \\
$K$ & number of ensemble members \\
LPIPS & learned perceptual image patch similarity \\
P0, P1 & projection models without and with distillation from the CT teacher (Chapter~\ref{ch:crossmodal}) \\
PSNR, PSNR$_{cal}$ & peak signal-to-noise ratio, after a per-image gain and offset fit \\
QP & the Qure.ai CT projections (training source of the second study) \\
SSIM & structural similarity \\
$w$ & twin weight of the lesion-invariance regulariser \\
$x$, $b$, $s$ & radiograph, bone image, soft-tissue image ($s=x-b$) \\
\end{longtable}

\cleardoublepage
\pagenumbering{arabic}

\chapter{Introduction}\label{ch:intro}
% Thesis introduction. Longer than paper/sections/introduction.tex, which is not input here. No result numbers except in the summary of answers,
% which are copied from the generated tables.
\section{Motivation}
This thesis develops an explainable and hallucination-aware framework in which chest CT supervises models that read chest radiographs. It
has two parts. The first, which takes most of the thesis, trains bone suppressors on projections rendered from CT and asks whether they
erase lesions; it also asks whether classifiers gain from the suppressed images. The second uses CT to teach a Swin transformer classifier
of radiographs directly, through a CT classifier and projections of the same patients, and asks whether that classifier becomes more
accurate, whether its explanations agree with those of the CT model, and whether it reports fewer findings that are not there. Both parts
ask the same question of a CT-supervised model: what it keeps, what it invents and where it looks.

A chest radiograph projects the whole thorax onto one plane. Ribs, clavicles and the spine lie over most of the lung fields, and a small nodule
or an early consolidation that sits behind a rib has to be seen through the bone shadow. Software that suppresses rib and clavicle shadows has
been shown to help readers find lung nodules \citep{freedman2011lung}, and computer-aided detection has been evaluated on such images as well
\citep{schalekamp2014cad}.

Dual-energy subtraction (DES) separates bone and soft tissue physically: two exposures at different tube voltages are combined so that either
material cancels \citep{kuhlman2006des}. It needs dedicated hardware and a second exposure, and it is rare outside larger centres. Learning-based
bone suppression produces a comparable soft-tissue image from a single standard radiograph. Two families exist today. One trains generative
models on real dual-energy pairs \citep{chen2024bsdiff,sun2026bsldm,sun2025gllcm,sun2026debonedit}; it needs paired data from the target
hospital. The other renders paired supervision from chest CT: each scan is split into bone and soft tissue, the parts are projected separately,
and a network learns to predict one component from their sum \citep{angaitkar2026anatomy}. The second family can be trained on public data,
and the bone image of every training sample is known exactly.

CT can also supervise a radiograph classifier more directly. Public data sets now pair chest CT volumes with report-derived findings
\citep{hamamci2026ctrate}, and the projections of those volumes \citep{angaitkar2026anatomy} pair each scan with a radiograph-like image of
the same patient. A classifier trained on CT sees each finding without the overlap of ribs, heart and mediastinum, and knowledge
distillation \citep{hinton2015distilling} offers a way to pass what it learned to an X-ray model through the paired images. Swin
transformers \citep{liu2022swinv2} can serve both modalities with one architecture, which makes the transfer and the comparison of their
heat maps straightforward.

\section{Problem statement}
A suppressor is usually judged by how closely its output matches a reference image: PSNR and SSIM against a dual-energy soft-tissue image, or
the error of the bone image on rendered data. These measures reward a correct bone estimate on average over the image. They do not ask what
happens to a lesion. A network that predicts the bone image is free to assign part of a nodule's signal to that image, because the nodule is
small, dense and, behind a rib, hard to tell from bone. The suppressed image $s=x-f(x)$ then shows less of the lesion than the radiograph did,
which defeats the purpose of suppression. Fidelity metrics hardly change, because the lesion covers a few hundred pixels of an image of a
million.

Measuring this needs a lesion whose signal is known, and real radiographs do not come with one. A second difficulty is that the safest model
by this measure is one that does nothing: the identity removes no bone and no lesion. Any measure of lesion loss must therefore be read next to
a measure of how much bone the model removes. Figure~\ref{fig:audit_maps} (Chapter~\ref{ch:audit}) shows the problem on one radiograph: the
same added nodule is moved into the bone image to very different degrees by three suppressors.

Classifiers raise the mirror-image problem. A classifier can report a finding the image does not contain, or miss one it does, and its
accuracy on a test set does not say which images or which changes provoke such errors. Heat maps are meant to show where a classifier looks,
but they are often not checked against expert annotations, and the common gradient-based maps were designed for convolutional networks
rather than transformers \citep{chattopadhay2018gradcampp}. When CT is used to teach an X-ray classifier, three things therefore need testing besides accuracy: whether the transfer helps
on real radiographs at all, whether the two models attend to the same place in the same patient, and whether the classifier hallucinates
less when the image changes in ways that add anatomy but no disease.

\section{Research questions}
The thesis answers six research questions; each is tied to pre-specified hypotheses of the four analysis plans (Chapters~\ref{ch:v1},
\ref{ch:v2}, \ref{ch:classifier} and~\ref{ch:crossmodal}).
\begin{description}
\item[RQ1] Is supervision with exact CT components enough to keep lesions, and how much of a lesion does a model trained on rendered LIDC-IDRI
 data absorb? (first plan, H1; Chapter~\ref{ch:v1})
\item[RQ2] Does a lesion-invariance (twin) regulariser lower absorption without losing fidelity on real radiographs? (second plan, H1$'$, H2$'$
 and H3$'$; Chapters~\ref{ch:v2}, \ref{ch:audit} and~\ref{ch:benchmark})
\item[RQ3] Does the counterfactual audit measure something real, that is, does evidence that does not use lesion insertion agree with it?
 (H4$'$a; Chapter~\ref{ch:benchmark})
\item[RQ4] Does it matter for downstream tasks, and how reliable is the output: does the ensemble spread track the error, and do conformal
 intervals cover it, also on a population no model was trained on? (H4$'$b and H5$'$; Chapters~\ref{ch:benchmark}, \ref{ch:reliability}
 and~\ref{ch:external})
\item[RQ5] Do classifiers gain from the suppressor's soft-tissue and bone images when they read them as separate streams, compared with a
 control of the same capacity, and for which findings? (third plan, H6--H8 and H10; Chapter~\ref{ch:classifier})
\item[RQ6] Does knowledge distilled from a Swin CT classifier into a Swin X-ray classifier improve the classification of real radiographs,
 do the CT and X-ray models attend to the same place in the same patient, and does the transfer make the classifier hallucinate less?
 (fourth plan, H11--H15; Chapter~\ref{ch:crossmodal})
\end{description}

\section{Contributions}
\begin{enumerate}
\item A counterfactual lesion audit that measures the fraction of an added lesion a suppressor moves into its bone image, on real radiographs,
 at sites shared by all models, with bone overlap taken from independent annotations, and reported beside suppression strength (first paper and
 thesis).
\item A twin lesion-invariance regulariser for CT-supervised suppressors, with a dose-response over its weight and an ablation of its parts
 (first paper and thesis).
\item A training recipe on public data only (rendered LIDC-IDRI CTs in the first study, the Qure.ai CT projections in the second), with a
 three-model ensemble, split-conformal intervals and a screen of modern encoders (first paper and thesis; the v1 study, the renderer and the
 physics are described in full only here).
\item Two pre-specified evaluations on locked sets, with downstream classification and detection, an external validation on paediatric
 radiographs, and every change to the analysis disclosed, including negative results (first paper and thesis; the negative results and the
 deviations log are thesis appendices).
\item A pre-specified test of whether classifiers gain from the soft-tissue and bone images, used alone, stacked or as streams of a multi-stream
 network, against a control of the same capacity, with an analysis by finding group, stream-wise attribution maps scored against expert boxes
 and conformal uncertainty (second paper and thesis).
\item An explainable, hallucination-aware Swin transformer framework for cross-modality chest CT and X-ray classification: a 2.5D Swin CT
 classifier on CT-RATE, its distillation into an X-ray classifier through paired projections, a test of heat-map agreement between the
 modalities, counterfactual hallucination tests and conformal false-negative control, evaluated under a fourth pre-specified plan against a
 DenseNet-121 reference (thesis).
\end{enumerate}

\section{Outline of the thesis}
Chapter~\ref{ch:background} introduces chest radiography, dual-energy subtraction, learning-based suppression, the architectures, metrics,
uncertainty methods and statistics used later. Chapter~\ref{ch:data} describes the data sets and the CT-to-radiograph physics of the first
study. Chapter~\ref{ch:v1} reports the first, pre-specified study on rendered LIDC-IDRI data and the real-image check that led to the second
plan. Chapter~\ref{ch:v2} presents the Qure.ai projections as an exactly additive training source, the twin regulariser, the exploratory
screening that fixed its operating point, the dose-response and the encoder screen. Chapter~\ref{ch:audit} defines and validates the audit and
reports the confirmatory audit. Chapter~\ref{ch:benchmark} gives the experimental design and the locked-set results for fidelity, rib
suppression, real nodules and downstream tasks. Chapter~\ref{ch:reliability} reports the ensemble uncertainty, and Chapter~\ref{ch:external}
the external validation on BDCXR-3257. Chapter~\ref{ch:classifier} tests whether classifiers gain from the suppressor's images (third plan),
and Chapter~\ref{ch:crossmodal} builds and tests the cross-modality Swin framework (fourth plan). Chapter~\ref{ch:discussion} discusses the
findings of the suppression studies and their limits, Chapter~\ref{ch:ethics} covers reproducibility, licences and ethics, and
Chapter~\ref{ch:conclusion} answers the research questions. Appendix~\ref{app:deviations} lists the deviations from the analysis plans,
Appendix~\ref{app:negative} the negative and null results, Appendix~\ref{app:extra} additional tables, and Appendix~\ref{app:classifier}
the deviations and additional figures of the classifier study. Two papers accompany this thesis: the first reports the second study, the
second the classifier study of Chapter~\ref{ch:classifier}. The v1 study, the physics details, the screening history, the cross-modality
study, the negative results and the ethics chapter are given in full only here.

\section{Notation}
\begin{center}
\begin{tabular}{@{}lp{0.68\linewidth}@{}}
\toprule
Symbol & Meaning \\
\midrule
$x$ & radiograph (or training composite), intensities in $[0,1]$ \\
$b$, $s=x-b$ & bone image and soft-tissue image \\
$f$, $f_k$ & suppressor (predicts the bone image), member $k$ of an ensemble of $K$ \\
$\bar f$, $\sigma$ & ensemble mean and per-pixel standard deviation (spread) \\
$n$, $\alpha$ & lesion signal added to an image, its amplitude \\
$A$ & absorption: fraction of $\alpha n$ moved into the bone estimate (Equation~\ref{eq:absorption}) \\
$w$ & weight of the twin lesion-invariance term (Equation~\ref{eq:twin}) \\
$M$ & lesion support used by the twin term \\
$X$, $S$, $\hat B$ & classifier inputs in Chapter~\ref{ch:classifier}: radiograph, predicted soft-tissue and bone image \\
P0, P1 & projection models of Chapter~\ref{ch:crossmodal}, without and with distillation from the CT teacher \\
F$_\mathrm{in}$, F$_\mathrm{p0}$, F$_\mathrm{p1}$ & NIH classifiers of Chapter~\ref{ch:crossmodal}, initialised from ImageNet, P0 and P1 \\
\bottomrule
\end{tabular}
\end{center}


\chapter{Background}\label{ch:background}
% PLAN (15-20 pages). Tutorial background first, then the related work of the paper, which is input unchanged at the end.
%   2.1 Chest radiography and bone overlap: image formation, attenuation, why ribs and clavicles hide lesions.
%   2.2 Dual-energy subtraction: principle, hardware, dose, SZCH-X-Rays as a public DES data set (sun2026bsldm, sun2026debonedit).
%   2.3 Learning-based bone suppression: MTANN (suzuki2006mtann), CNNs (yang2017cascade, gusarev2017deep), diffusion (chen2024bsdiff, sun2026bsldm,
%       sun2025gllcm, sun2026debonedit), CT-supervised (angaitkar2026anatomy).
%   2.4 Architectures used here: Swin Transformer V2 (liu2022swinv2), UNet++ (zhou2018unetpp), CBAM (woo2018cbam); optimiser (loshchilov2019adamw).
%   2.5 Image-quality and detection metrics: PSNR, SSIM (wang2004ssim), LPIPS (zhang2018lpips), CNR, FROC and CPM (niemeijer2011cpm), AUROC and
%       DeLong (delong1988auc).
%   2.6 Uncertainty: deep ensembles (lakshminarayanan2017ensembles), split conformal prediction (lei2018conformal, angelopoulos2023conformal).
%   2.7 Statistics: paired bootstrap, cluster bootstrap, Holm correction, non-inferiority, pre-specified analysis plans.
%   2.8 Related work as in the paper (input below).
% Section 2.3 describes each cited paper only as far as its abstract (PubMed or arXiv, read 2026-11-03) supports; check against the PDFs before submission.
This chapter collects the background that the later chapters assume. It is written for a reader who knows machine learning but not the physics
of chest radiography or the conventions of this field. Sections~\ref{sec:bg_cxr} and~\ref{sec:bg_des} explain how a chest radiograph is formed
and why bone overlap is a problem; Section~\ref{sec:bg_learned} surveys learned bone suppression; Section~\ref{sec:bg_arch} introduces the network
components and training techniques used in this study; Sections~\ref{sec:bg_metrics} to~\ref{sec:bg_stats} define the metrics, the uncertainty
estimates and the statistical procedures. The related work of the paper closes the chapter (Section~\ref{sec:related}). No result of this study
appears here.

\section{Chest radiography and bone overlap}
\label{sec:bg_cxr}
\paragraph{Projection imaging.} A chest radiograph is a projection image. An X-ray tube emits a beam that passes through the thorax and is
recorded by a detector on the other side. In the standard posteroanterior (PA) view the patient faces the detector and the beam enters from the
back; the anteroposterior (AP) view, used at the bedside, reverses the direction. The tube produces a polyenergetic spectrum: bremsstrahlung up to
a maximum photon energy set by the tube voltage (kVp), plus characteristic lines of the anode material, shaped by the filtration in the beam. Each
detector pixel records the energy that reaches it along the rays that end in that pixel. Everything the ray crosses, from the skin of the back to
the skin of the chest, contributes to the same pixel; depth information is lost. This superposition is the root of the problem studied in this
thesis.

\paragraph{Attenuation.} For a narrow monoenergetic beam of energy $E$ crossing a homogeneous slab of thickness $t$, the transmitted intensity
follows the Beer--Lambert law, $I=I_0\,e^{-\mu(E)\,t}$, where $\mu$ is the linear attenuation coefficient of the material. Through a
heterogeneous body the exponent becomes a line integral along the ray $L$:
\begin{equation}
\begin{gathered}
I(E)=I_0(E)\,\exp\!\Big(-\int_L \mu(\mathbf{r},E)\,\mathrm{d}r\Big),\\
p(E)=-\ln\frac{I(E)}{I_0(E)}=\int_L \mu(\mathbf{r},E)\,\mathrm{d}r .
\end{gathered}
\label{eq:bg_beerlambert}
\end{equation}
The log-transformed value $p$ is linear in the attenuation along the ray, so overlapping structures \emph{add}. It is convenient to write
$\mu=(\mu/\rho)\,\rho$, the product of the mass attenuation coefficient $\mu/\rho$, which depends only on the elemental composition and the
energy, and the density $\rho$. If the body is described by a few materials $c$ (for example bone, soft tissue and lung) with fixed compositions,
the line integral splits into a sum over materials,
\begin{equation}
\begin{gathered}
p(E)=\sum_c \Big(\frac{\mu}{\rho}\Big)_{\!c}(E)\;a_c,\\
a_c=\int_L \rho_c(\mathbf{r})\,\mathrm{d}r ,
\end{gathered}
\label{eq:bg_materials}
\end{equation}
where $a_c$ is the areal density (g\,cm$^{-2}$) of material $c$ along the ray. Tabulated mass attenuation coefficients for elements and mixtures
are available from NIST XCOM \citep{berger2010xcom}, and reference compositions of body tissues from ICRU Report~44 \citep{icru1989report44}.
Equation~\eqref{eq:bg_materials} is the basis of everything that follows: of dual-energy subtraction (Section~\ref{sec:bg_des}), of the
CT-derived training data of this study, and of the idea that a radiograph can be written as a soft-tissue image plus a bone image.

\paragraph{Polyenergetic beams and beam hardening.} A real beam has a spectrum $\phi(E)$, and an energy-integrating detector weights each photon
roughly by its energy. The recorded log signal is then
\begin{equation}
p=-\ln\frac{\int \phi(E)\,E\,\exp\!\big(-\sum_c (\mu/\rho)_c(E)\,a_c\big)\,\mathrm{d}E}{\int \phi(E)\,E\,\mathrm{d}E},
\label{eq:bg_poly}
\end{equation}
the model used in Chapter~\ref{ch:data}. Because low-energy photons are absorbed preferentially, the mean energy of the beam rises along its path
(\emph{beam hardening}) and $p$ is no longer exactly linear in the $a_c$: the sum of separately projected materials differs slightly from the
projection of the whole body. Chapter~\ref{ch:data} quantifies this additivity gap for the renderer of this study.

\paragraph{Why bone stands out.} At diagnostic energies (roughly 20 to 150~keV) two interactions dominate attenuation in tissue. Photoelectric
absorption rises steeply with atomic number and falls steeply with energy; Compton scattering depends mainly on electron density and changes
slowly with energy \citep{bushberg2020essential}.
Bone mineral contains calcium and phosphorus, whose atomic numbers are much higher than those of the hydrogen, carbon, nitrogen and oxygen that make
up soft tissue, so bone attenuates more than soft tissue at every energy and its attenuation falls faster with energy. The first property makes
ribs and clavicles bright on a radiograph; the second makes dual-energy subtraction possible.

\paragraph{From detector signal to stored image.} Scattered photons reach the detector from directions other than the primary ray and add a
smooth haze that lowers contrast; anti-scatter grids reduce but do not remove it. The number of detected photons is random, which adds quantum
noise. Before an image is stored, digital systems apply vendor-specific processing (logarithmic or other tone curves, look-up tables, multi-scale
contrast enhancement and edge sharpening) \citep{prokop2003processing,schaeferprokop2008digital}. A stored radiograph is therefore
an unknown, mostly monotone function of the line integral in Equation~\eqref{eq:bg_poly}, combined with local filtering. This is why the training
data of this study pass the exact components through random monotone tone curves (Section~\ref{sec:methods}) and why real-image metrics are also
reported after a per-image tone calibration (Section~\ref{sec:bg_metrics}).

\paragraph{Why ribs and clavicles hide lesions.} On a PA radiograph the posterior and anterior ribs, the clavicles and parts of the scapulae
project over the lung fields, and the anterior and posterior ribs cross each other. A pulmonary nodule is a small region of near-water density
inside the air-filled lung; in the log domain it adds a compact signal of a few millimetres to a few centimetres. A rib adds a signal of similar or
larger amplitude, with sharp cortical edges, at the same spatial scale. In a linear projection the nodule over a rib keeps its own signal, but it
sits on a background that is bright, structured and full of edges, and the eye (or a detection algorithm) has to separate the two. Background
structure of this kind, often called anatomical noise, is generally regarded as a stronger limit on nodule detection in chest radiography than
quantum noise \citep{samei1999detection}. Reader studies support the clinical relevance
of bone overlap: in a study in which fifteen radiologists read the radiographs of 368 patients first unaided and then with software that
suppresses ribs and clavicles, filters noise and equalises contrast, reader sensitivity rose from 49.5\,\% to 66.3\,\% while specificity fell from
96.1\,\% to 91.8\,\%, and 74\,\% of the aided detections were of cancers with at least 70\,\% overlap between bone and nodule \citep{freedman2011lung}.

\paragraph{Bone suppression as source separation.} In the log domain a radiograph can be written as $x=s+b$, a soft-tissue image $s$ plus a bone
image $b$ (up to the beam-hardening gap and the tone curve). Bone suppression estimates $b$, or equivalently $s=x-b$, from $x$ alone. This is a
separation problem with one equation and two unknowns per pixel; it can only be solved with prior knowledge of what bone looks like (elongated,
smoothly curved ribs with bright cortical edges, clavicles of known shape and position) and of what soft tissue looks like. The prior is weakest
exactly where it matters: a compact bright nodule that lies on a rib is consistent with both explanations. A suppressor that has learned that
bright compact structures in the rib region belong to bone can move part of the nodule into $\hat b$ and remove it from the soft image. The audit
of Chapter~\ref{ch:audit} measures this lesion absorption directly.

\section{Dual-energy subtraction}
\label{sec:bg_des}
\paragraph{Principle.} Dual-energy subtraction (DES) separates bone from soft tissue by imaging the same anatomy with two different spectra.
Write $p_L$ and $p_H$ for the log images obtained with a low-energy and a high-energy spectrum, and treat the body as two materials, soft tissue
($\mathrm{s}$) and bone ($\mathrm{b}$), with effective mass attenuation coefficients $\mu^{L}_{\cdot}$ and $\mu^{H}_{\cdot}$ for the two spectra.
Equation~\eqref{eq:bg_materials} then gives two equations in two unknowns:
\begin{equation}
\begin{gathered}
p_L=\mu^{L}_{\mathrm{s}}\,a_{\mathrm{s}}+\mu^{L}_{\mathrm{b}}\,a_{\mathrm{b}},\\
p_H=\mu^{H}_{\mathrm{s}}\,a_{\mathrm{s}}+\mu^{H}_{\mathrm{b}}\,a_{\mathrm{b}} .
\end{gathered}
\label{eq:bg_des_system}
\end{equation}
A weighted subtraction of the log images cancels one material:
\begin{align}
p_{\mathrm{soft}} &= p_H-w_{\mathrm{b}}\,p_L=\big(\mu^{H}_{\mathrm{s}}-w_{\mathrm{b}}\,\mu^{L}_{\mathrm{s}}\big)\,a_{\mathrm{s}},
 & w_{\mathrm{b}}&=\mu^{H}_{\mathrm{b}}/\mu^{L}_{\mathrm{b}}, \label{eq:bg_des_soft}\\
p_{\mathrm{bone}} &= p_H-w_{\mathrm{s}}\,p_L=\big(\mu^{H}_{\mathrm{b}}-w_{\mathrm{s}}\,\mu^{L}_{\mathrm{b}}\big)\,a_{\mathrm{b}},
 & w_{\mathrm{s}}&=\mu^{H}_{\mathrm{s}}/\mu^{L}_{\mathrm{s}}. \label{eq:bg_des_bone}
\end{align}
The system is solvable only if the ratios $w_{\mathrm{b}}$ and $w_{\mathrm{s}}$ differ, that is, if bone and soft tissue change their attenuation
differently with energy; the photoelectric effect in bone mineral provides this difference (Section~\ref{sec:bg_cxr}). The soft-tissue image
$p_{\mathrm{soft}}$ shows lung, vessels, heart and soft-tissue lesions without ribs and clavicles; the bone image $p_{\mathrm{bone}}$ shows the
skeleton and anything else that contains calcium.

\paragraph{What DES does and does not do.} The decomposition is a two-material model. A third material with a different energy dependence is
split between the two images according to its own coefficients, and calcium outside the skeleton is, by construction, bone. A calcified nodule
therefore appears in the DES bone image, which is one of the main clinical uses of the technique: a review of DES chest radiography names the
clearer depiction of calcification, which helps to characterise pulmonary nodules, as its main advantage and lists further uses in the hila,
mediastinum, airways, vessels, pleura and chest wall \citep{kuhlman2006des}. A non-calcified nodule, by contrast, belongs to the soft-tissue
image. This distinction defines what a learned suppressor should do with a soft-tissue lesion over a rib: leave it in $s$. In practice the
weights in Equations~\eqref{eq:bg_des_soft} and~\eqref{eq:bg_des_bone} are tuned empirically, because the spectra are polyenergetic and beam
hardening and scatter make the effective coefficients depend on thickness; cancellation is imperfect and residual bone edges can remain. The
subtraction of two images also removes most of the signal while their noise adds, so DES images are noisier than the standard radiograph and are
usually filtered, which changes their texture and tone.

\paragraph{Hardware, dose and motion.} Two acquisition schemes are in use \citep{schaeferprokop2008digital}. In \emph{dual-exposure} DES the tube voltage is switched between two exposures taken in quick succession; the energy separation is
good, but anything that moves between the two exposures (the heart, the diaphragm, the patient) is misregistered and leaves artefacts in both
images \citep{schaeferprokop2008digital}. In \emph{single-exposure} DES a layered detector with a filter between the layers
records the softer part of the spectrum in front and the harder part behind; there is no motion between the images, but the energy separation is
smaller and the images are noisier. Dual-exposure radiographs exist in research collections; for example, one CNN study trained on 504 cases of
two-exposure DES chest radiographs \citep{yang2017cascade}. DES needs dedicated equipment, and the review cited above states that it requires a
somewhat higher radiation dose than conventional radiography \citep{kuhlman2006des}; authors of learned methods cite equipment cost and
additional exposure as reasons to replace DES by software \citep{chen2024bsdiff}. Most chest radiographs, including all large public
collections used in this study, are single-energy images.

\paragraph{Public paired data.} Learning bone suppression from real images requires radiographs paired with DES images of the same exposure, and
such pairs are scarce. The data sets of the earlier learned methods (Section~\ref{sec:bg_learned}) are not described as public in their
abstracts. The SZCH-X-Rays data set was compiled for BS-LDM, whose abstract describes 818 pairs of high-resolution radiographs and soft-tissue
images from a partner hospital \citep{sun2026bsldm}; it was released with DeBoneDiT, whose abstract describes the release as 741 pairs of
radiographs and DES soft-tissue images and as the largest publicly available high-quality paired data set for bone suppression \citep{sun2026debonedit}.
The BS-LDM paper compiled 831 pairs and excluded 13; DeBoneDiT and GL-LCM report 741 pairs split into 592 training, 74 validation and 75
test pairs, with the same exclusion criteria. None of the three papers explains the difference, and only the 75 test pairs are public. The
images come from one partner hospital, acquired with dual-exposure DES on a GE Discovery XR656 at $2021\times2021$ pixels of 0.194~mm.
The released data set carries a non-commercial licence (Chapter~\ref{ch:data}). In this study the SZCH-X-Rays pairs are used only as
evaluation references, never for training. Three properties of DES references matter for that use. They are themselves estimates, with residual
bone, noise and possible motion artefacts; their tone and texture differ from the radiograph they are paired with; and they come from one
institution. Chapter~\ref{ch:benchmark} therefore reports metrics both as computed and after a per-image gain and offset fit, and treats
agreement with DES as one line of evidence among several.

\section{Learning-based bone suppression}
\label{sec:bg_learned}
This section follows the field roughly in chronological order, grouping closely related studies, and describes each method only as far as its
abstract supports. Input resolution, which most abstracts do not state, is left out. Table~\ref{tab:bg_learned} summarises the learned methods; Section~\ref{sec:related} places them
relative to this study.

\begin{table}[htbp]
\centering\small
\renewcommand{\arraystretch}{1.15}
\caption{Learned bone-suppression methods discussed in Section~\ref{sec:bg_learned}, described only as far as their abstracts support. Rows
follow the order of first appearance (preprint date where there is one); the year is that of the bibliography entry; data sizes are those stated
in the abstracts.}
\label{tab:bg_learned}
\begin{tabular}{@{}p{0.07\linewidth}p{0.21\linewidth}p{0.29\linewidth}p{0.30\linewidth}@{}}
\toprule
Year & Method & Model & Training pairs named in the abstract \\
\midrule
2006 & MTANN \citep{suzuki2006mtann} & multiresolution massive-training ANN (non-linear filter) & DES bone images as teaching images (number not stated) \\
2017 & CamsNet \citep{yang2017cascade} & cascade of multi-scale CNNs in the gradient domain & 504 two-exposure DES cases (404 training, 100 test) \\
2017 & \citet{gusarev2017deep} & two end-to-end networks & not stated in the abstract \\
2024 & BS-Diff \citep{chen2024bsdiff} & conditional diffusion model (U-Net) with an autoencoder-based enhancement module & radiographs and soft-tissue images of 120 patients \\
2026 & BS-LDM \citep{sun2026bsldm} & conditional latent diffusion with a VQ-GAN & SZCH-X-Rays (818 pairs); 241 JSRT pairs \\
2025 & GL-LCM \citep{sun2025gllcm} & latent consistency model with global and local paths & SZCH-X-Rays; JSRT \\
2026 & \citet{angaitkar2026anatomy} & networks that predict one CT-projected component & projections of chest CT; no real pairs \\
2026 & DeBoneDiT \citep{sun2026debonedit} & source-conditioned Brownian-bridge diffusion transformer & SZCH-X-Rays release (741 DES pairs) \\
\bottomrule
\end{tabular}
\end{table}

\paragraph{Massive-training neural networks (2006).} \citet{suzuki2006mtann} trained a multiresolution massive-training artificial neural
network (MTANN), a highly non-linear filter applied separately at three resolution levels, to map a radiograph to its DES bone image. The
network outputs a bone-like image, and subtracting it from the radiograph gives a soft-tissue-like image. On a validation set of 118 radiographs
with nodules and an independent test set of 136 digitised screen-film radiographs from 14 institutions, ribs and clavicles were substantially
suppressed while the visibility of nodules and vessels was maintained; the abstract reports no detection measure. This
predict-the-bone-and-subtract design is the one used in this study.

\paragraph{Reader and CAD studies (2011--2014).} The reader study of \citet{freedman2011lung}, summarised in Section~\ref{sec:bg_cxr}, showed
higher sensitivity with suppression software at some cost in specificity. A later study asked whether computer-aided detection (CAD) adds anything once readers
already have bone-suppressed images: with 300 subjects (111 with a solitary nodule) read by eight observers with PA and lateral radiographs and
bone-suppressed images, reading with CAD marks increased the JAFROC figure of merit, although observers dismissed most of the true-positive CAD
candidates they had missed \citep{schalekamp2014cad}. These studies use commercial software whose algorithm their abstracts do not describe.

\paragraph{Contrast of inserted nodules (2013, 2020).} Two studies measured what suppression does to a nodule whose position is known.
\citet{hogeweg2013suppression} suppressed elongated structures (ribs, clavicles, catheters) with a non-learned method that models intensity
profiles across the structure and separates them by blind source separation. In 100 radiographs with nodules extracted from CT and simulated near
the clavicles, the abstract reports that clavicle contrast fell from 2.46 to 0.55 after suppression, while nodule contrast fell from 1.83 to 1.35:
suppression removed most of the clavicle and some of the nodule. \citet{takagi2020dose} imaged an anthropomorphic phantom with a 12~mm simulated
nodule at 66 positions under several tube voltages and doses; in areas not overlapping the ribs, bone suppression significantly increased nodule
contrast and CNR, and bone-suppressed images at 70~kVp were equal to or better than those at 140~kVp. The abstract does not name the
suppression software. Both studies measure contrast of a known nodule; neither audits a learned suppressor against an exact, additive lesion
signal (Section~\ref{sec:related}).

\paragraph{Convolutional networks (2017).} \citet{yang2017cascade} trained a cascade of multi-scale convolutional networks (CamsNet) to map
the gradients of a radiograph to the gradients of its bone image, refined the prediction from coarse to fine scales, fused the predicted bone
gradients in a maximum-a-posteriori framework to reconstruct the bone image, and subtracted it from the radiograph. Training and testing used 504
cases of real two-exposure DES radiographs (404 for training and 100 for testing) with a pixel size of about 0.194~mm, and the trained model was
also applied to radiographs from other machines, including scanned films. A plausible reason for working in the gradient domain is that the bone
signal varies most at the rib edges, so gradients concentrate the learning problem there. \citet{gusarev2017deep} proposed two end-to-end architectures that treat bones as noise to be
removed and reported that the scheme does not introduce artefacts; their abstract does not describe the training data.

\paragraph{Diffusion models on DES pairs (2024--2026).} Diffusion models generate an image by reversing a gradual noising process; a
\emph{conditional} diffusion model generates the soft-tissue image given the radiograph. BS-Diff \citep{chen2024bsdiff} combines a conditional
diffusion U-Net with an enhancement module that includes an autoencoder, trained on paired radiographs and soft-tissue images of 120 patients
from the authors' hospital, and evaluates it with comparisons, ablations and clinical evaluations. BS-LDM \citep{sun2026bsldm} moves the diffusion
into the latent space of a vector-quantised GAN trained with a multi-level hybrid loss, adds offset noise in the forward process and a temporal
adaptive thresholding in the reverse process to preserve low-frequency soft-tissue information, and introduces SZCH-X-Rays; it also uses 241
pairs from the JSRT database \citep{shiraishi2000jsrt} converted to negative images. GL-LCM \citep{sun2025gllcm} replaces diffusion sampling by a
latent consistency model for speed and combines a global path with a local path guided by lung segmentation, adding a training-free local
guidance to reduce boundary artefacts; it is evaluated on SZCH-X-Rays and JSRT. DeBoneDiT \citep{sun2026debonedit} formulates suppression as a
source-conditioned Brownian-bridge diffusion that starts from the radiograph instead of from noise, uses a transformer denoiser in the latent
space of an autoencoder trained with a feature loss from a pretrained depth-estimation encoder, and releases SZCH-X-Rays. These methods produce
the soft-tissue image directly and are trained on real DES pairs from one institution; generative models of this kind can also produce content
that is not in the input (Section~\ref{sec:related}).

\paragraph{CT-supervised suppression (2026).} \citet{angaitkar2026anatomy} replace real pairs by digitally reconstructed radiographs: chest CT
volumes are decomposed into bone, non-lung soft tissue and lung, each component is projected separately, and the weighted components sum exactly
to the full projection. Models learn to predict one component (bone or lung) and recover the rest by subtraction, and they transfer to real
radiographs without real paired training data. The abstract reports detection gains concentrated on abnormalities with substantial bone overlap and
releases models and projections. Supervision of this kind is exact and scalable, because every CT volume yields a pair, but it carries a domain gap
between projections and real radiographs, and the lesions of the training CTs are whatever the CT collection happens to contain. This study builds
on these projections (Chapter~\ref{ch:v2}).

\paragraph{Summary.} Three sources of supervision appear in the literature: real DES pairs, which are realistic but scarce, noisy and
institution-specific; non-learned models of the bone, which need no pairs but have hand-set assumptions; and CT projections, which are exact and
plentiful but synthetic. Evaluation has used similarity to DES images, reader and CAD studies, downstream detection and the contrast of known
nodules. What a learned suppressor does to a soft-tissue lesion that lies on a rib, measured against an exactly known lesion signal, is the
question of this thesis.

\section{Architectures and optimisation}
\label{sec:bg_arch}
The suppressor of this study is an encoder--decoder network that maps a radiograph $x$ to a bone image $\hat b$; the soft image is
$\hat s=x-\hat b$. Predicting the bone image rather than the soft image keeps the output sparse (zero outside bone) and makes the identity
$x=\hat s+\hat b$ hold by construction. The encoder is a pretrained image backbone that produces a feature pyramid; the decoder merges the pyramid
back to full resolution. The confirmatory dose-response and the ablation use a Swin Transformer V2 Tiny encoder (SwinV2-T); the final model uses a
ConvNeXt V2 Tiny encoder (ConvNeXt-V2-T), chosen by the pre-declared rule described in Section~\ref{sec:experiments}. Both produce four feature maps
at strides 4, 8, 16 and 32, so the same decoder, after a $1\times1$ projection of each map, serves both.

\paragraph{Swin Transformer V2.} A vision transformer cuts an image into patches, embeds each as a token and lets every token attend to every
other; the cost is quadratic in the number of tokens, which is prohibitive for dense prediction at high resolution. The Swin Transformer computes
self-attention only inside non-overlapping windows of $M\times M$ tokens and shifts the window grid by half a window between consecutive blocks, so
that information crosses window borders; patch-merging layers halve the resolution and double the channels between stages, which gives a
hierarchical pyramid like that of a convolutional network \citep{dosovitskiy2021vit,liu2021swin}. Swin Transformer V2 \citep{liu2022swinv2} adds three changes to scale capacity and
resolution: a residual post-norm combined with cosine attention to stabilise training, a log-spaced continuous position bias to transfer models
pre-trained at low resolution to high-resolution inputs, and SimMIM self-supervised pre-training. In cosine attention the similarity of a query
$q_i$ and a key $k_j$ is
\begin{equation}
\mathrm{Sim}(q_i,k_j)=\frac{\cos(q_i,k_j)}{\tau}+B_{ij},
\label{eq:bg_cosattn}
\end{equation}
with a learnable temperature $\tau$ and a relative position bias $B_{ij}$ produced by a small network from log-spaced relative coordinates.
Because attention is computed within windows of fixed size, an encoder pretrained on small images can be applied to a larger input with the same
window size; only the number of windows grows. The continuous bias additionally allows the window size itself to change.

\paragraph{ConvNeXt V2.} ConvNeXt is a convolutional network modernised with design choices from transformers: a patchify stem, large ($7\times7$)
depthwise convolutions, layer normalisation and inverted-bottleneck MLPs with GELU \citep{liu2022convnext}. ConvNeXt V2 \citep{woo2023convnextv2} co-designs the
architecture with masked-autoencoder pre-training: it proposes a fully convolutional masked autoencoder (FCMAE) and a Global Response
Normalization (GRN) layer that enhances inter-channel feature competition. For a feature map $X\in\mathrm{R}^{H\times W\times C}$, GRN computes
the spatial $\ell_2$ norm of each channel, normalises it by the mean over channels and rescales the features:
\begin{equation}
\begin{gathered}
G_i=\lVert X_{:,:,i}\rVert_2,\\
N_i=\frac{G_i}{\frac1C\sum_{j=1}^{C}G_j},\\
X_{:,:,i}\leftarrow\gamma_i\,X_{:,:,i}\,N_i+\beta_i+X_{:,:,i},
\end{gathered}
\label{eq:bg_grn}
\end{equation}
with learnable $\gamma$ and $\beta$. The encoder used here was pre-trained with FCMAE and then fine-tuned on ImageNet \citep{deng2009imagenet}
(Section~\ref{sec:method_v2}).

\paragraph{Nested decoder (UNet++).} A U-Net decoder upsamples the deepest features step by step and concatenates, at each scale, the encoder
features of the same resolution (skip connections). UNet++ \citep{zhou2018unetpp} replaces the plain skips by nested, dense skip pathways: node
$X^{i,j}$ at depth $i$ and position $j$ along the skip receives all previous nodes at the same depth and the upsampled node from the next depth,
\begin{equation}
X^{i,j}=\mathcal{H}\Big(\big[X^{i,0},X^{i,1},\dots,X^{i,j-1},\,\mathcal{U}(X^{i+1,j-1})\big]\Big),
\label{eq:bg_unetpp}
\end{equation}
where $\mathcal{H}$ is a convolution block, $\mathcal{U}$ upsampling and $[\cdot]$ concatenation. The intermediate nodes narrow the semantic gap
between encoder and decoder features before they are fused. UNet++ also proposes deep supervision of the intermediate outputs; the decoder here
uses only the final node. Because the transformer and ConvNeXt pyramids start at stride 4, the network of this study adds a short tail with two
shallow convolutional skips from the input to recover full resolution (Section~\ref{sec:methods}).

\paragraph{CBAM.} The Convolutional Block Attention Module \citep{woo2018cbam} reweights a feature map $F$ first by channel and then by
position:
\begin{equation}
\begin{gathered}
M_c(F)=\sigma\big(\mathrm{MLP}(\mathrm{AvgPool}(F))+\mathrm{MLP}(\mathrm{MaxPool}(F))\big),\\
M_s(F)=\sigma\big(f^{7\times7}([\mathrm{AvgPool}_c(F);\mathrm{MaxPool}_c(F)])\big),
\end{gathered}
\label{eq:bg_cbam}
\end{equation}
with $F'=M_c(F)\odot F$ and $F''=M_s(F')\odot F'$. Here $\sigma$ is the logistic function, the MLP is shared between the two pooled descriptors,
$\mathrm{AvgPool}_c$ and $\mathrm{MaxPool}_c$ pool over channels and $f^{7\times7}$ is a $7\times7$ convolution. In a suppressor the spatial map
can emphasise rib regions; the module is cheap and is placed after the decoder blocks.

\paragraph{AdamW, schedule and clipping.} Adam keeps exponential moving averages of the gradient $g_t$ and of its square, $m_t$ and $v_t$, and
scales each parameter's step by $1/\sqrt{\hat v_t}$. With an $L_2$ penalty added to the loss, the penalty gradient is rescaled in the same way, so
parameters with large gradient variance are regularised less. AdamW \citep{loshchilov2019adamw} decouples weight decay from the adaptive update:
\begin{equation}
\theta_t=\theta_{t-1}-\eta_t\Big(\frac{\hat m_t}{\sqrt{\hat v_t}+\epsilon}+\lambda\,\theta_{t-1}\Big),
\label{eq:bg_adamw}
\end{equation}
where $\hat m_t$ and $\hat v_t$ are the bias-corrected moments and $\lambda$ the decoupled weight-decay rate. The learning rate $\eta_t$ rises
linearly during a short warm-up of $t_w$ steps and then follows a cosine decay to zero at step $T$,
\begin{equation}
\begin{gathered}
\eta_t=\eta_{\max}\,\tfrac12\Big(1+\cos\Big(\pi\,\frac{t-t_w}{T-t_w}\Big)\Big),\\
t_w\le t\le T .
\end{gathered}
\label{eq:bg_cosine}
\end{equation}
A pretrained encoder is often given a smaller learning rate than a randomly initialised decoder, so that the pretrained features are adapted
rather than overwritten. Gradient clipping rescales the full gradient $g$ to $g\,\min(1,c/\lVert g\rVert_2)$ when its norm exceeds $c$, which
limits the damage of rare very large gradients.

\paragraph{EMA weights.} An exponential moving average of the weights,
\begin{equation}
\bar\theta_t=d_t\,\bar\theta_{t-1}+(1-d_t)\,\theta_t ,
\label{eq:bg_ema}
\end{equation}
is kept alongside the trained weights and used for validation and evaluation. Averaging over the last part of the trajectory smooths the noise of
stochastic optimisation and usually gives a slightly better and more stable model than the last iterate. The decay $d_t$ is small at the start and
grows to a fixed maximum, so that the average is not dominated by the random initialisation; normalisation statistics (buffers) are copied rather
than averaged.

\paragraph{Mixed precision.} Mixed-precision training runs most of the forward and backward pass in 16-bit floating point (fp16), which halves
memory traffic and uses the fast matrix units of current GPUs, while the master weights and the optimiser state stay in 32-bit precision. The
fp16 format has a small range: very small gradients underflow to zero and values above about $6.5\times10^4$ overflow. Dynamic loss scaling
multiplies the loss by a factor $S$ before the backward pass, divides the gradients by $S$ before clipping and the optimiser step, skips any step
whose gradients contain infinities or NaNs, and adapts $S$ during training. Loss scaling does not protect activations in the forward pass, and
fp16 overflow can still make a run diverge; Section~\ref{sec:experiments} reports a divergence attributed to fp16 overflow and how it was handled.

\section{Metrics}
\label{sec:bg_metrics}
Throughout this section $y$ is a reference image (for example a DES soft-tissue image), $\hat y$ an estimate of it, both scaled to $[0,1]$, $p$
indexes pixels and $R$ is a set of pixels (the whole image or a region).

\paragraph{Tone calibration.} A DES image and a radiograph differ in global tone (Section~\ref{sec:bg_des}), and a suppressor trained on other
data cannot know the tone of the reference. Fidelity metrics are therefore also computed after a per-image least-squares gain and offset,
\begin{equation}
\begin{gathered}
\hat y_{\mathrm{cal}}=\mathrm{clip}_{[0,1]}\big(a^\ast\hat y+c^\ast\big),\\
(a^\ast,c^\ast)=\arg\min_{a,c}\sum_p\big(a\,\hat y_p+c-y_p\big)^2 ,
\end{gathered}
\label{eq:bg_calib}
\end{equation}
which removes differences in brightness and contrast but not in structure. Section~\ref{sec:methods} calls these metrics ``calibrated''.

\paragraph{PSNR.} The mean squared error over $R$ and the peak signal-to-noise ratio are
\begin{equation}
\begin{gathered}
\mathrm{MSE}_R=\frac{1}{|R|}\sum_{p\in R}\big(\hat y_p-y_p\big)^2,\\
\mathrm{PSNR}_R=10\log_{10}\frac{L^2}{\mathrm{MSE}_R}\ \mathrm{dB},
\end{gathered}
\label{eq:bg_psnr}
\end{equation}
with dynamic range $L=1$. PSNR is a monotone transform of the MSE: it weights every pixel equally and is dominated by large errors, so it is
sensitive to the global tone and to a few badly wrong pixels. Restricting $R$ to the region where bone was removed (the bone region of
Section~\ref{sec:methods}) focuses the metric on the pixels that suppression should change.

\paragraph{SSIM.} The structural similarity index \citep{wang2004ssim} compares local means, variances and covariance of two images $u$ and $v$ in
a sliding window:
\begin{equation}
\begin{gathered}
\mathrm{SSIM}(u,v)=\frac{(2\mu_u\mu_v+C_1)\,(2\sigma_{uv}+C_2)}{(\mu_u^2+\mu_v^2+C_1)\,(\sigma_u^2+\sigma_v^2+C_2)},\\
C_1=(K_1L)^2,\ C_2=(K_2L)^2,
\end{gathered}
\label{eq:bg_ssim}
\end{equation}
with the customary $K_1=0.01$ and $K_2=0.03$. It is the product of a luminance term
$l=(2\mu_u\mu_v+C_1)/(\mu_u^2+\mu_v^2+C_1)$, a contrast term $c=(2\sigma_u\sigma_v+C_2)/(\sigma_u^2+\sigma_v^2+C_2)$ and a structure term
$s=(\sigma_{uv}+C_2/2)/(\sigma_u\sigma_v+C_2/2)$. The local values form an SSIM map, whose mean over $R$ is reported. SSIM ranges up to 1
(identical images); it is more sensitive to structure and local contrast than PSNR, and it depends on the window, which differs between
implementations (the original uses an $11\times11$ Gaussian window; the scikit-image function used here uses a $7\times7$ uniform window by
default).

\paragraph{MS-SSIM.} Multi-scale SSIM \citep{wang2003msssim} evaluates the contrast and structure terms at $M$ scales, obtained by repeated
low-pass filtering and down-sampling by two, and the luminance term only at the coarsest scale:
\begin{equation}
\mathrm{MS\text{-}SSIM}(u,v)=\big[l_M(u,v)\big]^{\alpha_M}\prod_{j=1}^{M}\big[c_j(u,v)\big]^{\beta_j}\big[s_j(u,v)\big]^{\gamma_j},
\label{eq:bg_msssim}
\end{equation}
with $M=5$ and the exponents fitted to human judgements in the original paper ($\beta_j=\gamma_j=0.0448$, $0.2856$, $0.3001$, $0.2363$, $0.1333$
and $\alpha_M=\beta_M$). It is less dominated by fine texture than single-scale SSIM.

\paragraph{Band-limited SSIM.} A DES soft-tissue image is filtered and noisier than the radiograph, so a suppressor that keeps the fine texture
of the radiograph can be penalised by full-resolution SSIM for texture that has nothing to do with bone. Band-limited SSIM compares the two images
after the same Gaussian low-pass filter $G_\sigma$,
\begin{equation}
\mathrm{SSIM}_\sigma(\hat y,y)=\mathrm{SSIM}\big(G_\sigma\ast\hat y,\;G_\sigma\ast y\big),
\label{eq:bg_blssim}
\end{equation}
which keeps the scales of ribs and soft-tissue structures and discards the finest texture. The values of $\sigma$ are fixed in the analysis plan
and applied identically to every method.

\paragraph{LPIPS.} The learned perceptual image patch similarity \citep{zhang2018lpips} measures distance in the feature space of a pretrained
network. With $\hat u^{l}_{hw}$ and $\hat v^{l}_{hw}$ the channel-normalised activations of layer $l$ at position $(h,w)$,
\begin{equation}
d(u,v)=\sum_l\frac{1}{H_lW_l}\sum_{h,w}\big\lVert w_l\odot\big(\hat u^{l}_{hw}-\hat v^{l}_{hw}\big)\big\rVert_2^2 ,
\label{eq:bg_lpips}
\end{equation}
where the channel weights $w_l$ were fitted to human similarity judgements. Lower is better. This study uses the AlexNet variant and replicates
the grey image to three channels. The network and the human judgements come from natural images, so LPIPS is a secondary, perceptual check
rather than a measure of diagnostic content.

\paragraph{Contrast-to-noise ratio of a box.} For a lesion with bounding box $\mathcal{B}$, let $\mathcal{R}$ be the ring obtained by enlarging
the box by half its width on the left and right and by half its height above and below, minus the box itself (clipped to the image). The
contrast-to-noise ratio is
\begin{equation}
\mathrm{CNR}=\frac{\bar x_{\mathcal{B}}-\bar x_{\mathcal{R}}}{\mathrm{sd}_{\mathcal{R}}(x)+\varepsilon},
\label{eq:bg_cnr}
\end{equation}
with means $\bar x$ and standard deviation $\mathrm{sd}$ over the indicated pixels and a small $\varepsilon$. The effect of suppression is the
change $\Delta\mathrm{CNR}=\mathrm{CNR}(\hat s)-\mathrm{CNR}(x)$ on the same box. Two effects are mixed in this number: removing rib edges from
the ring lowers $\mathrm{sd}_{\mathcal{R}}$ and raises the CNR, while moving part of the lesion into the bone image lowers the numerator. A
positive change therefore does not prove that the lesion was preserved, and the CNR is used as evidence that complements, but does not replace,
the insertion audit. A box also contains background around the lesion, so its mean is a diluted measure of lesion contrast.

\paragraph{Detection: FROC and CPM.} A detector outputs boxes with scores. A predicted box is matched to a ground-truth box if their
intersection over union, $\mathrm{IoU}=|\mathcal{B}_1\cap\mathcal{B}_2|/|\mathcal{B}_1\cup\mathcal{B}_2|$, reaches a threshold; predictions are
matched greedily in order of decreasing score and each ground-truth box at most once. For a score threshold $t$, let $\mathrm{TP}(t)$ and
$\mathrm{FP}(t)$ be the numbers of matched and unmatched predictions with score at least $t$, $N_{\mathrm{gt}}$ the number of ground-truth
lesions and $N_{\mathrm{img}}$ the number of images. The free-response ROC (FROC) curve plots the sensitivity
$\mathrm{TP}(t)/N_{\mathrm{gt}}$ against the mean number of false positives per image $\mathrm{FP}(t)/N_{\mathrm{img}}$ as $t$ varies. The
competition performance metric (CPM) averages the sensitivity at seven operating points,
\begin{equation}
\begin{gathered}
\mathrm{CPM}=\frac17\sum_{r\in\{\frac18,\frac14,\frac12,1,2,4,8\}}\mathrm{Sens}(r),\\
\mathrm{Sens}(r)=\max_{t:\ \mathrm{FP}(t)/N_{\mathrm{img}}\le r}\frac{\mathrm{TP}(t)}{N_{\mathrm{gt}}},
\end{gathered}
\label{eq:bg_cpm}
\end{equation}
as introduced in the ANODE09 study and named CPM by \citet{niemeijer2011cpm} \citep[see also][]{vanginneken2010anode09,setio2017luna16}. Tied
scores are treated as one threshold, so the curve does not depend on how a sort orders ties. FROC analysis suits lesion detection because an
image can contain several lesions and several false marks, which an image-level ROC curve cannot represent.

\paragraph{Average precision.} With precision $P(t)=\mathrm{TP}(t)/(\mathrm{TP}(t)+\mathrm{FP}(t))$ and recall
$\mathrm{Rec}(t)=\mathrm{TP}(t)/N_{\mathrm{gt}}$, the interpolated precision at recall $r$ is $P_{\mathrm{int}}(r)=\max_{r'\ge r}P(r')$ and the
average precision is the area under the interpolated precision--recall curve,
\begin{equation}
\mathrm{AP}=\sum_k\big(r_k-r_{k-1}\big)\,P_{\mathrm{int}}(r_k),
\label{eq:bg_ap}
\end{equation}
summed over the distinct recall levels $r_k$ (all-point interpolation). The mean over classes at an IoU threshold of 0.5 is reported as
mAP@50. Unlike the CPM, AP penalises false positives through precision rather than per image.

\paragraph{AUROC.} For a binary label with $m$ positive cases with scores $X_1,\dots,X_m$ and $n$ negative cases with scores $Y_1,\dots,Y_n$, the
area under the ROC curve equals the Mann--Whitney statistic
\begin{equation}
\begin{gathered}
\hat\theta=\frac{1}{mn}\sum_{i=1}^{m}\sum_{j=1}^{n}\psi(X_i,Y_j),\\
\psi(a,b)=\begin{cases}1 & a>b\\ \tfrac12 & a=b\\ 0 & a<b\end{cases},
\end{gathered}
\label{eq:bg_auc}
\end{equation}
the probability that a random positive case scores higher than a random negative one. It does not depend on a threshold or on the class
prevalence. For several labels (for example the findings of a multi-label classifier) the macro-AUROC is the unweighted mean of the per-label AUROCs.

\section{Uncertainty estimation}
\label{sec:bg_uncertainty}
A suppressor returns one bone image; a user also needs to know where that image may be wrong. Two kinds of uncertainty are usually distinguished:
\emph{aleatoric} uncertainty, which comes from ambiguity in the data and cannot be reduced by more training data (the rib-or-nodule ambiguity of
Section~\ref{sec:bg_cxr} is of this kind), and \emph{epistemic} uncertainty, which comes from limited data and model knowledge and grows on
inputs unlike the training data. This study uses a deep ensemble to obtain a per-pixel uncertainty score and split conformal prediction to turn
that score into intervals with a stated coverage.

\paragraph{Deep ensembles.} \citet{lakshminarayanan2017ensembles} train several networks independently from different random initialisations
(and data orders) and combine their predictions; they also use proper scoring rules and, optionally, adversarial training. For $K$ members
$f_1,\dots,f_K$, the ensemble prediction and spread of the bone image at pixel $p$ are
\begin{equation}
\begin{gathered}
\bar f_p(x)=\frac1K\sum_{k=1}^{K}f_{k,p}(x),\\
\sigma_p(x)=\Big(\frac{1}{K-1}\sum_{k=1}^{K}\big(f_{k,p}(x)-\bar f_p(x)\big)^2\Big)^{1/2}.
\end{gathered}
\label{eq:bg_ensemble}
\end{equation}
The mean is usually more accurate than a single member, and the spread is large where the members disagree, which happens mostly where the data
did not pin the function down. In this study each member outputs only a bone image, not a variance, so $\sigma$ reflects disagreement between
members rather than noise in the target. With few members $\sigma$ is itself a noisy estimate, and nothing guarantees that it is calibrated: a
spread of $0.01$ need not mean an error of about $0.01$. Whether the spread ranks pixels by their error is checked with the Spearman correlation
between spread and absolute error (Section~\ref{sec:bg_stats}) and with sparsification curves (Section~\ref{sec:method_v2}).

\paragraph{Split conformal prediction.} Conformal prediction turns any score into a prediction set with a guaranteed coverage
\citep{lei2018conformal,angelopoulos2023conformal}. In the split version the model is fixed and a calibration set
$(X_1,Y_1),\dots,(X_n,Y_n)$, disjoint from the training data, is used to compute nonconformity scores. For a regression output with an
uncertainty scale, a natural score is the normalised residual
\begin{equation}
S_i=\frac{|Y_i-\bar f(X_i)|}{\sigma(X_i)+\varepsilon},
\label{eq:bg_score}
\end{equation}
with a small $\varepsilon$ that keeps the interval open where the members agree exactly. Let $\hat q$ be the
$\lceil (n+1)(1-\alpha)\rceil$-th smallest of $S_1,\dots,S_n$. The interval for a new input is
\begin{equation}
C(X_{n+1})=\big[\bar f(X_{n+1})-\hat q\,(\sigma(X_{n+1})+\varepsilon),\;\bar f(X_{n+1})+\hat q\,(\sigma(X_{n+1})+\varepsilon)\big].
\label{eq:bg_conformal}
\end{equation}
For the bone image $Y$ is a bone pixel value; because $\hat s=x-\hat b$ and $x$ is known, the soft image has an interval of the same width.

\paragraph{Exchangeability and marginal coverage.} The guarantee rests on one assumption: the calibration pairs and the test pair
$(X_1,Y_1),\dots,(X_{n+1},Y_{n+1})$ are \emph{exchangeable}, that is, their joint distribution does not change when the indices are permuted.
Independent, identically distributed data are exchangeable. Under exchangeability the rank of $S_{n+1}$ among $S_1,\dots,S_{n+1}$ is uniformly
distributed, from which it follows that
\begin{equation}
\Pr\big(Y_{n+1}\in C(X_{n+1})\big)\ \ge\ 1-\alpha ,
\label{eq:bg_coverage}
\end{equation}
and, if the scores are almost surely distinct, the coverage is also at most $1-\alpha+1/(n+1)$ \citep{lei2018conformal}. Three limits of this
statement matter for this study. First, the coverage is \emph{marginal}: the probability is over the random calibration set and the random test
point together, so the interval covers $1-\alpha$ of test points on average, not $1-\alpha$ of the pixels of every image or of every region (for
example the pixels over a lesion). Second, when pixels are pooled as calibration units, the pixels of one image are dependent, so exchangeability
of pixels is an approximation; the resulting statement is about pixels drawn as in calibration, not about images (Section~\ref{sec:method_v2}).
Third, the guarantee holds only for test data from the calibration distribution. Calibrating on CT projections and testing on real radiographs
breaks exchangeability, so coverage on real images is measured, not guaranteed, and is reported as coverage under shift
(Chapter~\ref{ch:reliability}).

\section{Statistical methods}
\label{sec:bg_stats}
All comparisons in this thesis are \emph{paired}: every method is scored on the same units (images, lesion sites, patients), and inference is
based on the per-unit differences $d_i=a_i-b_i$ between two methods. Pairing removes the large variation between units that both methods share and
makes small differences detectable.

\paragraph{Percentile bootstrap.} The bootstrap \citep{efron1993bootstrap} estimates the sampling distribution of a statistic by resampling the
observed data. For the mean difference $\hat\theta=\frac1n\sum_i d_i$, draw $B$ samples of size $n$ with replacement from $d_1,\dots,d_n$,
compute the mean $\hat\theta^{\ast}_b$ of each, and take the empirical quantiles
\begin{equation}
\big[\hat\theta^{\ast}_{(\alpha/2)},\ \hat\theta^{\ast}_{(1-\alpha/2)}\big]
\label{eq:bg_boot}
\end{equation}
as a $1-\alpha$ confidence interval. The percentile interval needs no normality assumption; it assumes that the units are independent and that the
sample is large enough for the bootstrap distribution to approximate the sampling distribution.

\paragraph{Cluster bootstrap.} Units are often nested: several lesion sites per image, several images per patient. Sites in one image share the
image's anatomy and the model's behaviour on it, so they are positively correlated, and resampling sites as if independent gives intervals that
are too narrow. The cluster bootstrap resamples whole clusters $g=1,\dots,G$ with replacement and recomputes the statistic over all units in the
resampled clusters,
\begin{equation}
\hat\theta^{\ast}_b=\frac{\sum_{g\in\mathcal{S}_b}\sum_{i\in g}d_i}{\sum_{g\in\mathcal{S}_b}n_g},
\label{eq:bg_clusterboot}
\end{equation}
where $\mathcal{S}_b$ is the $b$-th resample of cluster indices (with repetitions) and $n_g$ the number of units in cluster $g$. The between-cluster
variation then enters the interval as it should. In this study, for example, lesion sites are clustered by image and images by patient
(Section~\ref{sec:experiments}).

\paragraph{Wilcoxon signed-rank test.} \citet{wilcoxon1945individual} proposed ranking the absolute paired differences. Zero differences are
dropped; the remaining $|d_i|$ are ranked from 1 to $n$ (ties receive average ranks), and the statistic is the sum of the ranks of the positive
differences,
\begin{equation}
\begin{gathered}
W^{+}=\sum_{i:\,d_i>0}\mathrm{rank}(|d_i|),\\
\mathrm{E}[W^{+}]=\frac{n(n+1)}{4},\\
\mathrm{Var}[W^{+}]=\frac{n(n+1)(2n+1)}{24}
\end{gathered}
\label{eq:bg_wilcoxon}
\end{equation}
under the null hypothesis that the differences are distributed symmetrically about zero (variance without ties). Large samples use the normal
approximation, small ones the exact distribution. The test is robust to outliers and does not assume normality, but it tests a shift in location
(the pseudo-median of the differences), not the mean; the bootstrap interval of the mean and the Wilcoxon $p$-value therefore answer slightly
different questions and are reported together.

\paragraph{Holm correction.} Testing $m$ hypotheses at level $\alpha$ each inflates the probability of at least one false rejection (the
family-wise error rate). The Holm procedure \citep{holm1979simple} orders the $p$-values $p_{(1)}\le\dots\le p_{(m)}$ and rejects
$H_{(1)},\dots,H_{(k-1)}$, where $k$ is the first index with $p_{(k)}>\alpha/(m-k+1)$. Equivalently, the adjusted $p$-values
\begin{equation}
\tilde p_{(i)}=\max_{j\le i}\ \min\big(1,\,(m-j+1)\,p_{(j)}\big)
\label{eq:bg_holm}
\end{equation}
are compared with $\alpha$. Holm controls the family-wise error rate under any dependence between the tests and rejects at least as much as
the Bonferroni correction. It is applied within each pre-declared family of hypotheses; confidence intervals are not adjusted.

\paragraph{One-sided non-inferiority.} A new method that adds a regulariser is not expected to be better than a reference on every metric; the
question is whether it is \emph{not worse by more than a tolerable amount}. For a metric where higher is better and a margin $\delta>0$ fixed in
advance, the hypotheses are
\begin{equation}
H_0:\ \Delta\le-\delta\qquad\text{against}\qquad H_1:\ \Delta>-\delta ,
\label{eq:bg_noninf}
\end{equation}
where $\Delta$ is the mean paired difference (new minus reference). $H_0$ is rejected at level $0.05$ if the lower limit of a one-sided 95\,\%
confidence interval for $\Delta$, which equals the lower limit of a two-sided 90\,\% interval, lies above $-\delta$. Three points are often
confused. Failing to show non-inferiority does not show inferiority; non-inferiority does not mean equality; and a margin chosen after seeing the
data makes the test meaningless, which is why the margins of this study were fixed in the analysis plan (Section~\ref{sec:experiments}).

\paragraph{DeLong test for paired AUCs.} When two classifiers score the same cases, their AUCs are correlated, and comparing them as if
independent wastes power. \citet{delong1988auc} use the structure of Equation~\eqref{eq:bg_auc}. For classifier $r$ define the placement values
\begin{equation}
\begin{gathered}
V^{r}_{10}(X_i)=\frac1n\sum_{j=1}^{n}\psi(X^{r}_i,Y^{r}_j),\\
V^{r}_{01}(Y_j)=\frac1m\sum_{i=1}^{m}\psi(X^{r}_i,Y^{r}_j),
\end{gathered}
\label{eq:bg_placements}
\end{equation}
whose means are both $\hat\theta^{r}$. With $S_{10}$ and $S_{01}$ the sample covariance matrices of the placement values over the positive and
the negative cases, the covariance of the AUC estimates is $S=S_{10}/m+S_{01}/n$, and
\begin{equation}
z=\frac{\hat\theta^{1}-\hat\theta^{2}}{\sqrt{S_{11}+S_{22}-2S_{12}}}
\label{eq:bg_delong}
\end{equation}
is compared with a standard normal distribution. The covariance term $S_{12}$ is what makes the test paired: the more the two classifiers agree
on which cases are hard, the smaller the variance of their difference.

\paragraph{Spearman correlation.} The Spearman coefficient is the Pearson correlation of the ranks of two variables; without ties it equals
\begin{equation}
\rho_S=1-\frac{6\sum_i D_i^2}{n(n^2-1)},
\label{eq:bg_spearman}
\end{equation}
where $D_i$ is the difference between the ranks of the two variables for unit $i$. It measures monotone rather than linear association and is
invariant to monotone transformations of either variable, which suits the comparison of an uncertainty score with an error, whose scales differ.

\paragraph{Training seeds.} The result of training a network depends on its random seed (initialisation, data order, augmentation). A difference
between two arms measured with one seed each confounds the effect of the arm with seed noise. Arms are therefore trained with several seeds,
seed results are averaged within each unit before a test, and the spread over seeds is reported alongside (Section~\ref{sec:methods}).

\paragraph{Pre-specified plans, development and confirmatory data.} A study that tries many models, metrics, subsets and tests and reports the
best has a much higher false-positive rate than its nominal $\alpha$, and a model chosen because it scored best on a data set will, on average,
score lower on new data from the same source (selection bias, or the winner's curse) \citep{simmons2011falsepositive,varoquaux2022failures}.
Two safeguards follow. First, development and confirmatory data are separated: every choice of architecture, weight, threshold and metric is
made on development data (here the validation projections and the SZCH development pairs), and the confirmatory data are scored only after
the choices are fixed. Second, the confirmatory analysis is pre-specified: hypotheses, primary metrics, tests, margins, hypothesis families and
the data lists are written down and frozen, with file hashes, before the confirmatory data are scored (Section~\ref{sec:experiments}).
Unlike a pre-registration, the plans and their hashes were kept in the project's version control and not deposited with a public registry,
so their timing rests on the commit history. A pre-specified plan does not forbid exploratory analyses; it makes the difference between
confirmatory and exploratory results visible. Any change
to the plan after freezing is a deviation and is reported as such (Appendix~\ref{app:deviations}), and a hypothesis that fails is reported as
failed.

\section{Related work}
\label{sec:related}
% PLAN (four subsections, about 1.5 pages). Describe a cited paper only as far as its abstract or title supports; check each claim against the PDF
% before submission.
%   2.1 Bone suppression: classical and CNN methods (suzuki2006mtann, yang2017cascade, gusarev2017deep); diffusion methods on real DES pairs
%       (chen2024bsdiff, sun2026bsldm, sun2025gllcm, sun2026debonedit). SZCH-X-Rays was introduced with BS-LDM and released with DeBoneDiT.
%   2.2 CT-supervised suppression (angaitkar2026anatomy; CT-RATE hamamci2026ctrate as the CT source of the projections we train on). What we reuse
%       (projections, released weights as comparator) and what we change (backbone, resolution, regulariser, evaluation).
%   2.3 Evaluation of restoration models for lesion fidelity: lesion insertion and removal, nodule generation (the NODE21 generation track,
%       sogancioglu2024node21), hallucination in generative restoration. [citations to be found, see paper/REFS_TODO.md]
%   2.4 Uncertainty and encoders: deep ensembles (lakshminarayanan2017ensembles), split conformal prediction (lei2018conformal,
%       angelopoulos2023conformal); modern encoders screened here (liu2022swinv2, woo2023convnextv2, simeoni2025dinov3, perezgarcia2025raddino).
\subsection{Bone suppression}
Early learning-based rib suppression used massive-training artificial neural networks \citep{suzuki2006mtann}; later work used convolutional
networks, for example a cascade of multi-scale networks in the gradient domain \citep{yang2017cascade} and other deep models
\citep{gusarev2017deep}. Recent methods train diffusion models on paired radiographs and DES soft-tissue images
\citep{chen2024bsdiff,sun2026bsldm,sun2025gllcm,sun2026debonedit}; the SZCH-X-Rays data set was compiled for BS-LDM \citep{sun2026bsldm} and made
public with DeBoneDiT \citep{sun2026debonedit}. BS-Diff trained on 120 pairs, augmented to 840, at $256\times256$ pixels; BS-LDM on 818 SZCH
pairs and 241 processed JSRT pairs at $1024\times1024$; GL-LCM and DeBoneDiT on the 741-pair SZCH version and JSRT at $1024\times1024$
(DeBoneDiT in a latent space of $4\times256\times256$).

\subsection{CT-supervised suppression}
\citet{angaitkar2026anatomy} build digitally reconstructed radiographs from chest CT, project bone, non-lung soft tissue and lung
separately, and train models that predict one component. Their projections are released on the Hugging Face Hub and are built from CT-RATE volumes
\citep{hamamci2026ctrate}: the release holds \num{21887} projections whose volume names are those of CT-RATE. We train on these projections and use the released weights as the main comparator. Our models differ in encoder and resolution (ConvNeXt-V2-T
at 512~px against SwinV2-B at 1024~px) and in sample budget, and add the twin regulariser, the ensemble and the audit.

\subsection{Evaluating lesion fidelity}
Learned image-to-image models can remove or invent content: distribution matching in translation networks can drop or hallucinate lesions
\citep{cohen2018distribution}, deep reconstruction can be unstable to small structures \citep{antun2020instabilities}, and hallucinations of
learned reconstruction have been defined and measured \citep{bhadra2021hallucinations,tivnan2024hallucination}. Inserting known features or signals into
images is an established way to test learned MRI reconstruction \citep{cheng2020addressing} and the effect of learned denoising on signal
detection \citep{li2021assessing}. For bone
suppression, simulated nodules were inserted to measure contrast after a hand-built clavicle suppression method \citep{hogeweg2013suppression},
and a phantom with a simulated nodule at many positions was used to measure contrast after bone suppression \citep{takagi2020dose}. These
studies measure lesion contrast; our audit measures the fraction of a known additive lesion signal that a learned suppressor moves into its
bone image, on real radiographs, and is shared by every model compared.

\subsection{Uncertainty and encoders}
Deep ensembles give a simple predictive-uncertainty estimate \citep{lakshminarayanan2017ensembles}, and split-conformal prediction turns an error score into
intervals with a coverage guarantee under exchangeability \citep{lei2018conformal,angelopoulos2023conformal}. We use both for the bone image
(Section~\ref{sec:experiments}). For the encoder we start from SwinV2-T \citep{liu2022swinv2} and screen ConvNeXt V2 \citep{woo2023convnextv2},
DINOv3 \citep{simeoni2025dinov3} and RAD-DINO \citep{perezgarcia2025raddino} under the same recipe: a modern convolutional encoder, a general
self-supervised vision transformer and a vision transformer pre-trained on chest radiographs.



\chapter{Data and physics}\label{ch:data}
% PLAN (20-25 pages). The paper text on data and methods is input in full; this chapter adds what the paper has no room for.
%   3.1 Data: paper/sections/data.tex (input below, brings its own \section). Add per-set preparation: file formats, preprocessing, orientation
%       checks, exclusions.
%   3.2 Methods: paper/sections/methods.tex (input below, brings its own \section; compact version of CT decomposition, projection, detector
%       physics, suppressor, v1 supervision, audit definition, real-image protocol, statistics). When the thesis text outgrows it, split methods.tex
%       into one file per subsection (in the paper folder) and input them here separately; do not fork the text.
%   3.3 Thesis-only expansion (below): segmentation recipe in detail, rendering geometry and frame, attenuation tables (ICRU-44
%       compositions, NIST tables with XCOM berger2010xcom), spectra (SpekPy poludniowski2021spekpy), beam hardening and the additivity gap, tone curves and
%       exactness of X' = S' + B', detector effects (blur, noise, scatter), renderer tests (tests/test_renderer.py, tests/test_compose.py),
%       orientation check, LUNA16 versus DICOM voxel check, usable-CT rules and the padded-row decision (DEVIATIONS.md, 2026-10-27).
%   Every statement in Section sec:dp-details is read from the code (s4/, tools/, kernels/, tests/); no result of the study is typed here.
This chapter first reproduces the data and methods text of the paper and then adds the details that the paper has no room for.
Section~\ref{sec:dp-details} extends the paper text below without repeating it: each of its subsections gives the implementation behind one
step of Sections~\ref{sec:data} and~\ref{sec:methods} as it stands in the code.

\section{Data}
\label{sec:data}
% PLAN (about 1 page and one table). Physics and rendering are in methods.tex; this section lists what data are used for what, under which licence,
% and which sets are locked.
%   3.1 Training sources: LIDC-IDRI (v1 renders; lesion library for the twin pastes; synthetic test CTs) and the public Qure ct2xr-projections
%       (v2 source; exactly additive components, see Section method_v2). Licence consequence: weights trained on the projections inherit CC BY-NC-SA 4.0.
%   3.2 Real evaluation data: SZCH-X-Rays (paired DES), VinDr-RibCXR (rib masks, licensed, local only), JSRT, Node21, TBX11K, VinBigData, NIH ChestX-ray14.
%   3.3 Data discipline: development versus locked sets, split rules, what was seen before the freeze (DEVIATIONS.md 2026-10-27, SAP_v2.md).
%   Table~data: dataset, role, licence or access. Licences are copied only from pages that were read; everything else stays a \todo.
% Counts below are copied from SAP_v2.md, methods.tex and DEVIATIONS.md; regenerate them from results/ when the tables are final.
Table~\ref{tab:data} lists every data set with its role in this study and its licence or access condition.

\subsection{Training sources}
\paragraph{LIDC-IDRI.} The LIDC-IDRI collection \citep{armato2011lidc,armato2015lidctcia,clark2013tcia}, with some series read from the LUNA16 copy
\citep{setio2017luna16}, supplies the v1 renders, the lesion library used to paste lesions into training images, and the synthetic test CTs. The
selection, the usable-CT rules and the patient-level split are given in Section~\ref{sec:methods}, together with the detector physics (attenuation
tables from the NIST X-ray mass attenuation coefficient database \citep{hubbell2004xraymass}, tube spectra from SpekPy \citep{poludniowski2021spekpy}).

\paragraph{Qure.ai CT projections.} The v2 models are trained on the public Qure.ai \path{ct2xr-projections} data \citep{angaitkar2026anatomy}, which
hold, for each CT volume, a composite image and translated bone, soft-tissue and lung projections. The composite is an exact linear combination of
the components up to rounding, so the bone image of every training sample is known exactly (Section~\ref{sec:method_v2}).
The release holds \num{21887} projections, one per CT volume; the volume names are those of CT-RATE \citep{hamamci2026ctrate}. We split them by
patient (\num{14374} patients) into \num{21051} training, 406 validation and 430 test projections (\path{tools/qp_prep.py}).

\subsection{Real radiographs}
\paragraph{SZCH-X-Rays.} Paired radiographs and dual-energy soft-tissue images from one partner hospital \citep{sun2026bsldm,sun2026debonedit}. We use the
75-pair test subset; its 75 images are split into a development part (24 images) and a held-out part (51 images) by a hash of the file name
(\path{tools/szch_split.py}). The 75 images are the published test split of the 741-pair version of the data set (592/74/75), acquired with dual-exposure DES on a GE
Discovery XR656 \citep{sun2026debonedit}; an earlier version of 818 pairs was used for BS-LDM \citep{sun2026bsldm}, and the papers do not
explain the difference.

\paragraph{VinDr-RibCXR.} 245 radiographs with per-rib masks \citep{nguyen2021ribcxr}. The data are licensed and are processed only on the licence
holder's own computer and private cloud account; only metrics leave them.

\paragraph{Nodule and detection sets.} JSRT \citep{shiraishi2000jsrt} (154 nodules), Node21 \citep{sogancioglu2024node21} (1\,476 nodule boxes),
TBX11K \citep{liu2020tbx11k} and VinBigData (the Kaggle competition data; related to VinDr-CXR \citep{nguyen2022vindrcxr}) are used for real-lesion contrast
and for downstream detection; NIH ChestX-ray14 \citep{wang2017chestxray} is used for downstream classification with its official lists.
NIH uses the official test list (\num{25596} images) and a patient-level split of the remaining images into \num{77890} training and \num{8634}
validation images. The detection sets are split by a fixed hash of the image identifier: Node21 \num{3440}/\num{512}/\num{930}, TBX11K
\num{6200}/\num{688}/\num{2088} (its official validation set serves as test set) and VinBigData \num{10433}/\num{1469}/\num{2943} images for
training, validation and test. VinBigData boxes are kept where at least two radiologists marked the same finding; images with only single-reader
findings are left out. No image of any set was unreadable. VinBigData is the Kaggle competition built on the VinDr-CXR images \citep{vinbigdata2020kaggle}; its rules allow non-commercial research use.

\paragraph{BDCXR-3257 (external validation).} 3,257 paediatric chest radiographs (JPEG) from the Radiology and Imaging Department of
Dr.~M~R~Khan Shishu Hospital \& Institute of Child Health, Dhaka, collected in routine care and anonymised, with one label each: normal (880),
bacterial pneumonia (1,477) or viral pneumonia (900) \citep{durjoy2026bdcxr}, a task studied before on other paediatric sets \citep{kermany2018identifying}. Most children are under five years old
(1,132 images under one year). The data were used under the ethical approval REC-FSIT-2025/No:~11928 (Daffodil International University) and
the data use agreement between the hospital and the authors' research group; no radiograph of this set is shown in this paper. There are no
patient identifiers, so near-duplicate images are grouped by their correlation at $64$~px ($r\ge0.93$; 48 pairs, the rest single images) and
a group never spans two folds.

\subsection{Development and locked sets}
SAP.md fixed the locked sets before any locked-set evaluation. After the v1 ablation the SZCH pairs were looked at (DEVIATIONS.md, 2026-10-27), so SZCH
became a development set for decisions about the v2 recipe, and SAP\_v2.md was written and frozen before any held-out or locked evaluation of a v2
model. Decisions used the 24 development images, the synthetic validation data and the real-image audit on all 75 SZCH images; the audit on SZCH is
therefore development evidence only, and the confirmatory audit runs on VinDr-RibCXR (not opened by any model run before the plan was written) and
JSRT. Paired fidelity claims use the 51 held-out SZCH pairs. Before the freeze, screening tools printed held-out PSNR and SSIM next to the
development numbers for v1-family and screening runs; no rule, table or decision used them, and the tool was changed to print development
numbers only.

\begin{table}[htbp]
\centering\small
\renewcommand{\arraystretch}{1.15}
\caption{Data sets, their role in this study and the licence or access condition, as stated on the providers' pages (read on
3 November 2026).}
\label{tab:data}
\begin{tabular}{@{}p{0.24\linewidth}p{0.35\linewidth}p{0.33\linewidth}@{}}
\toprule
Data set & Role here & Licence or access \\
\midrule
LIDC-IDRI (TCIA) & v1 renders, lesion library, synthetic test CTs & CC BY 3.0; attribution with the data DOI \\
Qure.ai ct2xr-projections & v2 training source & CC BY-NC-SA 4.0 \\
CT-RATE & CT source of the projections (not used directly) & CC BY-NC-SA 4.0, gated \\
SZCH-X-Rays & paired DES evaluation (24 development, 51 held-out) & CC BY-NC 4.0 \\
VinDr-RibCXR & locked audit and rib benchmark (245 images) & licensed; local use, metrics only \\
JSRT & audit and nodule contrast & free after registration; not public domain, no redistribution or commercial use \\
Node21 & nodule contrast, detection & CC BY-NC-ND 4.0 \\
TBX11K & detection & CC BY 4.0 (data) \\
VinBigData / VinDr-CXR & detection & competition rules: non-commercial research use \\
NIH ChestX-ray14 & classification & unrestricted use with attribution (NIH Clinical Center) \\
BDCXR-3257 & external validation (pneumonia classification, audit) & CC BY 4.0 (figshare); used under the providers' data use agreement \\
\bottomrule
\end{tabular}
\end{table}


\section{Methods}
\label{sec:methods}

\subsection{CT decomposition and projection}
\paragraph{Data.} We used the LIDC-IDRI collection (1\,018 chest CT series from 1\,010 patients) with the radiologists' nodule contours.
Series were read from DICOM or from the LUNA16 MetaImage copy (voxel-identical on all 12 series we compared, including nodule slice locations); 1\,000 series
could be rendered (18 were skipped because of non-axial orientation or irregular slice spacing). A render was \emph{usable} if the lungs were not cut by the
scan, at most 50 of the 512 frame rows lay beyond the scan, the lung volume was between 1.5 and 9~L and the bone volume exceeded 300~cm$^3$
(844 usable CTs). Patients were assigned to train/validation/test (70/10/20\%) by a hash of the patient identifier (588/98/158 CTs; 1\,057/172/298
consensus nodules) before any model was trained; the split file hash is recorded in \path{freeze.json}.

\paragraph{Components.} Each CT was split into bone, non-lung soft tissue and lung. The body is the largest component of $\mathrm{HU}>-300$
(slice-wise hole filling, then a $1{\times}9{\times}9$ opening to detach the couch). Bone follows the iterative-relaxation recipe of the reference method \citep[Algorithm~1]{angaitkar2026anatomy}, with a 10~mm support cube,
the first pass starting from the reference mask and a single closing after the ladder:
a reference mask at 300~HU, thresholds $\tau\in\{200,150,100\}$~HU admitted only where at least 5\% of a 10~mm neighbourhood was already bone, union with the
reference mask, a radius-2 closing to fill marrow, and removal of components smaller than 1\% of the largest. Lung is air ($<-320$~HU) inside the body in
components at least 20\% the size of the largest, closed and hole-filled slice-wise (so vessels and nodules belong to the lung component) and
without bone. Bone voxels were replaced in the soft-tissue component by a normalised convolution of the surrounding soft tissue (Gaussian, 4~mm, with
12 and 36~mm fall-backs). Density is water-equivalent, $\rho=\mathrm{clip}(1+\mathrm{HU}/1000,0,3)$~g\,cm$^{-3}$.

\paragraph{Projection.} Each component was integrated along the anteroposterior axis (parallel beam), giving an areal density $a_c$ (g\,cm$^{-2}$).
The frame is a body crop: width 1.08 times the body width, height the lung extent plus 60~mm clipped to 0.75--1.2 of the width (short scans cannot
reach 1.2), area-resized to $512\times512$ and later squashed to a square exactly as real radiographs are. Padded rows are set to air.
The AP view is the horizontal mirror of the PA view.

\paragraph{Detector physics.} A pixel value is
\begin{equation}
P_c=-\ln\frac{\sum_E \phi(E)\,E\,\exp\!\big(-(\mu/\rho)_c(E)\,a_c\big)}{\sum_E\phi(E)\,E},
\label{eq:physics}
\end{equation}
with ICRU-44 \citep{icru1989report44} cortical bone, soft tissue and lung attenuation (table~4 of NIST SRD 126; \citealp{hubbell2004xraymass}) and W-anode spectra (SpekPy, 12$^\circ$, 2.5~mm Al) at
70--140~kVp (training draws 80--130~kVp, the synthetic evaluation uses 100~kVp). Components are therefore stored as areal densities and the
spectrum is applied at training time, so one render serves every tube voltage.
The sum of separately projected components differs from the projection of the whole body because of beam hardening; we report this
additivity gap (6.1\% at 100~kVp for 20~g\,cm$^{-2}$ of soft tissue and 3~g\,cm$^{-2}$ of bone).

\subsection{Training images with exact components}
A training sample is composed on the GPU from the stored areal densities in this order: lesion insertion (below), one random affine transform shared by all maps
(rotation $\pm15^\circ$, scale 0.9--1.15, shift $\pm6\%$, mirror), the physics above, component weighting
$S=\mathrm{norm}(P_{\mathrm{nl}}+g\,P_{\mathrm{lung}})$, $g\sim U(1,3)$, $B=w\,P_{\mathrm{bone}}/p_{99}$, $w\sim U(0.2,0.5)$,
$X=(S+B)/1.5$, and a monotone tone curve $\gamma$. The tone curve is either a random gamma or histogram matching to the intensity quantiles of a randomly
drawn real radiograph (4\,000 NIH ChestX-ray14 training-list images). Defining $X'=\gamma(X)$, $S'=\gamma(S)$ and $B'=X'-S'$ keeps the
decomposition exact under any monotone mapping. Gaussian blur is linear and applied identically to $X'$ and $B'$; quantum noise and scatter
modify only $X'$.

\subsection{Suppressor}
The suppressor is a SwinV2-Tiny encoder (ImageNet-1k, window 8, $512^2$ input) with a nested (UNet++) decoder over the four Swin stages, CBAM attention, a plain
U-Net tail with two shallow convolutional skips to full resolution and a $1\times1$ head that outputs the bone image directly; the suppressed image is
$x-\hat b$. Relative to the reference this is a smaller backbone at half the resolution, forced by compute. Training used AdamW (lr $2\times10^{-4}$,
weight decay 0.01, 500-step warm-up, cosine decay), batch 8, mixed precision, gradient clipping at 1, EMA weights and 200\,000 samples per arm.
The loss is $0.5\,\mathrm{L1}+0.5\,\mathrm{MSE}+0.5\,(1-\mathrm{SSIM})+2.5\,(\mathrm{L1}+\mathrm{MSE})_{\mathrm{ROI}}$ with the ROI taken from the lung component
(top structures: $P_{\mathrm{lung}}/p_{99.9}>0.3$). Every arm changes exactly one factor of the recipe arm (Table~\ref{tab:arms}); all other settings,
including the seed set, are identical.

\subsection{Lesion-faithful supervision}
Each lesion (LIDC consensus of at least two readers, 50\% voxel agreement) was projected as its excess density over the surrounding lung.
\emph{Lesion-weighted loss} adds $(\mathrm{L1}+\mathrm{MSE})$ on a dilated lesion mask with weight 1. \emph{Counterfactual lesion insertion (CLIA)}
adds, before the physics, 0--3 lesion patches drawn from \emph{training} CTs at random lung positions (including over bone) with random scale
(0.7--1.5) and an amplitude of 0.6--1.6 times that scale; the bone target is untouched, so the network is explicitly taught that a lesion on a rib is not bone.

\subsection{Lesion-absorption audit}
For a model with bone estimate $\hat B$ and a lesion signal $n$ in composite-image units, the absorption is
$A=\langle \hat B(x+\alpha n)-\hat B(x),\,\alpha n\rangle/\langle\alpha n,\alpha n\rangle$ over the lesion support;
$A=0$ means the lesion stays in the soft image and $A=1$ that it was removed with the bone. Because $A=0$ is trivially reached by doing nothing,
it is always reported next to the suppression gain. On synthetic test CTs the lesions are drawn from the test CTs only (never from training CTs) and the truth is exact.
On real radiographs (SZCH test set) lesions from held-out CTs are added at random positions in the lung fields (lung masks from a published chest-X-ray
segmentation network); the bone overlap of a site is measured on the real dual-energy bone signal, never on a model output.

\subsection{Evaluation on real images}
\paragraph{Paired dual-energy data.} The SZCH-X-Rays test set has 75 radiographs with dual-energy-subtraction soft-tissue references.
Tone differs between a radiograph and its DES image, so metrics are reported as is (``raw'') and after a per-image least-squares gain and offset of the estimate
onto the reference (``calibrated'', primary); calibrated PSNR and SSIM are also reported inside the region where DES removed the most signal
(top 10\% of calibrated radiograph minus DES) and outside it.
\paragraph{Rib benchmark.} On VinDr-RibCXR (245 images with per-rib masks; evaluated locally, never used for training) we report the relative reduction of Sobel
gradient energy on the rib-boundary band inside the lungs ($R_\mathrm{rib}$), on lung pixels away from ribs ($R_\mathrm{non}$), their difference (selectivity)
and the ROC-AUC of the gradient-reduction map for separating the two sets of pixels.
\paragraph{Controls and baselines.} Identity, CLAHE, a homomorphic filter and the publicly released CT-supervised weights of the reference method
(research licence) are scored with the same code.

\subsection{Statistics}
All comparisons are paired. Confidence intervals are percentile bootstrap intervals (10\,000 resamples) over patients when several units come from one patient;
tests are Wilcoxon signed-rank tests with Holm correction within a hypothesis family; non-inferiority uses the lower bound of a one-sided 95\% interval against a
margin fixed in the analysis plan; AUCs are compared with the paired DeLong test. Seeds are averaged within unit before testing and their spread is reported.
The analysis plan, the split and the code commit were frozen before any locked evaluation set was scored (\path{freeze.json}).


\section{Details that the paper omits}
\label{sec:dp-details}
The statements in this section are read from the code of the v1 data pipeline, whose files are named in typewriter type. A CT is segmented and projected once into areal-density maps (\path{s4/ct.py},
\path{s4/seg.py}, \path{s4/render.py}); the maps of all usable CTs are loaded onto the GPU as one bank (\path{s4/bank.py}); and every
training sample is composed from the bank at run time (\path{s4/xray.py}, \path{s4/compose.py}). The real radiographs of the downstream tasks
are prepared by a separate script, described first.

\subsection{Data preparation per set}
\label{ssec:dp-prep}
The downstream sets are converted once into fixed arrays by \path{s4/prep.py}, with one loader per set. SZCH-X-Rays and VinDr-RibCXR do not
pass through it; they are read directly by their evaluation scripts (\path{s4/eval_real.py}, \path{s4/eval_ribcxr.py}).

\paragraph{Reading and resizing.} Every image is read as one grey-level array in floating point (MetaImage through SimpleITK, PNG and JPEG
through OpenCV at their full bit depth, colour converted to grey), area-resized to $512\times512$~px without keeping its aspect ratio, as the
renders are (Section~\ref{ssec:dp-geom}), and min--max normalised to $[0,1]$. This is the suppressor input $x$. The variants used downstream---the
original, the suppressed image $x-\hat b$ of the frozen ensemble and optionally the bone estimate and the suppressed image of the released model of
the reference method \citep{angaitkar2026anatomy}, run at its native 1024~px---are clipped to $[0,1]$, area-resized to the target size $s$ and
stored as 8-bit arrays of shape $N\times s\times s$: $s=256$~px for the classification sets (NIH ChestX-ray14, BDCXR-3257) and $s=512$~px for the
detection and nodule sets (TBX11K, VinBigData, Node21, JSRT), whose boxes are rescaled with separate horizontal and vertical factors. For the
detection sets the script also stores, per box, the fraction covered by strong predicted bone (above 25\% of the 99th percentile of the bone map)
and a contrast-to-noise ratio against a surrounding ring of half the box size, for the original, the suppressed and a CLAHE image
\citep{zuiderveld1994clahe}.

\paragraph{Hash splits.} All hash splits in this work, including the LIDC patient split (\path{tools/make_split.py}), map a key $k$ to
\begin{equation}
  h(k) = \mathrm{int}_{16}\big(\mathrm{SHA}_8(\texttt{20261026:}k)\big) \bmod 100,
  \label{eq:dp-hash}
\end{equation}
where $\mathrm{SHA}_8(\cdot)$ is the first 8 hexadecimal digits of the SHA-256 digest of its argument,
and assign it to training if $h<70$, to validation if $70\le h<80$ and to the test set otherwise. The key is the patient identifier where a set has
one and the image identifier otherwise, so the assignment of an image depends on its own key only and survives changes of the file list.

\paragraph{Per set.} \emph{NIH ChestX-ray14} \citep{wang2017chestxray}: 14 multi-hot finding labels; the official test list is the test set, and
the official training and validation list is split by the hash of the patient identifier into validation ($70\le h<80$) and training (all other
values). \emph{TBX11K} \citep{liu2020tbx11k}: boxes from the COCO-format files; the official training images are split by the hash of the file name
into training and validation, and the official validation images form the test set. \emph{VinBigData} \citep{vinbigdata2020kaggle}: images from a
public 1024~px copy, boxes rescaled from the original image size; for each image and finding (``no finding'' is ignored) a box joins the first
cluster whose mean box has an intersection over union above 0.4 with it and that holds no box of the same radiologist. Clusters of at least two
radiologists become one mean box, single-reader clusters are dropped, an image whose findings are all single-reader is excluded, and an image
without findings stays as a normal image; the split is by the hash of the image identifier. \emph{Node21} \citep{sogancioglu2024node21}: the
processed 1024~px images (Zenodo record 5548363), read from a private copy or fetched from the archive by HTTP range requests, with the nodule
boxes halved to 512~px and a hash split by image name; the loader stops if fewer than 99\% of the listed images are present. \emph{JSRT}
\citep{shiraishi2000jsrt}: evaluation only, so every image is a test image; each benign or malignant nodule becomes a square box with the nodule
diameter as side, centred on the marked position (0.175~mm pixels in the 2048~px original, divided by four). \emph{BDCXR-3257}
\citep{durjoy2026bdcxr}: the loader checks the three class counts; the five folds are stratified by class and grouped by near-duplicate group
(Section~\ref{sec:data}), and in each fold whole groups of the training part move to validation until it holds at least 10\% of the training-part images of each
class.

\paragraph{Unreadable images.} An image that cannot be decoded is logged and skipped, never fatal: its array rows stay zero and its name is
listed as unreadable in \path{status.json}. Images that the annotations list but that are missing on disk are logged and left out (TBX11K,
Node21).

\subsection{Segmentation recipe}
\label{ssec:dp-seg}
A CT is loaded by \path{s4/ct.py} as a volume of Hounsfield units indexed $[z,y,x]$, with $z=0$ the most superior slice, clipped to
$[-1024,3071]$~HU; the voxel spacing $(\Delta z,\Delta y,\Delta x)$ travels with it, and every rule below that is stated in millimetres is
converted to voxels with it. \path{s4/seg.py} then builds the body, bone and lung masks in this order, because each mask uses the earlier ones.
The three projected components, bone, lung and the rest of the body, are disjoint.

\paragraph{Body.} (1)~Threshold at $\mathrm{HU}>-300$; (2)~fill the holes of every axial slice; (3)~open with a $1\times9\times9$ voxel
structuring element (in-plane only), which cuts thin connections such as the couch off the body; (4)~keep the largest connected component. The
order of steps 2 and 3 matters: filling first encloses the lungs on every slice with an intact chest wall, whereas an opening first could break a
thin chest wall, and the filling would then leave the lungs out (a failure recorded in the code comment of \texttt{body\_mask}).

\paragraph{Bone.} Bone follows the iterative relaxation of \citet[Algorithm~1]{angaitkar2026anatomy}. Let $R=\{\mathrm{HU}>300\}$ be the
reference mask and $f_M(v)$ the fraction of a mask $M$ inside a cube of side 10~mm centred on voxel $v$. Starting from $M_0=R$, three steps admit
lower thresholds only next to bone already found:
\begin{equation}
\begin{gathered}
M_k = \big(\{\mathrm{HU}>\tau_k\}\cap\{v: f_{M_{k-1}}(v)\ge\theta\}\big)\cup R,\\
(\tau_1,\tau_2,\tau_3)=(200,150,100)~\mathrm{HU},
  \quad \theta=0.05 .
\end{gathered}
\label{eq:dp-bone}
\end{equation}
A voxel of 120~HU is bone only if at least 5\% of its 10~mm neighbourhood was bone after the second step, so the mask grows from the cortex into
trabecular bone without admitting isolated soft-tissue voxels of the same attenuation. To keep this affordable, $f_M$ is computed on a grid
coarsened by $\max(1,\mathrm{round}(3~\mathrm{mm}/\Delta))$ along each axis and copied back; voxels of an incomplete last block receive the
fraction zero and can enter only through $R$. The mask is then closed with a ball of radius two voxels, which fills the marrow inside the
cortical shell; the closing runs on the bounding box of the mask, edge-padded by 3 voxels because a binary closing in SciPy erodes structures
that touch the array border. Components smaller than $\kappa=1\%$ of the largest one are removed (the code calls them speckle islands), and the
renderer intersects the result with the body mask.

\paragraph{Lung.} Air inside the body, $\{\mathrm{HU}<-320\}\cap M_{\mathrm{body}}$, is split into connected components, and every component with
at least 20\% of the voxels of the largest one is kept; this keeps both lungs, connected or not, and drops small gas pockets. On every slice that
contains lung the mask is closed with a disc of radius 5 pixels and hole-filled, so that vessels and nodules enclosed by aerated lung belong to the
lung component, and bone voxels are removed. The trachea and main bronchi are connected air and therefore part of the lung mask, which matters
for the frame (Section~\ref{ssec:dp-geom}).

\paragraph{Soft tissue.} The non-lung soft tissue is $M_{\mathrm{soft}}=M_{\mathrm{body}}\setminus(M_{\mathrm{lung}}\cup M_{\mathrm{bone}})$.
If the bone voxels were simply removed from it, the soft-tissue projection would carry a negative imprint of the skeleton, as if the bones were
air. Its density inside the bone mask is therefore filled from the surrounding soft tissue by normalised convolution,
\begin{equation}
  \rho_{\mathrm{fill}}(v) = \frac{\big(G_\sigma*(\rho\,\mathbf{1}_{M_{\mathrm{soft}}})\big)(v)}{\big(G_\sigma*\mathbf{1}_{M_{\mathrm{soft}}}\big)(v)},
  \label{eq:dp-fill}
\end{equation}
with a Gaussian $G_\sigma$ of $\sigma=4$~mm where the denominator exceeds 0.01, then $\sigma=12$ and 36~mm for voxels still unfilled deep inside
large bones, and the density of water for any voxel left. The lung component keeps the measured density of its voxels, vessels and nodules
included. The density $\rho=\mathrm{clip}(1+\mathrm{HU}/1000,0,3)$~g\,cm$^{-3}$ treats every voxel as water of a different density; since HU
measure attenuation relative to water at the effective energy of the CT beam, this mixes density and composition for bone, whose composition
enters only through the cortical-bone attenuation table (Section~\ref{ssec:dp-atten}). The simplification relative to the reference method is
recorded in \path{DEVIATIONS.md} (2026-10-26).

\paragraph{Sources of the rules.} The bone recipe follows the published method \citep[Algorithm~1]{angaitkar2026anatomy} in its structure:
a reference mask above 300~HU, a threshold ladder of 200, 150 and 100~HU admitted by support voting with $\theta=0.05$, and pruning of
components smaller than 1\,\% of the largest. It differs in three details: support is counted in a 10~mm cube instead of a window of $15^3$
voxels, the first pass starts from the reference mask instead of from all voxels, and the reference method's closings at every step (balls of
radius 1 and 2) are replaced by one radius-2 closing after the ladder. The body, lung and soft-tissue steps are choices of this work: the
reference method removes the bed and clips intensities, segments the lungs with a pretrained model, and fills the bone voxels of the soft-tissue
volume by diffusion inpainting; this work uses the morphology listed above (fill-then-open order, $1\times9\times9$ opening, 20\,\% rule,
radius-5 closing), a coarse grid and a normalised-convolution fill.

\paragraph{Failure modes.} The code handles three: the fill-then-open order of the body mask, the border erosion of SciPy's closing, and
trachea-only slices at the top of the scan, which the frame does not count as lung (Section~\ref{ssec:dp-geom}). Others follow from the thresholds
and are not detected. Any voxel above 300~HU becomes bone whatever its anatomy, so calcified nodules or lymph nodes, densely contrast-filled vessels
and metal would enter the bone target; lung disease denser than $-320$~HU belongs to the lung component only where the slice-wise closing encloses
it; and a couch that the opening does not detach stays in the body. Gross failures are screened by the usable-CT rules of
Section~\ref{sec:methods}, which \path{tools/make_split.py} applies to the metadata of every render and records as named reasons in the split
file; the lungs count as cut by the scan when their extent (Section~\ref{ssec:dp-geom}) starts in the first two or ends in the last two slices.
How often the threshold-based failures occurred was not measured: no step of the pipeline counts calcified nodules, contrast or metal in the
bone target, so their effect on the v1 models is unknown. The v2 models are trained on the Qure.ai projections and do not depend on this
segmentation.

\subsection{Rendering geometry}
\label{ssec:dp-geom}
\paragraph{Orientation and sources.} The DICOM loader accepts only axial series: the orientation of the first slice must equal $(1,0,0,0,1,0)$ to
within 0.02 in every direction cosine. It sorts the slices by position, puts the head first and skips series with fewer than 80 CT slices,
duplicate positions, a spacing that deviates anywhere from the median by more than 25\%, slices of unequal size or an unreadable file; a skipped
CT is logged and never stops the run. With this orientation the row index $y$ runs from anterior to posterior and the column index $x$ from the
patient's right to the patient's left, so a projection along $y$ shows the head at the top and the patient's right on the viewer's left, the
display convention of a posteroanterior radiograph.

Each series is read from the mounted DICOM copy (606 of the 1\,018 series, \path{DEVIATIONS.md}, 2026-10-26), otherwise from the LUNA16
MetaImage copy \citep{setio2017luna16}, otherwise downloaded from TCIA (\path{kernels/render/render_shard.py}); the source is recorded with
each render. The MetaImage loader accepts uncompressed 16-bit volumes only and locates a nodule slice from its $z$ position in millimetres, whereas
the DICOM loader matches the SOP instance identifier of the annotation. The agreement of the two copies was checked by
\path{kernels/orient_check/check.py} on 12 series spread over those present in both: shapes, spacings, the mean absolute HU difference for all
eight combinations of axis flips and the first slice of every consensus nodule located either way. The outcome is given in
Section~\ref{sec:methods}.

\paragraph{Projection and the AP view.} Each component is projected by a parallel-beam sum along the anteroposterior axis,
\begin{equation}
  a_c[z,x]=\frac{\Delta y}{10}\sum_y\rho_c[z,y,x]\quad(\mathrm{g\,cm^{-2}},\ \Delta y\ \text{in mm}),
  \label{eq:dp-project}
\end{equation}
which is the same for a beam entering from the front (AP) or the back (PA); the raw images of the two views differ by a horizontal mirror. The
renderer stores the PA convention, and the composer draws a horizontal mirror with probability 0.5 inside the affine transform that it shares
between all maps (\texttt{random\_affine} in \path{s4/compose.py}), so that both handednesses occur in training. Divergent-beam magnification,
which does differ between AP and PA, is not modelled.

\paragraph{Frame.} The frame (\texttt{frame\_box} in \path{s4/render.py}) is a box in the $[z,x]$ plane of the projection:
\begin{itemize}
\item Lung extent $L$: from the first to the last slice whose lung cross-section exceeds 8~cm$^2$. Trachea-only slices must not count, or every
scan that starts in the neck would look cut at the apex and the frame would be centred too high.
\item Width: $W=1.08$ times the left--right extent of the body mask over all slices, centred on the body.
\item Height: $H=\mathrm{clip}(L+60~\mathrm{mm},\,0.75\,W,\,1.2\,W)$, centred on the lungs.
\item Placement: if the scan is at least $H$ long, the frame is slid along $z$ until it lies inside the scan, so that no row is padded; otherwise
it stays centred and the rows beyond the scan are padded.
\end{itemize}
The code comment calls $0.75\le H/W\le1.2$ the aspect range of real radiographs and notes that short LIDC scans cannot reach the upper end.
No source for this range was found; it is a design choice of this work.
The box is area-resized to $512\times512$ pixels, so one pixel covers $W/512$~mm horizontally and $H/512$~mm vertically: the frame is squashed to
a square exactly as real radiographs are squashed by \path{s4/prep.py}, and a render and a radiograph of the same aspect ratio carry the same
distortion. The two-dimensional lung and body masks and the lesion signals are framed in the same way with zero padding; the background of a lesion
signal is the median density of a three-voxel ring around the nodule inside the lung and outside bone.

\paragraph{Padded rows.} Where the frame extends beyond the scan, the renderer repeats the first or last projected row (edge padding) and records
the number of padded frame rows at the top and bottom, $\mathrm{round}(512\,p/(z_1-z_0))$ for $p$ padded slices of a frame spanning $z_1-z_0$
slices. Edge padding turns the last projected row into vertical streaks. On 2026-10-27 two rules changed together (\path{DEVIATIONS.md}): of the
renders, 38\% needed more than the original limit of 25 padded rows because their scans are shorter than $0.75\,W$, so the usability limit was
raised to 50 rows in total (10\% of the frame); and the padded rows are set to air, zero areal density in all three components, when the renders
are loaded into the bank (\path{s4/bank.py}), instead of keeping the streaks. The renders on disk are unchanged. A row of air has no bone, so its
bone target stays zero under every tone curve (Section~\ref{ssec:dp-gap}).

\subsection{Attenuation and spectra}
\label{ssec:dp-atten}
\paragraph{Materials and tables.} The bone, non-lung soft-tissue and lung components are given the cortical bone, soft tissue and lung tissue of
ICRU Report~44 \citep{icru1989report44}. Cortical bone differs from the other two by its calcium and phosphorus content, which raises its
photoelectric absorption at low energy; soft tissue and lung tissue have nearly the same composition, and their tabulated mass attenuation
coefficients differ by less than 2\% over the energy grid used here, so the lung and soft-tissue components differ almost only through density.
Table~\ref{tab:icru44} lists the compositions as tabulated by NIST \citep{hubbell2004xraymass}. The renderer itself does not use the nominal
densities: it takes the density of every voxel from the CT (Section~\ref{ssec:dp-atten}), and the composition enters only through the
attenuation table of each material.
\begin{table}[htbp]
\centering\small
\caption{ICRU-44 compositions used through the NIST attenuation tables: density and elemental weight fractions (NIST SRD 126, Table~2).}
\label{tab:icru44}
\begin{adjustbox}{max width=\linewidth}\begin{tabular}{@{}lccccccccccccc@{}}
\toprule
Material & $\rho$ (g/cm$^3$) & H & C & N & O & Na & Mg & P & S & Cl & K & Ca \\
\midrule
Cortical bone & 1.920 & 0.034 & 0.155 & 0.042 & 0.435 & 0.001 & 0.002 & 0.103 & 0.003 & -- & -- & 0.225 \\
Soft tissue & 1.060 & 0.102 & 0.143 & 0.034 & 0.708 & 0.002 & -- & 0.003 & 0.003 & 0.002 & 0.003 & -- \\
Lung tissue & 1.050 & 0.103 & 0.105 & 0.031 & 0.749 & 0.002 & -- & 0.002 & 0.003 & 0.003 & 0.002 & -- \\
\bottomrule
\end{tabular}\end{adjustbox}
\end{table}
The coefficients $(\mu/\rho)_c(E)$ are the NIST tabulations for these materials (Table~4 of the NIST X-ray mass attenuation coefficient pages,
which rest on the XCOM cross sections \citep{berger2010xcom}).
These pages are the NIST X-ray mass attenuation coefficient tables \citep{hubbell2004xraymass}, not the XCOM program.
\path{tools/build_xray_tables.py} downloads the three pages, keeps them in \path{s4/data/nist_raw/} for provenance, keeps the rows between 8
and 200~keV, separates the duplicate energies at absorption edges by $10^{-6}$~keV and interpolates $\ln(\mu/\rho)$ linearly in $\ln E$ onto a
1~keV grid from 10 to 150~keV. The result is committed as \path{s4/data/xray_tables.json}, so no run depends on the NIST server. On this grid
the ratio of the bone to the soft-tissue coefficient is about 3.5 at 30~keV, 1.5 at 60~keV and 1.1 at 100~keV, which is why bone contrast falls
with tube voltage.

\paragraph{Spectra.} Tube spectra come from SpekPy \citep{poludniowski2021spekpy} for a tungsten anode at $12^\circ$, with 1~keV bins, the
\texttt{kqp} physics option and 2.5~mm of aluminium filtration, at 60 to 140~kVp in 10~kV steps. Each spectrum is interpolated onto the
10--150~keV grid and normalised to unit sum; the SpekPy version is recorded in the provenance field of the table file. Photons below 10~keV,
which the aluminium filter makes negligible, are not represented.

\paragraph{Detector and lookup tables.} The detector is energy-integrating, so each photon is weighted by its energy. With the normalised weight
$p(E)=\phi(E)\,E/\sum_{E'}\phi(E')\,E'$, the pixel value of component $c$ at areal density $a$ is
\begin{equation}
  P_c(a) = -\ln\sum_E p(E)\,\exp\!\big(-(\mu/\rho)_c(E)\,a\big),
  \label{eq:dp-pixel}
\end{equation}
as in Section~\ref{sec:methods}. For each voltage and material, $P_c$ is tabulated at 8\,192 points from 0 to 80~g\,cm$^{-2}$ and interpolated
linearly on the GPU, with larger areal densities clamped to the end of the table. The renders store areal densities, so one render serves every
voltage. The composer holds tables for 70 to 140~kVp; training draws one of 80, 90, \ldots, 130~kVp per sample with equal probability
(\path{s4/train.py}), and validation uses 100~kVp.

\subsection{Additivity gap and tone curves}
\label{ssec:dp-gap}
\paragraph{Why the components are not additive.} For the whole body, with areal densities $\mathbf a=(a_c)_c$ along a ray, the pixel value is
\begin{equation}
  P(\mathbf a) = -\ln\sum_E p(E)\,\exp\!\Big(-\sum_c(\mu/\rho)_c(E)\,a_c\Big),
  \label{eq:dp-whole}
\end{equation}
while the composer uses the separate $P_c(a_c)$ of Equation~\eqref{eq:dp-pixel}. For a monochromatic beam at $E_0$ both are linear and
$P(\mathbf a)=\sum_c(\mu/\rho)_c(E_0)\,a_c=\sum_cP_c(a_c)$: the components add exactly. For a polychromatic beam they do not. Writing
$p_a(E)\propto p(E)\exp(-(\mu/\rho)_c(E)\,a)$ for the spectrum after $a$~g\,cm$^{-2}$ of material $c$,
\begin{equation}
\begin{gathered}
P_c'(a)=\mathrm{E}_{p_a}\big[(\mu/\rho)_c\big]>0,\\
P_c''(a)=-\mathrm{Var}_{p_a}\big[(\mu/\rho)_c\big]\le0 ,
\end{gathered}
\label{eq:dp-concave}
\end{equation}
so each further gram per square centimetre attenuates less than the previous one, because the low-energy photons, which are attenuated most, go
first: beam hardening. The same mechanism couples the components: inside the body, the beam that reaches the bone has already been hardened by
the surrounding soft tissue, and the bone attenuates it less than it would attenuate the unhardened beam.

\emph{Claim.} If every $(\mu/\rho)_c$ is non-increasing in $E$ on the support of $p$, then $P(\mathbf a)\le\sum_cP_c(a_c)$: the sum of separately
projected components overestimates the attenuation of the whole body.

\emph{Proof.} The functions $f_c(E)=\exp(-(\mu/\rho)_c(E)\,a_c)$ are non-negative and non-decreasing. For two such functions $f,h$, Chebyshev's sum
inequality with weights $p$ gives $\sum_Ep\,f\,h\ge(\sum_Ep\,f)(\sum_Ep\,h)$; a product of non-negative non-decreasing functions is again of this
kind, so induction over the components gives $\sum_Ep\prod_cf_c\ge\prod_c\sum_Ep\,f_c$, and applying the decreasing function $-\ln$ to both sides
yields the claim.

The condition holds for the three tables of Section~\ref{ssec:dp-atten}, whose coefficients decrease monotonically on the 10--150~keV grid,
above the absorption edges of the constituent elements. The size of the gap for one representative ray is given in Section~\ref{sec:methods}, and
\path{tests/test_renderer.py} checks it against a fixed gate.
Figure~\ref{fig:additivity_gap} shows the gap over the range of the renders: it reaches about 8\,\% of the joint projection at 80~kVp, falls
with the tube voltage, and grows with the bone areal density.
\begin{figure}[htbp]
\centering
\safefig{figures/additivity_gap.pdf}
\caption{Beam-hardening gap of separately projected components, $(P_{\mathrm{bone}}+P_{\mathrm{soft}}-P)/P$, against the soft-tissue areal
density for three bone areal densities and three tube voltages (attenuation tables and spectra of the renderer; \texttt{tools/paper\_extras.py}).}
\label{fig:additivity_gap}
\end{figure}

\paragraph{What the composer does with the gap.} The composer never forms $P(\mathbf a)$. With per-sample quantiles $q_\alpha$ it adds weighted
component images,
\begin{equation}
\begin{gathered}
S=\frac{1}{1.5}\,\frac{s-q_{0.01}(s)}{q_{0.995}(s)-q_{0.01}(s)},\quad s=P_{\mathrm{nl}}+g\,P_{\mathrm{lung}},\\
B=\frac{1}{1.5}\,\frac{w\,P_{\mathrm{bone}}}{q_{0.99}(P_{\mathrm{bone}})},\\
X=S+B,
\end{gathered}
\label{eq:dp-compose}
\end{equation}
with $g\sim U(1,3)$, $w\sim U(0.2,0.5)$ and $1.5=1+\max w$ keeping $X$ within about $[0,1]$. The input is thus exactly additive in the logarithmic
domain by construction, and the bone target is exactly its bone term. The gap separates the composed input from a physically exact radiograph of
the same CT; it is not an error of the target. The design trades physical exactness of the input, which the random weights and the quantile
normalisation loosen further, for exactness of the target.

\paragraph{Tone curves.} A clinical radiograph is not the log-attenuation image itself: detector response, vendor processing and display apply a
monotone grey-level transform. The composer draws one transform $\gamma$ per sample (\texttt{apply\_tone} in \path{s4/compose.py}), piecewise
linear between knots $(\xi_j,\eta_j)$ and constant outside them:
\begin{itemize}
\item \emph{Gamma curve}: $\xi_j=1.5\,j/64$ for $j=0,\ldots,64$ and $\eta_j=\xi_j^{\,e}$ with $e\sim U(0.7,1.4)$.
\item \emph{Histogram matching}: $\xi_j$ are the 257 quantiles of $X$ at probabilities $j/256$ and $\eta_j$ the 257 stored quantiles of a real
radiograph drawn from the reference set of Section~\ref{sec:methods}, both plus a ramp from 0 to $10^{-5}$ that makes the knots strictly increasing.
\end{itemize}
In both cases $\eta_j$ is non-decreasing, so $\gamma$ is non-decreasing on the whole real line; the curve fitted to $X$ is also applied to $S$.
Training uses histogram matching unless an arm says otherwise; validation uses the gamma curve without detector effects or lesion insertion, so
that all arms are scored on the same images.

\paragraph{Exactness under the tone curve.} The composer sets
\begin{equation}
\begin{gathered}
X'=\gamma(X),\\
S'=\gamma(S),\\
B'=X'-S' .
\end{gathered}
\label{eq:dp-tone}
\end{equation}
\emph{Claim.} For every non-decreasing $\gamma$: (i) $X'=S'+B'$ exactly; (ii) $B'\ge0$; (iii) $B'=0$ wherever $B=0$; (iv) if $\gamma$ is
continuous and piecewise linear with largest slope $\lambda$, then wherever $B>0$, $B'=\bar\gamma'B$ with
$\bar\gamma'=B^{-1}\int_S^{S+B}\gamma'(t)\,\mathrm dt\in[0,\lambda]$.

\emph{Proof.} (i) holds by the definition of $B'$, for any $\gamma$. (ii) $P_{\mathrm{bone}}\ge0$ because $P_c(0)=0$ and $P_c'>0$
(Equation~\eqref{eq:dp-concave}); $w>0$ and the quantile in the denominator is bounded below by $10^{-6}$ in the code, so $B\ge0$, $X\ge S$ in
every pixel and $\gamma(X)\ge\gamma(S)$. (iii) $B=0$ gives $X=S$, hence $B'=0$. (iv) A continuous piecewise-linear $\gamma$ is absolutely
continuous, so $\gamma(S+B)-\gamma(S)=\int_S^{S+B}\gamma'(t)\,\mathrm dt$ with $0\le\gamma'\le\lambda$ almost everywhere.

Two consequences matter. First, $B'$ is in general not $\gamma(B)$ but the bone signal times the mean slope of the curve over the grey levels the
bone spans in that pixel, so the same bone appears stronger where the curve is steep; it is exactly the quantity whose subtraction leaves the
tone-mapped soft tissue, $X'-B'=\gamma(S)$, which $\gamma(B)$ would not be for a non-linear curve. Second, by (iii) the target is zero wherever the
bank holds no bone, the padded rows included.

\paragraph{Detector effects.} With detector effects on, the default for training, three effects follow the tone curve. Let $X'_b=G_\sigma*X'$
and $B'_b=G_\sigma*B'$, with $G_\sigma$ a separable Gaussian (replicate padding) whose $\sigma\sim U(0,1.2)$~px is drawn per sample with
probability 0.5 and applied when $\sigma>0.05$~px. Then
\begin{equation}
\begin{gathered}
X'_{\mathrm{det}} = X'_b + u_1\big(0.004+0.02\sqrt{\max(X'_b,0)}\big)\,\varepsilon + 0.08\,u_2\,\big(G_{20}*X'_b\big),\\
B'_{\mathrm{det}} = B'_b ,
\end{gathered}
\label{eq:dp-detector}
\end{equation}
with $\varepsilon\sim\mathcal N(0,1)$ per pixel and $u_1,u_2\sim U(0,1)$ per sample: noise whose standard deviation grows with the square root of
the intensity above a floor, and a low-frequency copy of the image (Gaussian of 20~px) with up to 8\% of its amplitude as scatter. Only the blur
touches the bone target, and because it is linear and identical on $X'$ and $B'$, $X'_b-B'_b=G_\sigma*S'$ and the decomposition stays exact.
Noise and scatter are added to the input only, so the image the network should leave after subtracting its bone estimate contains them: it is
taught to keep noise and scatter in the suppressed image. The loss uses the bone target only; the returned soft-tissue image is $S'$ before the
detector effects and serves validation, which runs without them. The code names no source for the noise and scatter constants; they are
augmentations of plausible form, not a calibrated detector model.

\subsection{Tests}
\label{ssec:dp-tests}
The pipeline has 32 unit tests in eight files under \path{tests/}. None needs data, a GPU or a network connection: every test builds its own
input (a CT phantom, a minimal LIDC annotation file, a MetaImage volume, a small fake projection pack or random scores), and the only repository
file they read is the committed \path{s4/data/xray_tables.json}. They run on the CPU from the repository root with
\begin{verbatim}
python -m pytest tests -q
\end{verbatim}
and need NumPy, SciPy, scikit-image, OpenCV, pydicom, PyTorch, timm, scikit-learn and pytest; \path{pytest.ini} sets the test path and writes
temporary files to \path{outputs/pytest_tmp}.

\paragraph{\texttt{test\_renderer.py} (7 tests).} A phantom of $120\times160\times200$ voxels ($2.5\times1.5\times1.5$~mm) with an elliptic body of
40~HU, two elliptic lungs of $-850$~HU, a spine cylinder of 600~HU and four rib rings of 700~HU gives closed-form answers. The tests check that
(1)~the bone mask holds more than 99\% of the voxels above 300~HU and none below $-500$~HU; (2)~the lung volume is within 15\% of the analytic
one; (3)~the projection of a cylinder of 1.6~g\,cm$^{-3}$ equals chord length times density to 8\%; (4)~component pixel values add exactly (to
$10^{-3}$) for a 60~keV monochromatic beam, and at 100~kVp the relative gap for 20~g\,cm$^{-2}$ of soft tissue and 3~g\,cm$^{-2}$ of bone stays
below a gate of 0.15 fixed in advance (the test prints the value; the paper reports it); (5)~every lookup table starts at zero and is strictly
increasing; (6)~for a nodule of 30~HU in lung of $-850$~HU the background density is 0.15~g\,cm$^{-3}$ to 0.03 and the peak excess areal density
is within 12\% of its closed form; (7)~a full render gives three finite, non-negative $512\times512$ maps with bone and uncut lungs, and a volume of
air raises the skip exception.

\paragraph{\texttt{test\_compose.py} (6 tests).} (1)~The GPU lookup agrees with the NumPy evaluation of Equation~\eqref{eq:dp-pixel} to
$2\times10^{-4}$; (2)~the affine transform without rotation, scaling, shift and mirror is the identity; (3)~lesion insertion never changes the
bone or non-lung soft-tissue map and adds lesion signal and mask; (4)~$X'=S'+B'$ holds to $10^{-6}$ with $B'\ge-10^{-6}$ at 80, 100, 120 and
140~kVp; (5)~histogram matching gives the composed image the median of the reference to 0.02 and keeps $B'$ non-negative; (6)~the suppressor runs
forward and backward at 256~px with finite loss and gradients, and the SSIM of an image with itself is 1.

\paragraph{\texttt{test\_io.py} (5) and \texttt{test\_qp.py} (2).} On LIDC annotation files written by the test, the series identifier is read,
two readers form one consensus nodule with the right reader count, malignancy and slice, one reader is not a consensus, and of three readers the
distant one is left out; a synthetic MetaImage volume is put head first, located by $z$ position, and rejected if shorter than 80 slices. On a fake
pack in the format of the v2 source (Section~\ref{sec:method_v2}), $X'=S'+B'$ holds to $10^{-5}$ with a non-negative bone image, lesion insertion
changes the input while the bone target keeps more than 90\% of its sum, and with detector effects on, $X'-B'-S'$ is no longer zero, because noise
and scatter reach the input only.

\paragraph{Statistics tests (12).} The other four files check the evaluation code on constructed cases: the paired DeLong test
\citep{delong1988auc} against the AUCs of scikit-learn and its rejection rate under the null (\path{test_delong.py}); FROC, average precision
and one-to-one box matching on known boxes (\path{test_froc.py}); the Holm correction \citep{holm1979simple}, clustered bootstrap intervals,
the Wilcoxon test \citep{wilcoxon1945individual} and the non-inferiority rule (\path{test_stats.py}); and conformal coverage and sparsification
error of the uncertainty estimates (\path{test_uncertainty.py}).

The tests cover the steps whose errors would silently change the supervision. Beyond the single phantom render they do not check the frame placement
and padded-row counts, the per-set loaders of \path{s4/prep.py} or the GPU bank; these are checked only through the counts and logs that their
runs print.
Two parts of the pipeline have no unit test: the frame box (padding counts on a short scan) and the VinBigData consensus rule.


\chapter{The v1 study: lesion-faithful supervision from LIDC-IDRI}\label{ch:v1}
% The v1 study is thesis-only. Numbers come from results/ablation_val and results/ablation_test (tools/aggregate_ablation.py, tables written by
% tools/paper_extras.py), from DECISIONS.md and DEVIATIONS.md (2026-10-27), and from results/final/stats_held.json.
\section{Question and analysis plan}
The first study asked whether a suppressor trained only on CT renders of the public LIDC-IDRI collection keeps lesions, and whether showing
it lesions during training helps. The analysis plan (SAP.md) was written and frozen with \path{freeze.json} on 27 October 2026, before any
locked set was evaluated. Its first hypothesis is synthetic, with exact truth:
\begin{description}
\item[H1] On synthetic test CTs, a lesion-aware arm absorbs less of an inserted lesion than the recipe arm (per-CT mean of $A$, paired,
 patient-cluster bootstrap interval, Holm over the lesion-aware arms), and its bone-image MAE is non-inferior within a margin of 5\,\% of the
 recipe's mean MAE (one-sided 5\,\%).
\end{description}
The same rule, applied to synthetic \emph{validation} CTs, selected the v1 model: an arm qualifies only if its absorption is lower with Holm
$p<0.05$ and its bone MAE is non-inferior. SAP.md locked the real sets: SZCH-X-Rays (75 dual-energy pairs), JSRT, Node21, the TBX11K
validation set and VinDr-RibCXR were never to be used for training, tuning or selection; the validation parts of NIH ChestX-ray14 and
VinBigData could inform pilots. Real-image results could never choose a model.

\section{Design}
Each arm changes one factor of the recipe of Chapter~\ref{ch:data} (the training recipe of the reference paper, re-implemented): the
lesion-aware arms add counterfactual lesion insertion (CLIA, up to three lesions from training CTs pasted into the composite), a loss term
weighted on lesion pixels, or both; the factor arms remove the frequency loss, the histogram-matching augmentation, the marking-ROI loss or
the rotation augmentation. Every arm was trained with three seeds for \num{200000} samples, at the learning rate fixed by a pilot on
synthetic validation CTs ($2\times10^{-4}$; Appendix~\ref{app:deviations}). CTs were split by patient (70/10/20) with a fixed hash; 844 CTs met
the usable-CT rules, and the selection used 98 validation CTs with 177 insertion sites, the test 158 CTs with 317 sites. Lesions for the audit
come from other CTs of the same split, never from training CTs.

\section{Outcome of the selection rule}
On the validation CTs (Table~\ref{tab:v1_ablation}) every lesion-aware arm absorbed less than the recipe arm (CLIA by 0.033, the
lesion-weighted loss by 0.036, the combination by 0.100), but only CLIA kept its bone MAE within the margin: the lesion-weighted loss sat on
the margin ($-0.00036$ against $-0.00036$) and the combination was about 6\,\% worse. By the pre-declared rule CLIA became the v1 model, and the
combination was kept as a pre-declared trade-off variant. Among the factor arms, removing the marking-ROI loss raised absorption (by 0.024); removing
histogram matching lowered both absorption (by 0.039) and the synthetic bone error, but that augmentation targets the appearance of real
radiographs, which synthetic validation does not measure; removing the frequency loss or the rotation changed absorption little.

On the test CTs (Table~\ref{tab:v1_h1}) H1 passes: CLIA absorbs 0.022 less than the recipe arm $[0.013,0.032]$, Holm $p<0.001$, with
non-inferior bone MAE. The lesion-weighted loss also lowers absorption (0.019) and is non-inferior on the test CTs; the combination lowers it most
(0.066) and misses the margin again. The 300k-sample finals of CLIA and recipe, trained later with three seeds each, replicate the contrast.
The Holm adjustment in these tables runs over all seven contrasts of the ablation, which is more conservative than the three named in the plan.
\begin{table}[htbp]
\centering\small
\caption{v1 ablation on synthetic validation CTs (selection data; three seeds each, 200k samples): every arm against the recipe arm.}
\label{tab:v1_ablation}
\fittable{tables/v1_ablation_val.tex}
\end{table}
\begin{table}[htbp]
\centering\small
\caption{First analysis plan, H1: lesion-aware arms against the recipe arm on synthetic LIDC test CTs (three seeds each), and the 300k-sample
finals as a replication.}
\label{tab:v1_h1}
\fittable{tables/v1_h1.tex}
\end{table}

\section{Real-image check and the resulting deviation}
After the ablation the v1 arms were compared with the released weights of the reference model on the 75 SZCH-X-Rays pairs, the comparison
that SAP.md had locked for the final evaluation. The v1 models had the higher bone-region PSNR$_{cal}$ (22.04 against 21.66\,dB) and the lower
full-resolution SSIM$_{cal}$ (0.850 against 0.869), and in the real-image audit on all 75 images they absorbed far more of an inserted lesion
($A$ 0.306 against 0.134 over 450 sites). Because these images had now been looked at, SZCH-X-Rays could no longer serve as a clean locked set
for decisions about a new recipe. It became a development set: 24 images, chosen by a hash of the file name (\path{tools/szch_split.py}),
were used for paired-metric decisions from then on, the other 51 were held out for final fidelity claims, and the audit on all 75 images became
development evidence only. The confirmatory audit moved to VinDr-RibCXR and JSRT, which no model had touched. Appendix~\ref{app:deviations}
lists this deviation with the others.

\section{Consequences for the v2 design}
The v1 study answered its own question: on rendered data with exact truth, lesion insertion lowers absorption a little at no cost in bone
error. It did not solve the problem it was meant to solve. On real radiographs the v1 models absorbed more than twice what the reference model
absorbed, and the absorption hardly depended on the lesion amplitude (Chapter~\ref{ch:audit}): it behaved like a uniform leak of lesion
signal into the bone estimate rather than a failure on hard cases. Three changes followed. The training source moved to the Qure.ai CT
projections, which are translated towards the appearance of real radiographs and are exactly additive (Chapter~\ref{ch:v2}). Lesion insertion,
which had acted only through the supervised target, was replaced by a term that ties the predictions with and without the lesion directly.
And a second analysis plan, SAP\_v2.md, was written before any held-out or locked evaluation of the new models. On the final locked sets the v1
model still absorbs 0.289 of an inserted lesion on VinDr-RibCXR, against 0.149 for the reference model and 0.051 for the v2 ensemble
(Chapter~\ref{ch:audit}).


\chapter{The v2 method: Qure projections and the twin regulariser}\label{ch:v2}
% v2 chapter. The method text of the paper is input; thesis-only parts are the design rationale, the exploratory screening (numbers copied from
% DECISIONS.md, all development data) and the thesis-length reading of the dose-response, ablation and encoder results (generated tables).
\section{Why a second analysis plan}
The first study left three things open (Chapter~\ref{ch:v1}). Its models kept lesions on rendered data but not on real radiographs; lesion
insertion acted only through the supervised target, so the network was never told that the bone estimate should not change when a lesion is
added; and the comparison with the reference model on SZCH-X-Rays had been looked at, so that set could no longer decide anything about a new
recipe. The second plan, SAP\_v2.md, was written on 29 October 2026 after the new recipe had been developed on development data only: the 24
SZCH development pairs, the audit on all 75 SZCH images, and the synthetic validation data. It names five confirmatory hypotheses on sets that no
decision had used: the dose-response of absorption on the twin weight on VinDr-RibCXR (H1$'$), non-inferior fidelity on the 51 held-out SZCH
pairs (H2$'$), lower absorption than the reference model and the no-twin control on VinDr-RibCXR (H3$'$), evidence that does not use insertion
(H4$'$: real-nodule contrast and downstream tasks), and reliability (H5$'$). The plan, the configuration and the weights were frozen on
1~November 2026 (\path{freeze_v2.json}) before any of these sets was evaluated.

\section{Training on public CT projections with a twin regulariser}
\label{sec:method_v2}
% PLAN (about 1.5 pages). This section is new text for the v2 recipe; methods.tex (v1 recipe, physics, audit definition) is untouched.
%   4.1 Qure projections as an exactly additive source: what a row contains, why the bone image is exact (composite = a*soft + b*bone + c*lung + d),
%       the augmentation order of s4/qp.py, and why monotone tone curves keep the decomposition exact.
%   4.2 Twin regulariser: definition, loss, the weight w, and the variants that form the ablation (stop-gradient, support size, whole image, number
%       of pasted lesions, insertion without the tie term). Relation to the v1 counterfactual insertion (CLIA) in methods.tex.
%   4.3 Ensemble and uncertainty: K=3 members, mean bone image, spread, split-conformal scale, what is calibrated where (SAP_v2.md, H5').
%   4.4 Training budget and what is identical across arms (same optimiser and schedule as methods.tex).
% Verify every formula below against s4/train.py, s4/qp.py and s4/uncertainty.py at the freeze commit before submission.
\subsection{Qure projections as an exactly additive source}
Each row of the Qure.ai data holds a composite radiograph $T$ and the translated bone, soft-tissue and lung projections $B_t$, $S_t$ and $L$
(8-bit, $512\times512$, one sample per CT volume at 100~kVp). We verified that $T=a\,S_t+b\,B_t+c\,L+d$ holds up to rounding, with a per-sample bone
weight $b$ (\path{tools/qp_prep.py}); the additive bone image of the normalised input $X=T/255$ is therefore $B=b\,B_t/255$ and the remainder is
$S=X-B$. Over all \num{21887} rows the least-squares fit leaves a mean absolute residual of 0.36 grey levels per projection (largest 3.2 of 255).
The rows are split by patient into \num{21051} training, 406 validation and 430 test projections.

Training samples are composed as in Section~\ref{sec:methods}: bone re-mixing ($X'=S+wB$, exact), one random affine transform shared by all maps,
lesion insertion into $S$ with the bone target untouched, a monotone gamma tone curve ($X'=g(S+B)$, $S'=g(S)$, $B'=X'-S'$), linear blur applied to
$X'$ and $B'$, and noise and scatter applied to $X'$ only. Defining $B'=X'-S'$ keeps the decomposition exact under any monotone tone curve.

\subsection{Twin lesion-invariance regulariser}
Let $f$ be the suppressor, $x$ a training image and $b$ its exact bone image. A \emph{twin} of $x$ is $x^{+}=\min(x+n,1)$, where $n\ge0$ is the
signal of a few lesion patches taken from the lesion library of \emph{training} CTs and pasted at random positions in the lung fields, in composite
units and after all augmentation. The bone target of $x^{+}$ is the same $b$. Let $M$ be the lesion support dilated by $3$~px. The loss is
\begin{equation}
\begin{gathered}
\mathcal{L}=\tfrac12\big[\mathcal{L}_{\mathrm{sup}}(f(x),b)+\mathcal{L}_{\mathrm{sup}}(f(x^{+}),b)\big]+w\,\mathcal{L}_{\mathrm{inv}},\\
\mathcal{L}_{\mathrm{inv}}=\frac{\sum_p M_p\,\big|f(x^{+})_p-f(x)_p\big|}{\max\big(1,\sum_p M_p\big)},
\end{gathered}
\label{eq:twin}
\end{equation}
where $\mathcal{L}_{\mathrm{sup}}$ is the suppression loss of Section~\ref{sec:methods} and $w$ is the twin weight. Half of every batch consists of twins, so
the cost per step does not change. The supervised term already asks for the unchanged bone image on $x^{+}$; $\mathcal{L}_{\mathrm{inv}}$ ties the
two predictions directly on the lesion support. Each twin carries between one and two pasted lesions (uniform), each with a random scale and
amplitude; the supervised term is averaged over all base images and twins of the batch, and $\mathcal{L}_{\mathrm{inv}}$ is normalised per image by
its support and then averaged over the twins.

The ablation (Section~\ref{sec:experiments}) changes one part at a time: the twin weight $w$; a stop-gradient on $f(x)$; a support dilated further with a
$31\times31$ window, or the whole image, instead of $M$; up to four instead of up to two pasted lesions; and lesion insertion into the base image
without the tie term (the counterfactual insertion of Section~\ref{sec:methods}). All of these arms are in the ablation of
Section~\ref{sec:res_twin} (development data); the twin weights form the confirmatory dose-response.

\subsection{Deep ensemble and uncertainty}
The reported model is a deep ensemble of $K=3$ suppressors trained with seeds $0$, $1$ and $2$ \citep{lakshminarayanan2017ensembles}; seed $3$ is a reserve.
The bone image is the ensemble mean and the soft image is $x$ minus that mean. The ensemble standard deviation is the per-pixel spread. A split-conformal
scale $q$ is the $\lceil (n+1)(1-\alpha)\rceil/n$ quantile of $|b-\bar f|/(\sigma+\varepsilon)$ pooled over calibration pixels of the Qure-projection
validation CTs, so that $[\bar f-q(\sigma+\varepsilon),\,\bar f+q(\sigma+\varepsilon)]$ has nominal coverage $1-\alpha=0.9$
\citep{lei2018conformal,angelopoulos2023conformal}. Calibration treats pixels as exchangeable, so the guarantee is marginal over pixels, not
over images. The sparsification curve is the mean absolute error of the pixels that remain after removing the fraction $f\in[0,0.95]$ with the
largest spread; AUSE is the area between this curve and the oracle curve that removes pixels by their error (0 is perfect).

\subsection{Training budget}
The suppressor keeps the nested decoder and CBAM attention of Section~\ref{sec:methods} \citep{zhou2018unetpp,woo2018cbam}. The dose-response and
ablation arms use the SwinV2-T encoder of that section \citep{liu2022swinv2}; ours uses a ConvNeXt-V2-T encoder \citep{woo2023convnextv2}
(FCMAE pre-training, ImageNet-22k then ImageNet-1k fine-tuning \citep{deng2009imagenet}; \texttt{timm} \citep{wightman2019timm}
\path{convnextv2_tiny.fcmae_ft_in22k_in1k}) with a learning
rate $0.25\times$ that of the decoder. Every v2 run is trained at $512$~px with AdamW \citep{loshchilov2019adamw} (learning rate $2\times10^{-4}$,
batch 8), the schedule and loss of Section~\ref{sec:methods}, and \num{300000} samples (\num{37500} steps); evaluation uses the EMA weights. The
encoder screen uses \num{100000} samples (Section~\ref{sec:experiments}). The seeds per arm are listed in the generated tables.


\section{Exploratory screening rounds}
The operating point of the regulariser was found in a sequence of screening rounds. All of them are exploratory: single runs or few seeds, the
SZCH development pairs and the audit on all 75 SZCH images (450 sites), never a locked set. The numbers below are copied from the decision log
of the project; they are development evidence and are not used for any confirmatory claim.

\paragraph{Round 1: training source.} With SwinV2-T at \num{100000} samples, switching the source from LIDC renders to the Qure.ai projections
lowered the real-image absorption from 0.300 to 0.216, at a cost of about 0.5\,dB of development bone-region PSNR$_{cal}$; adding lesion
insertion on top (0.218) or mixing the two sources (0.286, 0.269 with insertion) did not do better, and a twin term of weight 1 on a mixed
source moved absorption only slightly (0.260). The reference model absorbed 0.134 on the same sites.

\paragraph{Round 2: twin weight.} On the Qure source alone, absorption fell monotonically with the twin weight: 0.216, 0.186, 0.151, 0.103 and
0.020 at $w=0,1,4,16,64$. Development bone PSNR$_{cal}$ stayed flat up to $w=16$ (21.69--21.79\,dB) and fell to 21.27\,dB at $w=64$; the bone
error on the Qure-projection validation data rose with the weight (MAE 0.0044 to 0.0082 at $w=16$). Seed repeats put the seed spread at about
0.02 in absorption and 0.1\,dB in PSNR; three times the training length and the larger SwinV2-B encoder each moved absorption by about one
seed spread. The twin weight was the lever.

\paragraph{Rounds 3 and 4: seeds, budget and the operating point.} At $w=64$ two seeds disagreed strongly (absorption 0.020 and 0.061,
PSNR 21.27 and 20.80\,dB), so that weight was not usable; $w=16$ to 32 was stable. At \num{300000} samples four seeds of SwinV2-T with
$w=16$ gave absorption 0.059--0.071 (mean 0.067) at a development PSNR of 21.73\,dB, tighter than at \num{100000} samples, while $w=32$ cost
about 0.4\,dB for a further 0.009 of absorption. The operating point was fixed at $w=16$.

\paragraph{Wave 1: dose-response and ablation at the full budget.} The twin weights $w\in\{0,4,8,16,32,64\}$ and the ablation arms were
trained at \num{300000} samples with at least two seeds; their development results fed no decision about the final configuration and are
reported only in the ablation table (Table~\ref{tab:ablation_twin}); the dose-response itself is confirmatory and is evaluated on the locked
sets below.

\paragraph{Encoder screen and the final configuration.} ConvNeXt-V2-T, DINOv3 ViT-S/16 and RAD-DINO ViT-B/14 were screened at
\num{100000} samples (Table~\ref{tab:app_screen}; rule R1, Chapter~\ref{ch:benchmark}). ConvNeXt-V2-T met R1 after a rerun with a reduced
encoder learning rate, RAD-DINO missed it on PSNR by 0.05\,dB but was continued as exploratory, and both met the replacement rule R2 at
\num{300000} samples. The pre-declared tie-break (mean absorption within 0.005, then the higher bone PSNR) chose ConvNeXt-V2-T at two seeds
each; with the third seeds the absorption gap became 0.0055, just outside the band, and the choice was not reopened. The final model is the
ensemble of three ConvNeXt-V2-T suppressors with $w=16$; SwinV2-T and RAD-DINO are reported as secondary encoders.

\section{Dose-response and ablation}
\label{sec:res_twin}
On the locked VinDr-RibCXR audit the absorption of the SwinV2-T family falls from 0.183 at $w=0$ to 0.043 at $w=64$
(Table~\ref{tab:dose_response}, Figure~\ref{fig:dose_response}). The slope of $A$ on $\log_2(1+w)$ is $-0.0233$ $[-0.0249,-0.0217]$; the
upper bound is below zero, so H1$'$ passes. On VinDr-RibCXR the fall is steepest at small weights and continues, more slowly, up to $w=64$; on
the Qure-projection test split with exact truth the slope is $-0.0174$ $[-0.0190,-0.0160]$ and absorption hardly changes above $w=16$ (0.118, 0.122 and 0.112). On the
held-out SZCH pairs the bone-region PSNR$_{cal}$ of the family stays between 21.22 and 21.28\,dB up to $w=16$ and drops to 20.92 and 21.01\,dB
at $w=32$ and 64 (Figure~\ref{fig:pareto}). The confirmatory data thus reproduce what the development data had shown: the regulariser lowers
absorption in proportion to the logarithm of its weight, and its fidelity cost appears only above $w=16$.

The ablation (Table~\ref{tab:ablation_twin}, development data) shows which parts carry the effect. A stop-gradient on the branch without the
pasted lesion removes it entirely ($A$ 0.185 against 0.067 for the full regulariser and 0.190 without it): the network must be able to move both
predictions towards each other. A consistency term over the whole image also removes it (0.197), presumably because the per-pixel weight of
the lesion is diluted by the much larger background. Dilating the support by 31\,px keeps part of the effect (0.137); up to four lesions per
twin instead of up to two keep all of it (0.074). Inserting lesions without the twin term (CLIA) gives 0.191, the same as no twin at all,
which matches the v1 finding on real images: showing the model lesions is not the same as telling it to leave them alone.
\begin{table}[htbp]
\centering\small
\caption{Dose-response of the lesion absorption on the twin weight (SwinV2-T family): seeds per weight, mean $A$ on VinDr-RibCXR and on the
Qure-projection test split, and the pre-specified slope with its image-cluster bootstrap interval.}
\label{tab:dose_response}
\fittable{tables/dose_response.tex}
\end{table}
\begin{figure}[htbp]
\centering
\safefig{figures/dose_response.pdf}
\caption{Dose-response of the lesion absorption $A$ on VinDr-RibCXR on the twin weight $w$ (SwinV2-T), one marker per seed, with the
pre-specified slope on $\log_2(1+w)$.}
\label{fig:dose_response}
\end{figure}
\begin{table}[htbp]
\centering\small
\caption{Ablation of the twin regulariser at $w=16$ (SwinV2-T, 300k samples, mean over seeds): one factor changed per row. Development evidence:
audit on all 75 SZCH images, bone PSNR$_{cal}$ on the 24 development pairs.}
\label{tab:ablation_twin}
\fittable{tables/ablation_twin.tex}
\end{table}

\section{Encoder results on the locked sets}
The twin effect does not depend on the encoder (Table~\ref{tab:encoders}). On VinDr-RibCXR it lowers absorption by 0.127 for ConvNeXt-V2-T,
0.114 for SwinV2-T and 0.112 for RAD-DINO, with intervals well away from zero, and on JSRT by 0.073, 0.061 and 0.064. RAD-DINO reaches the same
absorption as the selected ConvNeXt-V2-T on VinDr-RibCXR (0.053 against 0.051); single-seed results of all three families are in
Table~\ref{tab:app_seeds}. These comparisons are exploratory: the plan's confirmatory contrasts concern the selected model only.
\begin{table}[htbp]
\centering\small
\caption{Twin regulariser per encoder on the locked audit sets (exploratory): absorption difference to the matched no-twin control and its
interaction with bone overlap.}
\label{tab:encoders}
\fittable{tables/encoders.tex}
\end{table}


\chapter{The counterfactual lesion audit}\label{ch:audit}
% Audit chapter: the paper section is input, then thesis-only implementation notes, validation and the confirmatory results (generated tables).
\section{Counterfactual lesion audit}
\label{sec:audit}
% PLAN (about 1.5 pages). The definition of A is already in methods.tex (Lesion-absorption audit); this section states the protocol and its limits.
%   5.1 Definition and reading of A: slope of the change in the bone estimate on the inserted signal; A=0 keeps the lesion, A=1 removes it; A=0 is
%       reached by doing nothing, so A is always reported next to a suppression gain; A can leave [0,1].
%   5.2 Sites and lesions: lesions from held-out CTs (library built from the TEST split, s4/lesion_lib.py), six sites per image at alpha=1, sites drawn
%       once with a fixed seed and shared by all models (s4/audit_real.py), lung masks from a published network, bone overlap from independent masks.
%   5.3 Why the audit is not circular: the regulariser pastes lesions to teach invariance and the audit pastes lesions to measure it. Mitigations:
%       held-out lesion library, different image domain, independent rib masks, and evidence that does not use insertion (H4' in SAP_v2.md).
%   5.4 Variants: exact-truth audit on the Qure-projection test split (secondary), SZCH audit (development evidence only), amplitude sweep.
%   5.5 Statistics (SAP_v2.md H1', H3'): image-cluster bootstrap, slope of A on log2(1+w), Holm over contrasts.
\subsection{Definition}
For a suppressor $f$ that outputs the bone image, an image $x$ and a lesion signal $n\ge0$ in composite-image units, the absorption at amplitude
$\alpha$ is
\begin{equation}
A=\frac{\big\langle f(x+\alpha n)-f(x),\;\alpha n\big\rangle}{\big\langle \alpha n,\;\alpha n\big\rangle},
\label{eq:absorption}
\end{equation}
with both inner products taken over the lesion support. $A=0$ means that the lesion stays in the soft-tissue image and $A=1$ that it was removed with
the bone; $A$ is a regression slope and can fall outside $[0,1]$. A model that does nothing has $A=0$, so $A$ is reported together with a suppression
gain and never alone. Two suppression measures are paired with $A$: the bone-region PSNR on the held-out SZCH pairs and, on VinDr-RibCXR, the
reduction of gradient energy on rib-boundary pixels (the rib benchmark).

\subsection{Sites and lesions}
Lesions come from the CTs of the LIDC-IDRI test split only, so no audited lesion was seen in training. On each real radiograph six sites are drawn at random
positions inside the lung fields at $\alpha=1$; the sites depend only on the image, its lung mask, the lesion library and a fixed seed, and are shared by
every model. Lung masks come from the PSPNet anatomical segmentation model of TorchXRayVision \citep{cohen2022torchxrayvision}, eroded by
49~px at $1024$~px so that a site lies inside the lung field; the model comes from the ChestX-Det project \citep{lian2021chestxdet}, whose paper does not fully
document its training data. The library holds 296 nodules
(4.0 to 37.6~mm, median 7.1~mm); six sites per image give 1,470 sites on VinDr-RibCXR, 1,482 on JSRT and 450 on SZCH. Because the PSPNet masks
differ slightly between GPUs, evaluations run on cloud GPUs read the lung masks exported on the machine of the final table, so every model is
scored on identical sites.
The bone overlap of a site is measured with independent information, never with a model output: the radiologist rib masks of VinDr-RibCXR
\citep{nguyen2021ribcxr} for the confirmatory audit, and the real dual-energy bone signal on SZCH. A site \emph{overlaps bone} when at least 15\,\% of its
support lies on an annotated rib. On JSRT no bone reference exists; the plan named the consensus of the bone maps of ours and of the reference
model, but the audit took the overlap from the reference model's bone map alone (secondary analysis; recorded as a deviation).

\subsection{Why the audit is not circular}
The twin regulariser pastes lesions to teach invariance, and the audit pastes lesions to measure it, so a model could pass the audit by learning the
insertion artefact. Four design choices limit this. The audit lesions come from CTs that were never used for training; the audit images are real
radiographs, whereas the training composites are synthetic; the confirmatory audit uses rib masks that no model has seen; and the claim that the
regulariser helps is also tested with evidence that does not use insertion: contrast of real nodules on JSRT and Node21 and detection recall of real
lesions grouped by bone overlap (Section~\ref{sec:experiments}). These choices limit the circularity but cannot remove it: a model that recognises
pasted signal in general could still score well, which is why the evidence without insertion is reported beside the audit (Section~\ref{sec:limitations}).

\subsection{Variants}
\emph{Exact-truth audit.} On the Qure-projection test split the lesion is added to the soft component, the true bone image is known, and the
bone error inside the lesion support can be measured against it in addition to $A$. \emph{SZCH audit.} It ran on all 75 SZCH images during development
and is reported as development evidence only (Section~\ref{sec:data}). \emph{Amplitude sweep.} $A$ is recomputed for
$\alpha\in\{0.25,0.5,1,2,4\}$ on the 450 SZCH sites (development data, supplementary material). For $\alpha\le1$ the absorption is flat
(v1 model 0.301--0.306, no-twin model 0.215--0.218, reference model 0.134--0.136), so the leak is linear in the lesion amplitude; at $\alpha=2$ and
4 it falls for every model, where the added signal starts to saturate at the top of the intensity range.

\subsection{Statistics}
Site-level differences between models are paired; confidence intervals are image-cluster bootstrap intervals; the dose-response over the twin weight is
the slope of $A$ on $\log_2(1+w)$ with an image-cluster bootstrap interval, and the contrasts of ours against the reference model and against the
no-twin control are adjusted with the Holm procedure (SAP\_v2.md). Bootstrap intervals use \num{10000} resamples for mean absorption and contrasts,
\num{2000} for the overlap interaction and 500 for the dose-response slope, always resampling whole images with replacement.


\section{Implementation notes}
The audit is implemented in \path{s4/audit_real.py}. For each radiograph it draws six sites with a random generator seeded once (seed 7):
for every site a lesion is drawn from the library and a centre is drawn among the pixels of the eroded lung mask; the lesion patch, rendered
at $512$~px from a LIDC-IDRI test CT, is scaled to the $1024$~px frame and added in image units. The sites depend only on the image, its lung
mask, the library and the seed, so every model, run separately, is scored on the same sites, and the paired differences between models are
site by site. A site whose support (pixels above 2\,\% of the lesion maximum) has fewer than 30 pixels is skipped. For each model the bone
estimate of the radiograph is computed once, and for each site the estimate of the radiograph with the lesion; $A$ is the regression slope of
Equation~\ref{eq:absorption} over the support, and a leak term records the change of the bone estimate far from the lesion. The bone overlap of
a site is the fraction of its support that lies on bone: the radiologist rib masks on VinDr-RibCXR, the strongest 20\,\% of the dual-energy
bone signal on SZCH-X-Rays, and a thresholded bone map of the reference model on JSRT and BDCXR-3257.

Two practical points matter for reproducibility. First, the lung masks come from a neural network whose output differs slightly between GPUs,
which moves the sites; evaluations that ran on cloud GPUs therefore read the lung masks exported on the machine that produced the final
tables, and the identity of the sites was checked file by file (Chapter~\ref{ch:ethics}). Second, the audit of a three-model ensemble runs
the three members on each modified image, so it is the most expensive evaluation of the study; on one T4 GPU a batch of four models, two of them
ensembles, took 16 minutes for the 245 VinDr-RibCXR images.

\section{Validation of the audit}
\paragraph{Controls.} The identity has $A=0$ by construction, at every site and on every data set, which the tables confirm. It also removes
no bone, which is why absorption is always shown next to a suppression measure.

\paragraph{Exact truth.} On the Qure-projection test split the lesion is added to the soft component and the true bone image is known
(Table~\ref{tab:app_qp}). The ranking of the models is the same as on real radiographs: ours absorbs 0.102, the matched control 0.220 and the
reference model 0.199, and the contrast of the lesion in our soft estimate is 0.99 of its true contrast against 0.94 for the control. The cost
of the regulariser shows here as well: the bone error of ours (0.0041) is higher than that of the control (0.0029), though lower than that of
the reference model on these images (0.0082).

\paragraph{Linearity.} On the development SZCH sites the amplitude of the inserted lesion was varied over $\alpha\in\{0.25,0.5,1,2,4\}$
(Table~\ref{tab:app_amplitude}). Up to $\alpha=1$ the absorption is flat for every model (v1 0.301--0.306, no twin 0.215--0.218, reference
model 0.134--0.136), so $A$ measures a property of the model, not of the lesion amplitude; above that it falls for all of them, where the
added signal begins to saturate at the top of the intensity range. The audit is run at $\alpha=1$.

\section{Confirmatory audit on VinDr-RibCXR}
On the 1,470 sites of VinDr-RibCXR (Tables~\ref{tab:audit_ribcxr} and~\ref{tab:audit_contrasts}) ours absorbs $A=0.051$ $[0.045,0.057]$,
the reference model 0.149 $[0.138,0.160]$ and the no-twin control 0.178 $[0.167,0.188]$. The paired differences are $-0.098$
$[-0.108,-0.089]$ and $-0.127$ $[-0.134,-0.119]$, Holm-adjusted $p<0.001$ for both: H3$'$ passes. The interaction between sites on an annotated
rib and the other sites is $0.003$ $[-0.017,0.027]$ against the reference model, so the advantage is the same with and without rib overlap. On
JSRT the same pattern holds ($A$ 0.048, 0.118 and 0.121) with a larger advantage at bone-overlapping sites ($-0.060$ $[-0.074,-0.046]$), but
there the overlap was defined by the reference model's own bone map. The SZCH audit, development evidence, agrees (0.054, 0.130 and 0.195).
Figure~\ref{fig:audit_maps} shows one site on a held-out SZCH image, chosen by a fixed rule.
\begin{table}[htbp]
\centering\small
\caption{Counterfactual lesion audit on VinDr-RibCXR: mean absorption per method with image-cluster bootstrap intervals (lower is better).}
\label{tab:audit_ribcxr}
\fittable{tables/audit_ribcxr.tex}
\end{table}
\begin{table}[htbp]
\centering\small
\caption{Paired site-level contrasts on the three audit sets, Holm over the two contrasts, and the interaction with bone overlap
($\geq$15\,\% vs.\ $<$15\,\%). SZCH is development evidence; the JSRT overlap comes from the reference model's bone map.}
\label{tab:audit_contrasts}
\fittable{tables/audit_contrasts.tex}
\end{table}
\begin{figure}[htbp]
\centering
\safefig{figures/audit_maps.pdf}
\caption{The audit on one held-out SZCH-X-Rays image (CC BY-NC 4.0): a 7\,mm nodule from a LIDC test CT is added; the change of each bone
estimate (red = lesion signal moved into the bone image) and the resulting soft-tissue image. Eight placements were drawn with a fixed seed and
the one with the median difference between ours and the reference model is shown.}
\label{fig:audit_maps}
\end{figure}

\section{Failure cases and limits of the audit}
The audit measures the leak of an added signal under a linear, additive model of image formation. Real lesions change attenuation inside the
patient, interact with scatter and with the tone curve of the detector, and may differ in texture from the CT patches used here; $A$ is
therefore a property of the suppressor for additive signals, not a measure of clinical detectability. Lesions near the lung edge can extend
beyond the eroded lung mask, very small lesions have supports near the 30-pixel limit where the regression slope is noisy, and saturated pixels
clip the added signal, which is why the amplitude is kept at $\alpha=1$. The bone overlap of a site is only as good as its reference: rib masks
on VinDr-RibCXR, the dual-energy signal on SZCH, and a model's bone map elsewhere. Finally, the audit and the twin regulariser both insert
lesions; the audit uses lesions from CTs never seen in training and real radiographs never used for training, and Chapter~\ref{ch:benchmark}
reports evidence that uses no insertion at all, but a model that learned to recognise pasted signal in general could still score well.


\chapter{Benchmark and downstream tasks}\label{ch:benchmark}
% Benchmark chapter: experimental design (paper section) and the locked-set results in thesis length. Result tables are the generated files in
% ../paper/tables; numbers in the text are copied from results/final/stats_held.json. The paper's results.tex is NOT input (labels would clash).
\section{Experiments}
\label{sec:experiments}
% PLAN (about 2 pages and one table). Design only: no result appears here. Everything below restates SAP.md / SAP_v2.md, which were frozen before the
% locked evaluations (tools/freeze.py, tools/freeze_v2.py); if this text and SAP_v2.md differ, SAP_v2.md wins and this text is corrected.
%   6.1 Models and comparators (ours, no-twin control, v1, Qure released weights, identity, CLAHE, homomorphic filtering).
%   6.2 Arms (Table arms) and the dose-response over the twin weight.
%   6.3 Protocols per hypothesis H1'-H5' (what is computed, on which set, with which test).
%   6.4 Metrics. 6.5 Downstream protocol. 6.6 Encoder screen (selection rule fixed before the screen).
\subsection{Models and comparators}
\emph{Ours} is the ensemble of three ConvNeXt-V2-T UNet++/CBAM suppressors (seeds 0--2) trained on the Qure projections with twin weight $w=16$
(Section~\ref{sec:method_v2}). It replaced the SwinV2-T configuration of the analysis plan under the pre-declared replacement rule (R2) before the
freeze: ConvNeXt-V2-T and RAD-DINO both met it on development data, and the pre-declared tie-break (mean absorption within 0.005, then the
higher bone-region PSNR) chose ConvNeXt-V2-T at two seeds each. With the third seeds the absorption gap is 0.0055, just outside the band; the
choice was not reopened, and both encoders are reported. The SwinV2-T and RAD-DINO pairs (twin and no twin, three seeds each) are secondary.
Comparators are: the \emph{no-twin control} (same encoder, data and recipe, $w=0$, three seeds); the v1 model trained on LIDC-IDRI only (\texttt{clia},
pre-specified pick of the v1 study); the publicly released weights of \citet{angaitkar2026anatomy} (SwinV2-B at $1024$~px; research licence);
and three controls that need no training: identity, CLAHE \citep{zuiderveld1994clahe} and a homomorphic filter
\citep{oppenheim1968homomorphic}. All methods are scored with the same code on the same images.

\subsection{Arms}
Table~\ref{tab:arms} lists the arms. Each arm changes one factor of its parent. The confirmatory dose-response varies the twin weight over
$w\in\{0,4,8,16,32,64\}$ in the SwinV2-T family with 3, 2, 2, 4, 2 and 2 seeds ($w=1$ appears only in the screening). The insertion-only arm and the
last four rows form the ablation of the regulariser (two seeds each, SwinV2-T, $w=16$).

\begin{table}[htbp]
\centering\small
\renewcommand{\arraystretch}{1.15}
\caption{Training arms. ``Lesions'' is the maximum number of lesions inserted into the base image before the physics (counterfactual insertion);
``Twin'' is the twin weight $w$, varied over $\{4,8,16,32,64\}$ in the \texttt{qp\_twin}$w$ arms ($w=1$ in the screening only). The parent of an arm is \texttt{recipe} for the
LIDC arms, \texttt{qp} for the first Qure arms and \texttt{qp\_twin16} for the last four rows.}
\label{tab:arms}
\begin{tabular}{@{}p{0.25\linewidth}p{0.08\linewidth}p{0.09\linewidth}p{0.08\linewidth}p{0.26\linewidth}@{}}
\toprule
Arm & Source & Lesions & Twin & Change \\
\midrule
\texttt{recipe} & LIDC & 0 & none & v1 baseline \\
\texttt{nodule\_loss} & LIDC & 0 & none & lesion-weighted loss \\
\texttt{clia} & LIDC & 3 & none & counterfactual insertion \\
\texttt{clia\_nodule} & LIDC & 3 & none & insertion and lesion-weighted loss \\
\texttt{qp} & Qure & 0 & none & Qure projections (no-twin control) \\
\texttt{qp\_clia} & Qure & 3 & none & insertion only, no tie term \\
\texttt{qp\_twin}$w$ & Qure & 3 & varied & twin regulariser \\
\texttt{qp\_twin16\_sg} & Qure & 3 & 16 & stop-gradient on the un-pasted branch \\
\texttt{qp\_twin16\_wide} & Qure & 3 & 16 & support dilated with a $31\times31$ window \\
\texttt{qp\_twin16\_all} & Qure & 3 & 16 & whole-image support \\
\texttt{qp\_twin16\_n4} & Qure & 3 & 16 & four lesions per twin \\
\bottomrule
\end{tabular}
\end{table}

\subsection{Protocols per hypothesis}
The analysis plan (SAP.md for the v1 family, SAP\_v2.md for v2) was frozen before any held-out or locked evaluation, on 1~November 2026
(\path{freeze_v2.json}: hashes of both plans, the SZCH development and held-out lists, the split files and the 52 weight files, at commit
\texttt{dd52b6a}). Paired differences are tested with the Wilcoxon signed-rank test \citep{wilcoxon1945individual} and given percentile bootstrap intervals
\citep{efron1993bootstrap}. Multiplicity is handled by the Holm procedure \citep{holm1979simple} on the $p$-values within each family; confidence intervals are unadjusted
95\,\% intervals; effect sizes are always reported with them, and a hypothesis that fails is reported as failed. The external validation on
BDCXR-3257 has its own addendum, frozen before any model saw those images (Section~\ref{sec:bdcxr_protocol}).
\begin{description}
\item[H1$'$, dose-response.] On VinDr-RibCXR (245 images, six sites each, lesions from LIDC test CTs, $\alpha=1$), the per-site absorption $A$ averaged over
 seeds falls with the twin weight: the slope of $A$ on $\log_2(1+w)$, with an image-cluster bootstrap interval, passes if its upper bound is below zero.
 The same analysis on the Qure-projection test split (exact truth) is secondary.
\item[H2$'$, fidelity.] On the 51 held-out SZCH pairs, bone-region PSNR of ours against the reference model is non-inferior within a margin of 0.5~dB
 (one-sided 95\,\% lower bound of the paired mean difference, image-level bootstrap); SSIM has a margin of 0.01 and is reported either way; ours against the
 no-twin control uses a margin of 0.25~dB (the cost of the regulariser). Metrics are computed after a per-image gain and offset fit to the DES reference
 (``calibrated'', primary) and without it (``raw'').
\item[H3$'$, lesion absorption.] On VinDr-RibCXR sites, $A$(ours)$<A$(reference model) and $A$(ours)$<A$(no-twin control), as paired site-level
 differences with image-cluster bootstrap intervals, Holm over the two contrasts, and an interaction between sites that overlap an annotated rib and sites
 that do not.
\item[H4$'$, evidence without insertion.] (a) Contrast-to-noise change of real nodules on JSRT and Node21, ours against the reference model and against
 the no-twin control, for all lesions and for lesions with at least 15\,\% bone overlap. (b) Downstream: NIH ChestX-ray14 macro AUROC per input
 (original, ours, reference model, CLAHE, original plus ours; three seeds); TBX11K, Node21 and VinBigData detection by FROC and mAP@50; recall at two false
 positives per image for lesions with and without bone overlap, as a difference in differences against the original input.
\item[H5$'$, reliability.] Spearman correlation between ensemble spread and absolute bone error on the SZCH held-out pairs (pseudo-truth bone from the
 least-squares-aligned DES image); split-conformal coverage calibrated on Qure-projection validation CTs and measured on the Qure-projection test split
 (nominal 90\,\%) and, as measured under shift, on SZCH.
\end{description}

\subsection{Metrics and downstream models}
Paired fidelity uses PSNR and SSIM \citep{wang2004ssim} on the bone region and on the whole image, band-limited SSIM, MS-SSIM \citep{wang2003msssim}, LPIPS \citep{zhang2018lpips}
and a high-frequency energy ratio; the same smoothing is applied to every method and never tuned. The rib benchmark reports the relative reduction of
Sobel gradient energy on the rib-boundary band inside the lungs and on lung pixels away from ribs, their difference (selectivity) and the ROC-AUC of the
gradient-reduction map. Detection is scored by FROC sensitivity at fixed false-positive rates, averaged as the competition performance metric
\citep{vanginneken2010anode09,niemeijer2011cpm}, and by mAP@50; AUCs are compared with the paired DeLong test \citep{delong1988auc}. The downstream classifier is a DenseNet-121
\citep{huang2017densenet} and the detector a Faster R-CNN with a ResNet-50 FPN backbone \citep{ren2017fasterrcnn}, with identical settings for every
input variant (\path{s4/cls.py}, \path{s4/det.py}). The classifier sees $256$~px images, is trained for 10 epochs with AdamW
\citep{loshchilov2019adamw} (learning rate $3\times10^{-4}$, weight decay $10^{-4}$, batch 64, cosine schedule) and keeps the epoch with the
best validation macro-AUROC. The detector sees $512$~px images, is trained for 12 epochs with SGD (learning rate 0.01, momentum 0.9, weight
decay $10^{-4}$, batch 8) and keeps the epoch with the best validation CPM. A multi-channel input (original plus soft, or original, soft and bone)
stacks the images as channels. Detection metrics treat tied scores as one threshold, so they do not depend on the order in which a sort
returns ties.

\subsection{Encoder screen}
To check whether a newer encoder helps under the same recipe, ConvNeXt V2 \citep{woo2023convnextv2}, DINOv3 \citep{simeoni2025dinov3} and RAD-DINO
\citep{perezgarcia2025raddino} were trained at a reduced budget under the same recipe, and an encoder moves to the full budget only if
it meets a rule fixed in advance (RUNPLAN.md, rule R1): at \num{100000} samples, absorption $A\le0.070$ on the development audit with
bone-region PSNR$_{cal}\ge21.60$~dB on the development pairs, or PSNR$_{cal}\ge21.95$~dB with $A\le0.100$. The thresholds come from the
SwinV2-T range at that budget and were written after one SwinV2-B run had been seen; the decision at the full budget used rule R2 (see Models and comparators above). Pretrained encoders other than SwinV2 are trained with a learning rate $0.25\times$ that of the decoder;
ConvNeXt-V2-T diverged at the full rate (fp16 overflow) and was rerun at $0.25\times$, an extension of the rule written for vision transformers.
DINOv3 ViT-S/16 did not improve on SwinV2-T, RAD-DINO ViT-B/14 missed R1 on PSNR by 0.05~dB but was continued at the full budget as an
exploratory encoder, and ConvNeXt-V2-T met R1.

\subsection{External validation protocol (BDCXR-3257)}
\label{sec:bdcxr_protocol}
An addendum to the analysis plan (SAP\_bdcxr.md) was frozen before any model of this study saw a BDCXR-3257 image; its one change, the
duplicate grouping (Section~\ref{sec:data}), was made from the images alone before any model ran and is recorded in the deviations log. All
suppressors are the frozen models; only the downstream classifiers are trained. (E1, primary) Pneumonia classification with three one-vs-rest
outputs, inputs as in H4$'$b (original, ours, reference model, CLAHE, original plus ours), the classifier settings of the downstream protocol,
five folds stratified by class over the near-duplicate groups (fold seed 0) with 10\,\% of each training part, by whole groups, for epoch selection,
and three seeds per fold and input (75 runs). Out-of-fold probabilities, averaged over seeds, give one prediction per image. The primary measure
is the pooled macro-AUROC; ours against the original image is non-inferior if the lower bound of the 95\,\% interval of the difference (group
bootstrap) exceeds $-0.01$, and superior if it exceeds 0; per class, paired DeLong tests with Holm over the three classes. (E2, secondary) The
lesion audit of Section~\ref{sec:audit} on a fixed class-stratified subset of 300 images (seed 7), six sites per image, with bone overlap from the
reference model's bone map; images whose lung mask covers less than 2\,\% of the image are skipped and counted (none was). (E3, descriptive) The
per-image ensemble spread over the lung field, compared with the held-out SZCH pairs and correlated with age.


\section{Fidelity on held-out paired radiographs}
On the 51 held-out SZCH pairs (Tables~\ref{tab:szch_heldout} and~\ref{tab:h2}) the bone-region PSNR$_{cal}$ of ours is 21.46\,dB and that of
the reference model 21.51\,dB; the paired difference is $-0.06$\,dB with a one-sided 95\,\% lower bound of $-0.15$\,dB, well inside the
non-inferiority margin of $-0.5$\,dB. Against the no-twin control the difference is $+0.01$\,dB (lower bound $-0.02$\,dB against a margin of
$-0.25$\,dB). Both parts of H2$'$ that concern bone-region fidelity pass. SSIM$_{cal}$ fails: ours is lower than the reference model by 0.021,
with a lower bound of $-0.023$ against a margin of $-0.01$, as the plan had anticipated from the development data.

The SSIM difference was examined on development data before the freeze (Appendix~\ref{app:negative}). The dual-energy soft-tissue image has
less high-frequency energy than a radiograph, so full-resolution SSIM against it rewards smooth outputs. On the held-out pairs, ours keeps
about 9\,\% more high-frequency energy than the dual-energy image in non-bone regions, the reference model about 1\,\% less; with a Gaussian
band limit of $\sigma=4$ the SSIM of the two models is 0.946 and 0.948, so almost all of the gap is fine texture. Whether a $1024$~px output,
as the reference model has, would close the rest was not tested. LPIPS points the same way (0.234 for ours, 0.166 for the reference model).
BS-LDM, trained on dual-energy pairs from the same hospital as the test images, reaches 28.57\,dB bone-region PSNR$_{cal}$; it is shown
for orientation only.
\begin{table}[htbp]
\centering\small
\caption{Held-out SZCH-X-Rays pairs: bone-region and whole-image metrics for ours and the comparators, calibrated (primary) and raw.}
\label{tab:szch_heldout}
\fittable{tables/szch_heldout.tex}
\end{table}
\begin{table}[htbp]
\centering\small
\caption{H2$'$ non-inferiority on the 51 held-out SZCH-X-Rays pairs (paired differences, image bootstrap, one-sided 95\,\% lower bound).}
\label{tab:h2}
\fittable{tables/h2_noninferiority.tex}
\end{table}
\begin{figure}[htbp]
\centering
\safefig{figures/pareto.pdf}
\caption{Lesion absorption on VinDr-RibCXR against bone-region PSNR$_{cal}$ on the 51 held-out SZCH-X-Rays pairs: the SwinV2-T family at
twin weights $0\ldots64$, ours and the no-twin control (ConvNeXt-V2-T, $K=3$), the reference model, the v1 model and the identity.}
\label{fig:pareto}
\end{figure}
\begin{figure}[htbp]
\centering
\safefig{figures/qualitative_szch.png}
\caption{Held-out SZCH-X-Rays cases (CC BY-NC 4.0, \citealp{sun2026bsldm}), chosen by rule: the held-out images at the 25th, 50th and 75th
percentile of the bone-region PSNR$_{cal}$ of ours. Radiograph, dual-energy soft-tissue reference, the soft-tissue estimates of ours and of the
reference model, and the bone estimate of ours.}
\label{fig:qualitative}
\end{figure}

\section{Rib benchmark}
On VinDr-RibCXR the rib benchmark measures how much gradient energy a suppressor removes on rib-boundary pixels and on other lung pixels
(Table~\ref{tab:rib_benchmark}). The reference model removes the most on both (48\,\% and 22\,\%), ours less (30\,\% and 8\,\%), and the
no-twin control lies in between (37\,\% and 14\,\%). Selectivity, the difference of the two, is 0.264 for the reference model and 0.220 for
ours, and the AUC that separates rib pixels from other lung pixels by the amount removed is the same (0.671 and 0.672). Ours is the more
conservative suppressor: it removes less rib edge and much less of everything else. Read with the audit, ours keeps lesions better per unit of
rib suppression (absorption a third of the reference model's at about 60\,\% of its rib-edge removal), but a share of the reference model's
higher absorption may simply reflect stronger suppression. The benchmark says nothing about lesions directly: it is a measure of what is
removed, not of what is kept.
\begin{table}[htbp]
\centering\small
\caption{Rib benchmark on VinDr-RibCXR (metrics only): gradient-energy reduction on rib-boundary pixels and on other lung pixels, selectivity
and AUC.}
\label{tab:rib_benchmark}
\fittable{tables/rib_benchmark.tex}
\end{table}

\section{Real nodules}
Real nodules answer RQ3 only in part (Table~\ref{tab:nodule_cnr}). On JSRT (154 nodules) the contrast-to-noise ratio of the nodule box rises
by 0.041 in the soft image of ours and by 0.018 in that of the reference model; the difference, $+0.023$ $[0.008,0.041]$, grows to $+0.032$
$[0.014,0.052]$ for the 123 nodules with at least 15\,\% bone overlap, where the reference model lowers contrast. On Node21 (1,474 boxes) the
reference model raises contrast more ($+0.103$ against $+0.078$; difference $-0.025$ $[-0.033,-0.018]$), also for the 1,109 bone-overlapping
boxes ($-0.017$ $[-0.026,-0.009]$). Against the no-twin control ours raises contrast in all four comparisons ($+0.007$ to $+0.027$, Holm
$p<0.001$). The twin regulariser thus helps real nodules relative to the same model without it, consistently; whether ours beats the reference
model depends on the data set, and on Node21 the reference model's stronger suppression raises box contrast more.
\begin{table}[htbp]
\centering\small
\caption{Contrast-to-noise change of real nodules on JSRT and Node21 (soft-tissue image minus original), all lesions and lesions with at least
15\,\% bone overlap; image-cluster bootstrap intervals.}
\label{tab:nodule_cnr}
\fittable{tables/nodule_cnr.tex}
\end{table}

\section{Downstream classification}
On NIH ChestX-ray14 (Table~\ref{tab:downstream_cls}) the three-seed ensemble of DenseNet-121 classifiers reaches a macro-AUROC of 0.818 on the
original images and 0.817 on the soft-tissue images of ours ($-0.001$ $[-0.004,0.003]$, patient-cluster bootstrap). Per class the paired
DeLong tests find four differences after Holm correction: ours is lower for effusion, infiltration and pleural thickening ($-0.004$, $-0.007$,
$-0.013$) and higher for nodule ($+0.011$). Adding the soft image as a second channel helps a little ($+0.006$ $[0.003,0.008]$), and so does the
reference model's soft image on its own ($+0.004$ $[0.002,0.007]$); CLAHE lowers the macro-AUROC ($-0.004$ $[-0.006,-0.002]$). The soft image
of ours is no substitute for the original, but the nodule class, the one the method was designed for, moves in its favour.
\begin{table}[htbp]
\centering\small
\caption{Downstream classification on NIH ChestX-ray14: macro-AUROC of the three-seed ensemble, per-seed mean $\pm$ SD and the difference to
the original image (patient-cluster bootstrap).}
\label{tab:downstream_cls}
\fittable{tables/downstream_cls.tex}
\end{table}

\section{Downstream detection}
For detection (Table~\ref{tab:downstream_det}, Figure~\ref{fig:downstream_forest}) the change in CPM with the soft image of ours alone is
$+0.001$ $[-0.023,0.023]$ on Node21, $+0.006$ $[-0.009,0.020]$ on TBX11K and $-0.009$ $[-0.016,-0.002]$ on VinBigData. The three-channel input
(original, soft and bone) gains $+0.028$ $[0.009,0.047]$ on Node21 and $+0.008$ $[0.002,0.014]$ on VinBigData; the reference model's soft image
gains $+0.029$ $[0.016,0.043]$ on TBX11K. The bone-overlap subgroup, recall at two false positives per image for lesions with and without at
least 15\,\% bone overlap compared as a difference in differences against the original input, shows no effect whose interval excludes zero
(Node21 $-0.006$ $[-0.071,0.061]$, VinBigData $+0.009$ $[-0.034,0.053]$); on TBX11K it cannot be estimated, because every test box overlaps bone
by at least 20\,\%. H4$'$b therefore finds no benefit of the soft image used alone; the gains come from giving the detector the soft and bone
images in addition to the original.
\begin{table}[htbp]
\centering\small
\caption{Downstream detection on Node21, TBX11K and VinBigData: CPM and mAP@50 (mean $\pm$ SD over three seeds), differences to the original
image and the difference in differences (DiD) of recall at two false positives per image (image bootstrap).}
\label{tab:downstream_det}
\fittable{tables/downstream_det.tex}
\end{table}
\begin{figure}[htbp]
\centering
\safefig{figures/downstream_forest.pdf}
\caption{Downstream effects of each input against the original image with 95\,\% intervals: macro-AUROC on NIH ChestX-ray14, CPM on the
detection sets.}
\label{fig:downstream_forest}
\end{figure}

\section{Comparison with the reference model and the controls}
The reference model is compared as released: a SwinV2-B suppressor that predicts at $1024$~px, trained by its authors on far more samples than
ours, run with its own pre- and post-processing. Our models predict at $512$~px with a tiny encoder. Differences between the two therefore mix
the effect of the regulariser with differences in resolution, capacity, data and training length; the matched no-twin control, which shares
everything with ours except the regulariser, is the comparison that isolates the regulariser, and it absorbs slightly more than the reference
model. The controls frame the scale of the effects: the identity removes nothing and absorbs nothing, CLAHE changes contrast everywhere and
does worse than the original image downstream, and homomorphic filtering removes low-frequency content without separating bone from soft
tissue. BS-LDM, trained on real dual-energy pairs from the test hospital, is far ahead in fidelity on that hospital's images and far behind
in the development audit; it answers a different question and is not part of a pre-specified comparison.


\chapter{Reliability}\label{ch:reliability}
% Reliability chapter (H5'). Numbers from results/final/unc_heldout.json and stats_held.json (H5 with image-bootstrap intervals); the
% exploratory recalibration (SAP_v3 H9) from paper/tables/v3_bone_intervals.tex (tools/hybrid_stats.py --tables_from).
\section{Design}
The deep ensemble of Chapter~\ref{ch:v2} gives, for every pixel, the mean $\bar f$ of the three bone estimates and their standard deviation
$\sigma$, the spread. Two questions follow (H5$'$). Does the spread rank the actual error, so that pixels where the members disagree are the
pixels where the bone estimate is wrong? And can the spread be turned into intervals that cover the true bone value at a stated rate? For the
second, a split-conformal scale $q$ is calibrated on the Qure-projection validation CTs, where the bone image is exact, so that
$[\bar f-q(\sigma+\varepsilon),\,\bar f+q(\sigma+\varepsilon)]$ has nominal coverage 0.90. Coverage is measured on the Qure-projection test
split (exact truth) and, as a measurement under domain shift, on the 51 held-out SZCH pairs against a pseudo-truth bone image: the radiograph
minus the dual-energy soft-tissue image after a least-squares gain and offset. The method is described in Section~\ref{sec:method_v2}.

\section{Results}
On the Qure-projection test split the intervals cover 0.900 $[0.897,0.903]$ of the pixels at the nominal 0.90, with a mean width of 0.018, and
the spread ranks the absolute error well (Spearman 0.76; Table~\ref{tab:reliability}). The calibration holds where it was made. On the held-out
SZCH pairs it does not: the same intervals cover 0.139 $[0.134,0.145]$ of the pixels, and the spread ranks the error only weakly (mean
per-image Spearman 0.11 $[0.07,0.15]$; Figure~\ref{fig:reliability}). Every one of the 51 images has coverage far below the nominal level, so
the failure is systematic rather than driven by a few images. H5$'$ holds on synthetic data and fails under the shift to real radiographs.
\begin{table}[htbp]
\centering\small
\caption{Ensemble spread and split-conformal intervals ($K=3$, nominal coverage 0.90), with image-bootstrap intervals.}
\label{tab:reliability}
\fittable{tables/reliability.tex}
\end{table}
\begin{figure}[htbp]
\centering
\safefig{figures/reliability.pdf}
\caption{Left: per-image coverage of the conformal intervals on the Qure-projection test split and on the held-out SZCH-X-Rays pairs (bars:
means; dashed: nominal 0.90). Right: per-image Spearman correlation of ensemble spread and absolute error on SZCH-X-Rays (line: mean).}
\label{fig:reliability}
\end{figure}

\section{Recalibration on real pairs (exploratory)}
After this failure was known, the third analysis plan (H9, Chapter~\ref{ch:classifier}) added a post hoc check on the same, already-seen
pairs: the intervals were recalibrated on the 24 SZCH development pairs and tested on the 51 held-out pairs. One variant uses a single global
calibration; a Mondrian variant calibrates separately per stratum of predicted bone overlap and spread tercile. Pooled coverage was 0.886
with the global calibration and 0.895 with the Mondrian calibration, whose intervals were less than half as wide (mean width 0.215 against
0.447; Table~\ref{tab:bone_recal}); the median image reached 0.886 and 0.894. Figure~\ref{fig:bone_recal} also shows a per-image variant that
takes the calibration quantile over per-image quantiles. A few real pairs therefore restore nominal coverage on average. The check is
exploratory: it was designed after the failure was seen, it uses pairs from one hospital, and the coverage of single images still varies.
\begin{table}[htbp]
\centering\small
\caption{Exploratory recalibration on real pairs (post hoc): bone-image intervals calibrated on the 24 SZCH-X-Rays development pairs and tested
on the 51 held-out pairs at nominal 0.90. Mondrian: calibration per stratum of predicted bone overlap and ensemble spread.}
\label{tab:bone_recal}
\fittable{tables/v3_bone_intervals.tex}
\end{table}
\begin{figure}[htbp]
\centering
\safefig{figures/fig_bone_uncertainty.pdf}
\caption{Exploratory recalibration on real pairs (post hoc). (a) Per-image coverage on the 51 held-out SZCH-X-Rays pairs for global and
Mondrian calibration, with quantiles pooled over pixels or taken per image; dashed: nominal 0.90. (b) Reliability of the Mondrian intervals
across nominal levels; cross: the synthetic calibration. (c) Sparsification of the absolute bone error by the ensemble spread, with the
error-sorted oracle.}
\label{fig:bone_recal}
\end{figure}

\section{Interpretation and limits}
Three things explain why the intervals fail on real images, and they cannot be separated with the data at hand. First, the spread rises only
modestly on real radiographs (mean interval width 0.026 against 0.018) while the error rises far more: the members agree with each other more
than they agree with the dual-energy reference, because they share their training data and their biases. Second, the pseudo-truth bone image is not the true bone image;
it carries the differences in tone and texture between a radiograph and a dual-energy soft-tissue image, so part of the measured error is
reference error. Third, split-conformal calibration guarantees coverage only for data exchangeable with the calibration data, and only marginally
over pixels, not per image; rendered projections and real radiographs are not exchangeable. The spread is still useful as a flag: on the
paediatric radiographs of Chapter~\ref{ch:external} it is 2.7 times its value on adult pairs. Calibrated intervals on real radiographs need
calibration data with real truth; the exploratory recalibration above shows that 24 dual-energy pairs from one site were enough there, not
whether they would transfer to other sites.


\chapter{External validation on paediatric radiographs}\label{ch:external}
% External validation on BDCXR-3257 (SAP_bdcxr.md). Numbers from results/bdcxr/stats.json (tools/bdcxr_stats.py). No BDCXR image may appear.
\label{sec:bdcxr}
\section{Purpose and data}
All models of this study were trained on adult data, and every locked set but one is adult. To see how the suppressor behaves on a population
it was never trained on, it was applied, frozen, to BDCXR-3257: 3,257 chest radiographs of children from the Radiology and Imaging Department of
Dr.~M~R~Khan Shishu Hospital \& Institute of Child Health, Dhaka, labelled normal (880), bacterial pneumonia (1,477) or viral pneumonia (900)
\citep{durjoy2026bdcxr}. Most children are under five years old (1,132 images are of children under one year). The data were used under the
ethical approval REC-FSIT-2025/No:~11928 and a data use agreement between the hospital and the author's research group; no image of this set is
shown in this thesis (Chapter~\ref{ch:ethics}). The set has no patient identifiers, no lesion boxes and no dual-energy pairs, so fidelity
cannot be measured on it; it supports a classification task, the lesion audit and the uncertainty signal.

\section{Protocol}
An addendum to the analysis plan, SAP\_bdcxr.md, was frozen on 3 November 2026 before any model of this study saw a BDCXR-3257 image. It
fixes three questions. E1 (primary): does the soft-tissue image of ours change pneumonia classification? Five folds stratified by class, three
seeds per fold and input, the classifier and settings of the downstream protocol of Chapter~\ref{ch:benchmark}, and the pooled out-of-fold
macro-AUROC; ours is non-inferior to the original image if the lower bound of the 95\,\% interval of the difference exceeds $-0.01$. E2
(secondary): the lesion audit on a fixed class-stratified subset of 300 images. E3 (descriptive): the ensemble spread over the lung field
compared with the adult SZCH pairs.

One change was made, before any model ran: the plan grouped near-duplicate images by a 64-bit image hash so that a group never spans two folds,
but on these images the hash chained 1,891 of the 3,257 images into one group. Near-duplicates were instead grouped by the correlation of the
$64$~px images ($r\ge0.93$, the lowest threshold of those tried at which no group grew beyond two images). This gave 48 pairs, eight of them
with different labels, and five folds of 650 to 653 test images each. Because the set has no patient identifiers, repeated images of one
child taken on different days cannot be detected; the grouping only removes near-identical images.

\section{Results}
\paragraph{E1: pneumonia classification.} The pooled out-of-fold macro-AUROC is 0.814 with the original image and 0.816 with the soft-tissue
image of ours (Table~\ref{tab:bdcxr_cls}). The difference, $+0.002$ $[-0.005,0.010]$, meets the non-inferiority margin and shows no superiority;
none of the per-class DeLong tests is significant after Holm correction. The reference model's soft image gives $+0.001$ $[-0.008,0.010]$,
original plus ours $-0.001$ $[-0.008,0.005]$, and CLAHE $-0.006$ $[-0.012,0.001]$, which misses the margin. Pneumonia against normal is
separated better than bacterial against viral pneumonia (AUROC about 0.88 against 0.77 for every input), and the per-fold spread of the
macro-AUROC (SD about 0.01) is as large as any difference between inputs.

\paragraph{E2: lesion audit.} On 1,800 sites in 300 paediatric images, with no image lost to a failed lung mask, ours absorbs $A=0.084$, the
reference model 0.217 and the no-twin control 0.251 (Table~\ref{tab:bdcxr_audit}); the paired differences are $-0.134$ $[-0.144,-0.124]$ and
$-0.167$ $[-0.175,-0.160]$, and the advantage is larger at sites that overlap bone. All three models absorb more than on adult images, and the
order is the same. The inserted lesions are adult nodules from LIDC-IDRI CTs, placed in children's lungs, so this is a stress test of the
suppressor rather than a statement about paediatric lesions.

\paragraph{E3: uncertainty.} The median lung-field spread of the ensemble is 0.0076 (interquartile range 0.0067 to 0.0088) on BDCXR-3257 and
0.0028 (0.0026 to 0.0033) on the adult SZCH pairs, 2.7 times higher, and it is higher in younger children (Spearman with age $-0.22$). The
ensemble knows, in this coarse sense, that these images are unlike its training data.
\begin{table}[htbp]
\centering\small
\caption{BDCXR-3257: pneumonia classification (normal, bacterial, viral), five folds over near-duplicate groups $\times$ three seeds; pooled
out-of-fold macro-AUROC, per-fold mean $\pm$ SD and the difference to the original image (group bootstrap).}
\label{tab:bdcxr_cls}
\fittable{tables/bdcxr_cls.tex}
\end{table}
\begin{table}[htbp]
\centering\small
\caption{BDCXR-3257: lesion audit on a fixed class-stratified subset of 300 images (1,800 sites).}
\label{tab:bdcxr_audit}
\fittable{tables/bdcxr_audit.tex}
\end{table}

\section{Reading}
On a population and a hospital the models never saw, suppression costs nothing in pneumonia classification and gains nothing either, which
matches the adult downstream results: pneumonia is a diffuse finding that bone shadows hide less than they hide a small nodule. The ordering
of absorption between ours, the reference model and the control carries over to children, and the uncertainty signal rises where it should.
The set is a single centre with image-level labels and no patient identifiers, so these results describe this set; they do not establish
paediatric use.


\chapter{Bone-suppressed inputs for chest radiograph classifiers}\label{ch:classifier}
% Classifier study: the third analysis plan (SAP_v3.md, hash-frozen in freeze_v3.json on 3 November 2026). Methods, results and discussion
% are the sections of the second paper (paper/main_cls.tex: cls_methods, cls_results, cls_discussion), input here once; this file adds only
% the question. \companion{...} in those sections becomes a chapter reference (thesis/main.tex). Every number comes from paper/tables/v3_*.tex.
\section{Question}
\label{sec:cls_intro}
Chapter~\ref{ch:benchmark} gave the soft-tissue image to downstream classifiers in two simple ways: used alone it did not improve them, and
added beside the radiograph it gave small gains; on paediatric radiographs it was non-inferior (Chapter~\ref{ch:external}). Two questions
remain. First, a suppressor computes its soft-tissue and bone images from the radiograph, so they hold no information a classifier could not
in principle extract itself. A gain can come only from the prior the suppressor learned from CT or from the way a network uses the extra
images, and a network that reads three images instead of one also spends three times the computation; a fair test therefore needs a control
that runs the same network on the radiograph alone. Second, suppression could help some findings and harm others, revealing small lesions
behind ribs while removing cues for diffuse or pleural findings, and a macro-average over classes would hide such a trade-off.

This chapter reports a third pre-specified study that addresses both. Five input configurations, among them a multi-stream network that
reads the radiograph, soft-tissue and bone images as separate streams and a control that feeds the same network the radiograph three times,
were trained with five seeds on NIH ChestX-ray14, VinBigData and BDCXR-3257, 175 runs in all. The plan continues the numbering of the earlier
plans: H6 is primary, H7, H8 and H10 are secondary, and H9, which concerns the suppressor's own intervals, is reported in
Chapter~\ref{ch:reliability}. The classifiers see the images as $X$, $S$ and $\hat B=X-S$, the radiograph and the suppressor's soft-tissue and
bone images. The study is also the second paper of this thesis; its sections are reproduced here.

\input{sections/cls_methods}
\input{sections/cls_results}
\input{sections/cls_discussion}


\chapter{An explainable, hallucination-aware Swin framework for chest CT and X-ray}\label{ch:crossmodal}
% Cross-modality study (v4): the fourth analysis plan SAP_v4.md and its design DESIGN_v4.md, hash-frozen in freeze_v4.json on 8 November 2026
% before the first v4 training run; later changes are in DEVIATIONS.md (Appendix app:deviations). Tables: paper/tables/v4_*.tex, written by
% tools/xm_tables.py from results/final_xm_stats.json (tools/xm_stats.py), results/v4/delong.json, results/v4/pointing_summary.json and
% results/v4/uq_cls_*.json; figures: paper/figures/fig_v4_*.pdf (tools/fig_v4.py). Numbers in the text are copied from those files.
\section{Question}
\label{sec:xm_intro}
The earlier chapters used CT in one way only: to render training images for a bone suppressor. No classifier ever saw a CT volume. This
chapter asks whether CT can teach a chest radiograph classifier directly. CT resolves in three dimensions what a radiograph superimposes,
so a classifier trained on CT learns its findings without the overlap of ribs, heart and mediastinum. If that knowledge passes to an X-ray
model through images of the same patients, the X-ray model might classify real radiographs better, look at the right place and report
fewer findings that are not there.

The chapter builds and tests such a framework: an explainable and hallucination-aware Swin transformer for cross-modality chest CT and
X-ray classification. One Swin family, SwinV2-T \citep{liu2022swinv2}, serves both modalities. A CT teacher learns six findings from
volumes of CT-RATE \citep{hamamci2026ctrate}. An X-ray model learns the same findings from frontal projections of the same volumes, once on
its own (P0) and once with knowledge distillation from the teacher (P1) \citep{hinton2015distilling}. Both then initialise classifiers of
the 14 NIH ChestX-ray14 diseases \citep{wang2017chestxray}, beside a classifier started from ImageNet weights. Five hypotheses were fixed
before any v4 model was trained, in a fourth analysis plan that is pre-specified and hash-frozen (8 November 2026) but not deposited in a
public registry; they continue the numbering of the earlier plans. H11 (primary) and H12 ask whether distillation from CT improves the
classification of real radiographs; H13 compares the Swin classifier with the DenseNet-121 of Chapter~\ref{ch:classifier}; H14 asks
whether the CT and the X-ray models attend to the same place in the same patient; and H15 asks whether the CT-distilled classifier
hallucinates less when an image gains anatomy but no disease.

\section{Materials and methods}
\label{sec:xm_methods}

\subsection{Data}
\label{sec:xm_data}
\paragraph{Paired CT and projections.} The Qure.ai ct2xr projections \citep{angaitkar2026anatomy} are frontal projections of CT-RATE
volumes, one per volume, each with six findings taken from its scan: lung opacity, pleural effusion, atelectasis, cardiomegaly,
consolidation and lung lesion. They were split by patient as in Chapter~\ref{ch:v2} (\num{21051} training, 406 validation and 430 test
projections). The CT subset holds one reconstruction per patient: every validation and test patient of that split and \num{1600} training
patients drawn with a fixed seed, \num{2156} volumes in all (\num{1600} training, 259 validation, 297 test). Volumes that CT-RATE lists as
not showing the chest and volumes not scanned head first and supine were never selected, and their projections were left out of the
projection models as well, which leaves \num{20533} training, 396 validation and 415 test projections. Each volume was reduced inside the
cloud kernel to 24 axial slices spread over the lung extent, with three Hounsfield windows as channels (lung, mediastinum, bone) and
cropped to the body at 256\,px; the lung box and an anteroposterior attenuation projection were kept for the heat maps. CT-RATE is
released under CC BY-NC-SA~4.0; following the design, nothing derived from it left the cloud accounts except metrics and the
$32\times32$ heat-map frames, with the logged exceptions of Appendix~\ref{app:deviations} (teacher outputs, predictions and the weights of
the distilled models, which stayed private on the author's computer).

\paragraph{Real radiographs.} NIH ChestX-ray14 and VinBigData were used with the splits of Chapter~\ref{ch:classifier}
(Section~\ref{sec:cls_data}): NIH \num{77890} training, \num{8634} validation and \num{25596} test images with the 984 expert boxes of the
test list; VinBigData \num{2943} test images, used for zero-shot transfer only. A fixed label map links the six projection findings to
the real data sets: pleural effusion, atelectasis, cardiomegaly and consolidation map to the classes of the same name, and lung lesion to
NIH Mass or Nodule (VinBigData Nodule/Mass). These five are the exactly mapped shared findings of H11 and H12. Lung opacity maps exactly to
VinBigData Lung Opacity but only approximately to NIH Infiltration; it is reported but not tested.

\subsection{Models and training}
\label{sec:xm_models}
All models use SwinV2-T (\texttt{swinv2\_tiny\_window8\_256}, ImageNet-1k weights from timm \citep{wightman2019timm}) and the class-query
decoder of Chapter~\ref{ch:classifier} (Section~\ref{sec:cls_methods}); learning rate $3\times10^{-4}$ with the encoder at 0.25 of it, and
five seeds per model (Fig.~\ref{fig:xm_design}).

\begin{figure}[htbp]
\centering
\safefig{figures/fig_v4_design.pdf}
\caption{Study design. A 2.5D CT teacher learns six findings on the CT subset; projection models learn them on the Qure.ai projections
without (P0) and with (P1) distillation from the teacher; NIH classifiers start from ImageNet (F$_\mathrm{in}$), P0 (F$_\mathrm{p0}$) or P1
(F$_\mathrm{p1}$). All models are SwinV2-T with the class-query decoder, five seeds each; the DenseNet-121 of
Chapter~\ref{ch:classifier} is the CNN reference.}
\label{fig:xm_design}
\end{figure}

\paragraph{CT teacher.} A 2.5D network: the X-ray encoder runs on each of the 24 slices with the three windows as colour channels; every
slice yields $8\times8$ tokens with a two-dimensional sine code and a learned slice-depth embedding, and the class-query decoder attends
over the tokens of all slices. A full 3D Swin did not fit a T4 GPU at a useful batch size, and the 2.5D design keeps the ImageNet weights
and the architecture of the X-ray model. It was trained for 20 epochs with four volumes per step, keeping the best validation epoch.

\paragraph{Projection models.} P0 learns the six findings on the projections for 15 epochs (batch 32). P1 is trained the same way but,
on the \num{1600} training rows whose volume is in the CT subset, also matches the teacher's logits (binary distillation at temperature
2, weight 1) and aligns its pooled decoder output with the teacher's through a linear head (cosine loss, weight 0.5); half of every batch
is drawn from these paired rows. P1 seed $k$ distils from the teacher of slot $k$ (below).

\paragraph{NIH classifiers.} F$_\mathrm{in}$ starts from ImageNet, F$_\mathrm{p0}$ from P0 and F$_\mathrm{p1}$ from P1 (encoder and
decoder loaded, class heads new); each is trained on the 14 NIH diseases for 10 epochs (batch 32) with the radiograph as its only input,
keeping the best validation epoch. The five DenseNet-121 radiograph-only runs of Chapter~\ref{ch:classifier} are the CNN reference; they
were trained for 20 epochs at the full encoder learning rate, the pre-specified rule for batch-normalised encoders, so H13 compares
encoders trained with different budgets, neither tuned further.

\paragraph{Seed handling.} CT teacher seed 2 did not learn: its validation macro-AUROC stayed between 0.46 and 0.54 and its logits were
constant over all volumes. Before any P1 model of seed 2 was trained, a rule was fixed and logged: train seeds 5 and 6 with the unchanged
recipe and give slot 2 to seed 5 unless its logits were constant (per-class standard deviation below 0.01 over the training and validation
volumes in every class), otherwise to seed 6. Seed 5 learned and took slot 2; seed 6 was not used. The teachers of slots 0--4 are
therefore seeds 0, 1, 5, 3 and 4.

\subsection{Explainability}
\label{sec:xm_xai}
\paragraph{X-ray.} Grad-CAM++ on the last encoder stage \citep{chattopadhay2018gradcampp}, Integrated Gradients
\citep{sundararajan2017axiomatic} and the attention of the last decoder layer were computed as in Section~\ref{sec:cls_xai}, and the
pointing game was played on the 984 NIH boxes with the same tolerance of 15\,px at 256\,px \citep{zhang2018topdown}, for each seed of
F$_\mathrm{in}$, F$_\mathrm{p0}$, F$_\mathrm{p1}$ and the DenseNet-121 reference.

\paragraph{CT and cross-modal agreement (H14).} For each CT subset test scan positive for a finding, Grad-CAM++ was computed on every slice
of the teacher, the slice maps were stacked into a saliency volume and summed along the anteroposterior axis, which gives a frontal map in
the geometry of the paired projection. The CT map and the X-ray map of the paired projection (P1, and P0 for comparison) were cropped to
their own lung boxes, resized to $32\times32$, scaled to sum 1 and averaged over the five seed pairs (teacher of slot $k$ with X-ray seed
$k$). Agreement is the Pearson correlation of the matched maps; its null pairs the CT map with the X-ray maps of the other scans positive
for the same finding. The pre-specified orientation check used the frames alone, never a heat map: the CT attenuation frames correlated
better with the projections as they were (0.472) than mirrored (0.435), so no X-ray map was mirrored.

\subsection{Hallucination tests}
\label{sec:xm_halluc}
A classifier hallucinates when it reports a finding the image does not contain, or misses one it does. The tests use the
\num{9861} NIH test images without any finding and each model's seed-mean probabilities, with a per-class threshold at 90\,\% specificity
on the validation split. \emph{Finding-free perturbations:} rib emphasis (the radiograph plus half of the bone image of the suppressor of
Chapter~\ref{ch:v2}, which adds anatomy but no disease), Gaussian noise ($\sigma=0.02$) and a gamma of 0.8; the hallucination rate is the
share of images in which some class moves from at most to above its threshold. \emph{Counterfactual lesions:} a nodule from the lesion
library of Chapter~\ref{ch:audit} is pasted into a lung field of images whose Mass and Nodule scores are at most their thresholds; sensitivity
is the share in which either score crosses its threshold, and the score change is reported by nodule-size tertile. \emph{Off-target
attribution:} the share of Grad-CAM++ mass outside the lung field for positive predictions of lung findings. Seed groups of one model run on
different accounts were combined first (probabilities averaged with the number of seeds as weights, thresholds from the averaged
validation probabilities).

\subsection{Statistics}
\label{sec:xm_stats}
A model's prediction is the mean of its seeds' probabilities. H11 compares F$_\mathrm{p1}$ with F$_\mathrm{in}$ in the mean test AUROC of
the five exactly mapped shared findings; H12 makes the same comparison for F$_\mathrm{p1}$ against F$_\mathrm{p0}$ and for F$_\mathrm{p0}$
against F$_\mathrm{in}$; H13 compares F$_\mathrm{in}$ with the DenseNet-121 reference in the macro-AUROC over all 14 NIH diseases. All
three use a paired patient-cluster bootstrap (\num{10000} resamples) with two-sided percentile $p$-values (the $+1$ rule); H11 and the two
H12 contrasts form one Holm family. H14 is the mean per-map difference between matched and mismatched correlation with a scan-cluster
bootstrap, and H15 the mean per-image difference in the rib-emphasis flip between F$_\mathrm{p1}$ and F$_\mathrm{in}$ with a
patient-cluster bootstrap. The plan lists as descriptive: the CT teacher, the projection models and their zero-shot transfer to NIH and
VinBigData, per-class DeLong tests \citep{delong1988auc}, conformal risk control of the false-negative rate at 0.10
\citep{angelopoulos2022crc} for the Swin classifiers (calibrated on the NIH validation predictions as in Section~\ref{sec:cls_uq}) and the
pointing game. How the DeLong tests are grouped (Holm over the 15 shared-finding tests and over the 14 classes of the H13 contrast) and how
the pointing game is summarised (hit rate of the seed-mean hits with an image-cluster bootstrap interval) were fixed after the hypothesis
results were known and are logged as such.

\section{Results}
\label{sec:xm_results}
All 32 training runs finished (seven CT teachers, ten projection models and 15 NIH classifiers), as did the 20 zero-shot transfers, the
ten cross-modal map runs, seven hallucination test runs and 20 pointing-game runs. One zero-shot run first failed because it started
before its model had been written and was rerun unchanged. Table~\ref{tab:xm_ct_seeds} and Fig.~\ref{fig:xm_ct_training} show the seven
CT teachers.

\begin{table}[htbp]
\centering\small
\caption{Every CT teacher run (unchanged recipe, 20 epochs): best validation and test macro-AUROC over the six findings of the CT subset.
Slot $k$ is the teacher of P1 seed $k$; seed 2 collapsed and seed 5 replaced it by the rule fixed beforehand.}
\label{tab:xm_ct_seeds}
\fittable{tables/v4_ct_seeds.tex}
\end{table}

\begin{figure}[htbp]
\centering
\safefig[88mm]{figures/fig_v4_ct_training.pdf}
\caption{Validation macro-AUROC of the CT teacher per epoch for all seven seeds. Seed 2 never left chance level, seed 6 learned late and
was not used, and seed 0 stalled near 0.60 for 12 epochs before it learned.}
\label{fig:xm_ct_training}
\end{figure}

\subsection{CT teacher and projection models}
\label{sec:xm_res_transfer}
The CT teacher reached a seed-ensemble test macro-AUROC of 0.821 on the 297 CT subset test volumes (Table~\ref{tab:xm_transfer}). It was
strong for cardiomegaly (0.951) and pleural effusion (0.942) and weak for lung lesions (0.608). The recipe was unstable across seeds: the
single teachers ranged from 0.737 to 0.822 among the five used, and two of seven runs collapsed or nearly did.

On the projections, distillation lowered accuracy slightly. The seed ensemble of P1 reached 0.806 against 0.814 for P0, and P1 was below
P0 in every seed (Table~\ref{tab:xm_seeds}, Fig.~\ref{fig:xm_transfer}a). Without any training on real radiographs, the projection
models transferred moderately to VinBigData (P0 0.739, P1 0.732) and weakly to NIH (P0 0.648, P1 0.640); distillation again gave the
lower value in both.

\begin{table}[htbp]
\centering\small
\caption{CT teacher, projection models and their zero-shot transfer: test AUROC of the seed ensembles over the six findings, with the
macro-average (descriptive). NIH lung opacity is the approximate mapping to Infiltration.}
\label{tab:xm_transfer}
\fittable{tables/v4_transfer.tex}
\end{table}

\begin{table}[htbp]
\centering\small
\caption{Test macro-AUROC of every single model (one seed each); the other tables use the seed ensembles. Projection models and zero-shot
transfer: six findings; NIH classifiers: 14 diseases.}
\label{tab:xm_seeds}
\fittable{tables/v4_seeds.tex}
\end{table}

\begin{figure}[htbp]
\centering
\safefig{figures/fig_v4_transfer.pdf}
\caption{Macro-AUROC of single seeds (open circles) and seed ensembles (diamonds). (a) Projection test split; (b, c) zero-shot transfer
of the projection models to NIH and VinBigData over the six shared findings; (d) NIH classifiers over the 14 diseases.}
\label{fig:xm_transfer}
\end{figure}

\subsection{Distillation and real radiographs (H11, H12)}
\label{sec:xm_res_h11}
Distillation from CT did not improve the classification of real radiographs (Table~\ref{tab:xm_hyp}, Fig.~\ref{fig:xm_forest}a). On the
five shared findings F$_\mathrm{p1}$ differed from F$_\mathrm{in}$ by $-0.0001$ $[-0.0016, 0.0014]$ (Holm $p=0.912$), so a gain larger
than 0.0014 is not compatible with the data. F$_\mathrm{p1}$ did not differ from F$_\mathrm{p0}$ ($-0.0008$ $[-0.0022, 0.0006]$, Holm
$p=0.742$), nor F$_\mathrm{p0}$ from F$_\mathrm{in}$ ($+0.0007$ $[-0.0008, 0.0023]$, Holm $p=0.742$). Over all 14 diseases the three Swin
classifiers reached 0.819 (F$_\mathrm{in}$), 0.822 (F$_\mathrm{p0}$) and 0.821 (F$_\mathrm{p1}$) (Table~\ref{tab:xm_macro}). Per shared
finding no difference survived the correction; the largest were $+0.0044$ for cardiomegaly with F$_\mathrm{p0}$ against F$_\mathrm{in}$
(Holm $p=0.093$) and $-0.0042$ for consolidation with F$_\mathrm{p1}$ against F$_\mathrm{in}$ (Holm $p=0.122$; Table~\ref{tab:xm_shared}).

\begin{table}[htbp]
\centering\footnotesize
\caption{Pre-specified hypotheses: differences of the seed ensembles with 95\,\% bootstrap intervals (\num{10000} resamples). Holm: H11
with the two H12 contrasts. The P0 row of H14 is descriptive.}
\label{tab:xm_hyp}
\fittable{tables/v4_hypotheses.tex}
\end{table}

\begin{figure}[htbp]
\centering
\safefig{figures/fig_v4_forest.pdf}
\caption{Pre-specified contrasts with 95\,\% bootstrap intervals. (a) AUROC contrasts of H11--H13; (b) H14, matched minus mismatched
correlation of the CT and X-ray heat maps; (c) H15, difference in the share of finding-free images flipped by rib emphasis. Filled
markers: $p<0.05$ (Holm-adjusted for H11 and H12).}
\label{fig:xm_forest}
\end{figure}

\begin{table}[htbp]
\centering\small
\caption{NIH test macro-AUROC of the seed ensembles over the 14 diseases and mean AUROC of the five shared findings (descriptive).}
\label{tab:xm_macro}
\fittable{tables/v4_macro.tex}
\end{table}

\begin{table}[htbp]
\centering\footnotesize
\caption{Shared findings on the NIH test split: AUROC of the seed ensembles and paired DeLong differences, Holm over the 15 tests
(descriptive).}
\label{tab:xm_shared}
\fittable{tables/v4_shared_delong.tex}
\end{table}

\subsection{Swin against DenseNet-121 (H13)}
\label{sec:xm_res_h13}
The DenseNet-121 reference was better than F$_\mathrm{in}$ over the 14 diseases by $0.0049$ $[0.0015, 0.0085]$ ($p=0.005$; the difference
F$_\mathrm{in}$ minus DenseNet-121 is $-0.0049$). The per-class tests, which were not planned in this form, locate almost all of it in one
class (Table~\ref{tab:xm_per_class}, Fig.~\ref{fig:xm_per_class}): for hernia, the rarest NIH class (86 positive test images),
DenseNet-121 reached 0.901 against
0.833 ($-0.0679$, Holm $p=0.0008$), which alone accounts for $-0.0049$ of the 14-class mean. The only other class that survived the
correction was infiltration ($-0.0062$, Holm $p=0.003$); the remaining 12 classes differed by less than 0.007 in either direction. On the
classes other than hernia the two encoders were practically equivalent.

\begin{table}[htbp]
\centering\footnotesize
\caption{NIH per-class test AUROC of the seed ensembles, with the paired DeLong difference F$_\mathrm{in}$ minus DenseNet-121 and Holm over
the 14 classes (descriptive).}
\label{tab:xm_per_class}
\fittable{tables/v4_per_class.tex}
\end{table}

\begin{figure}[htbp]
\centering
\safefig{figures/fig_v4_per_class.pdf}
\caption{(a) NIH per-class test AUROC of the seed ensembles. (b) F$_\mathrm{in}$ minus DenseNet-121 per class with 95\,\% DeLong
intervals; filled markers: Holm $p<0.05$ over the 14 classes. Hernia carries almost all of the H13 difference.}
\label{fig:xm_per_class}
\end{figure}

\subsection{Cross-modal agreement (H14)}
\label{sec:xm_res_h14}
The CT and X-ray heat maps did not agree. Over 419 maps of 227 CT subset test scans, the matched CT and P1 maps correlated at $-0.052$ on
average and the mismatched ones at $-0.060$, a difference of $+0.0076$ $[-0.0068, 0.0223]$ ($p=0.301$; Table~\ref{tab:xm_hyp},
Fig.~\ref{fig:xm_forest}b). With P0 the matched maps correlated less than the mismatched ones ($-0.0087$ $[-0.0228, 0.0056]$). Per
finding all mean correlations were negative, from $-0.008$ to $-0.140$, and the matched-minus-mismatched differences changed sign from one
finding to the next (Table~\ref{tab:xm_crossmodal}, Fig.~\ref{fig:xm_crossmodal}). The mean maps show why (Fig.~\ref{fig:xm_maps}). The
projected CT maps vary almost only from top to bottom, in horizontal bands that follow the importance of individual slices, while the
X-ray maps of most findings concentrate along the edges of the lung box; for cardiomegaly they peak near the centre, where the heart is,
and for lung lesions they are broad and weighted to the upper zones. Maps with such different structure are unlikely to correlate
positively, whatever the patient.

\begin{table}[htbp]
\centering\small
\caption{Cross-modal agreement per finding: mean Pearson correlation of the projected CT map with the X-ray map of the same scan (matched)
and with the X-ray maps of the other scans positive for the same finding (mismatched), maps averaged over the five seed pairs.}
\label{tab:xm_crossmodal}
\fittable{tables/v4_crossmodal.tex}
\end{table}

\begin{figure}[htbp]
\centering
\safefig[88mm]{figures/fig_v4_crossmodal.pdf}
\caption{Mean correlation of the CT map with the X-ray map of the same scan (filled) and of other scans positive for the same finding
(open), for P1 and P0; in parentheses the number of maps.}
\label{fig:xm_crossmodal}
\end{figure}

\begin{figure}[htbp]
\centering
\safefig{figures/fig_v4_maps.pdf}
\caption{Mean heat maps over the CT subset test scans positive for each finding: the CT teacher's Grad-CAM++ summed along the
anteroposterior axis (top) and the Grad-CAM++ maps of P1 and P0 on the paired projections, all cropped to the lung box at $32\times32$ and
scaled to sum 1 before averaging; each panel has its own colour scale. Only averages over scans are shown.}
\label{fig:xm_maps}
\end{figure}

\subsection{Hallucination tests (H15)}
\label{sec:xm_res_h15}
Rib emphasis flipped 13.7\,\% of the finding-free NIH test images for F$_\mathrm{in}$ and 12.9\,\% for F$_\mathrm{p1}$, a difference of
$-0.0080$ $[-0.0159, -0.0003]$ ($p=0.044$; Table~\ref{tab:xm_hyp}, Fig.~\ref{fig:xm_forest}c). H15 is secondary and was not corrected for
multiplicity. The comparison it was built for does not single out distillation: F$_\mathrm{p0}$, trained on the same projections without
the teacher, flipped 12.7\,\% and DenseNet-121 12.8\,\% (Table~\ref{tab:xm_halluc}, Fig.~\ref{fig:xm_halluc}). The clearest difference
between the models came from noise: it flipped 78.1\,\% of the images for DenseNet-121 and 27.3--29.3\,\% for the Swin classifiers, while a
gamma change flipped about 6\,\% for all four. Inserted nodules were rarely detected (sensitivity 2.8--3.2\,\%); the score change grew
with nodule size, from practically zero for nodules of 4.0--6.2\,mm to between $+0.015$ and $+0.019$ for those of 8.5--37.6\,mm
(Table~\ref{tab:xm_nodule}). Between 62\,\% (DenseNet-121) and 73\,\% (F$_\mathrm{p0}$) of the Grad-CAM++ mass of positive lung-finding
predictions fell outside the lung fields.

\begin{table}[htbp]
\centering\small
\caption{Hallucination tests on the \num{9861} finding-free NIH test images: share with some class above its threshold on the clean image,
share flipped by each perturbation, inserted-nodule sensitivity and the share of Grad-CAM++ mass outside the lung fields.}
\label{tab:xm_halluc}
\fittable{tables/v4_halluc.tex}
\end{table}

\begin{table}[htbp]
\centering\small
\caption{Counterfactual nodules: mean change of the Mass/Nodule score after insertion, by nodule-size tertile.}
\label{tab:xm_nodule}
\fittable{tables/v4_nodule_size.tex}
\end{table}

\begin{figure}[htbp]
\centering
\safefig{figures/fig_v4_halluc.pdf}
\caption{Hallucination tests. (a) Share of finding-free NIH test images in which some class crosses its threshold after a perturbation
that adds no disease. (b) Mean change of the Mass/Nodule score after inserting a nodule, by size tertile. (c) Share of Grad-CAM++ mass
outside the lung fields for positive lung-finding predictions.}
\label{fig:xm_halluc}
\end{figure}

\subsection{Pointing game}
\label{sec:xm_res_pointing}
The Swin classifiers' Grad-CAM++ maps rarely pointed into the expert boxes: 0.185 of the 984 targets for F$_\mathrm{in}$, 0.151 for
F$_\mathrm{p0}$ and 0.154 for F$_\mathrm{p1}$, against 0.385 for DenseNet-121 (Table~\ref{tab:xm_pointing}, Fig.~\ref{fig:xm_pointing}).
The decoder attention (0.479--0.572) and Integrated Gradients (0.615--0.645) did much better and differed little between the models. All
maps hit focal boxes (masses and nodules) less often than the other boxes. Grad-CAM++ of the Swin models often spreads along the image
borders (Fig.~\ref{fig:xm_xai}), as in the cross-modal maps.

\begin{table}[htbp]
\centering\small
\caption{Pointing game on the 984 NIH expert boxes (164 focal: Mass, Nodule): hit rate averaged over the five seeds, with a 95\,\%
image-cluster bootstrap interval (descriptive).}
\label{tab:xm_pointing}
\fittable{tables/v4_pointing.tex}
\end{table}

\begin{figure}[htbp]
\centering
\safefig{figures/fig_v4_pointing.pdf}
\caption{Pointing game. (a) Hit rate over all 984 boxes with 95\,\% image-cluster bootstrap intervals. (b) Focal (filled) and other
(open) boxes.}
\label{fig:xm_pointing}
\end{figure}

\begin{figure}[htbp]
\centering
\safefig{figures/fig_v4_xai_examples.pdf}
\caption{Heat maps of the seed-0 models on NIH test images with expert boxes (white): the first image of each box class among the 24
targets whose maps were saved; these 24 contain no mass or nodule box.}
\label{fig:xm_xai}
\end{figure}

\subsection{Uncertainty}
\label{sec:xm_res_uq}
Conformal risk control held the false-negative rate below its target for all three Swin classifiers: calibrated on the \num{8634} NIH
validation images, the test rate was 0.068--0.070 against a nominal 0.10, at the price of large label sets (6.33--6.45 of the 14 classes
per image; Table~\ref{tab:xm_conformal}, Fig.~\ref{fig:xm_conformal}). Selective classification by ensemble entropy behaved identically
for the three (area under the risk-coverage curve 0.0358).

\begin{table}[htbp]
\centering\small
\caption{Conformal risk control of the per-image false-negative rate at 0.10 on the NIH test split and selective classification by
$1-$entropy (descriptive).}
\label{tab:xm_conformal}
\fittable{tables/v4_conformal.tex}
\end{table}

\begin{figure}[htbp]
\centering
\safefig{figures/fig_v4_conformal.pdf}
\caption{(a) Risk-coverage curves of the Swin classifiers (Brier score of the images kept, most confident first). (b) Test false-negative
rate under conformal risk control at 0.10 (dashed); numbers: mean label-set size.}
\label{fig:xm_conformal}
\end{figure}

\section{Discussion}
\label{sec:xm_discussion}
Distillation from a CT teacher did not make a better radiograph classifier. On the shared findings the CT-distilled classifier matched
the one started from ImageNet within $\pm0.0016$ AUROC, and it matched the classifier pretrained on the same projections without the
teacher; on the projections themselves and in zero-shot transfer distillation was slightly worse. With five seeds, seed ensembles and a
plan fixed before training, the null result is informative: the interval excludes the gains that would justify a CT branch for
radiograph classification with this design.

Three observations suggest why. First, the teacher had little to give. On the six findings it was hardly better on CT (0.821) than P0 was
on the projections of the same test patients (0.814), and it was weakest exactly where CT should help most, lung lesions (0.608). Second, the
recipe was fragile: one of seven teachers collapsed, another nearly did, and the five used ranged from 0.737 to 0.822, so P1 seeds learned
from teachers of very different quality. Third, whatever P1 gained on the projections had to survive the move to real radiographs; the
zero-shot results (0.640--0.648 on NIH) show how wide that gap is, and the ten epochs of NIH fine-tuning may simply overwrite the
initialisation, as the equal results of F$_\mathrm{p0}$ and F$_\mathrm{in}$ suggest.

The explanation half of the framework failed in an instructive way. H14 asked whether the two modalities attend to the same place; they do
not even produce maps of the same kind. The 2.5D teacher's Grad-CAM++ mostly ranks slices, so its projected map varies from top to bottom
but hardly across, and the last stage of SwinV2-T gives an $8\times8$ grid whose Grad-CAM++ spreads along the image borders. The pointing
game shows the same weakness on real radiographs: Grad-CAM++ hit the boxes less than half as often for Swin as for DenseNet-121, while
decoder attention and Integrated Gradients localised about as well as for the CNN. For Swin classifiers the latter two are the maps to
report.

The hallucination tests separate robustness from accuracy. The Swin classifiers were not more accurate than DenseNet-121, but they were
far less sensitive to image noise (flip rate about 0.28 against 0.78), a property the accuracy numbers do not show. Under rib emphasis the
CT-distilled classifier flipped fewer images than the ImageNet-started one ($p=0.044$, uncorrected), but the classifier pretrained on the
projections without the teacher did the same, so the effect, if real, belongs to pretraining on projections rather than to CT. No model
reacted to inserted nodules of the two smaller size tertiles (up to 8.5\,mm), and even the largest ones moved the scores by less than
0.02.

\paragraph{Limitations.} The CT recipe was seed-unstable, and the replacement of seed 2 by seed 5 followed a rule fixed beforehand but
used a teacher trained later. The CT subset is small (\num{1600} training volumes) and its labels come from reports, as do those of the
projections. The NIH test split had been used in the earlier chapters with other models; H13's localisation in hernia comes from the
per-class tests, which were planned only as descriptive, and hernia has few positive test images. The cross-modal agreement depends on
Grad-CAM++ for both modalities, the weakest of the three maps on real radiographs; agreement of Integrated Gradients or attention maps was
not part of the plan. The example heat maps come from the first 24 saved targets, which contain no mass or nodule. VinBigData fine-tuning
(about 50 GPU hours) and the design's idea of using CT--X-ray disagreement as an abstention signal were not carried out. The teacher
outputs, the predictions and the weights of the distilled models left the cloud accounts for the author's computer, against the data
rule of the design; this is logged, and nothing derived from CT-RATE was released.

\section{Summary}
\label{sec:xm_summary}
A Swin framework for chest CT and X-ray was built and tested under a pre-specified, hash-frozen plan. Distilling a CT teacher into an X-ray
model did not improve the classification of real radiographs (H11, H12), the Swin classifier matched DenseNet-121 except in one rare class
(H13), and the CT and X-ray heat maps did not agree (H14). The Swin classifiers were much more robust to noise than DenseNet-121. Under
rib emphasis the CT-distilled classifier flipped slightly fewer finding-free images than the one started from ImageNet (H15, uncorrected),
as did the classifier pretrained on the projections without the teacher. Grad-CAM++ is a poor explanation for Swin classifiers;
Integrated Gradients and decoder attention localise about as well as for a CNN.


\chapter{Discussion and limitations}\label{ch:discussion}
% Discussion chapter: the paper's discussion and limitations are input, then thesis-only material.
\section{Discussion}
\label{sec:discussion}
% Written from the final tables (results/final/stats_held.json, results/bdcxr/stats.json). Failed hypotheses are named first.
\paragraph{Principal findings.} Three pre-specified results did not hold: the soft-tissue image alone did not improve any downstream task
(Table~\ref{tab:downstream_cls}, Table~\ref{tab:downstream_det}), SSIM fell short of the reference model by more than its margin
(Table~\ref{tab:h2}), and conformal intervals calibrated on synthetic data covered 14\,\% instead of 90\,\% of real bone pixels
(Table~\ref{tab:reliability}); recalibrated on 24 real pairs in an exploratory analysis, they covered 89\,\% (Table~\ref{tab:bone_recal}). Real-nodule contrast was mixed: higher than with the reference model on JSRT, lower on Node21
(Table~\ref{tab:nodule_cnr}). The central results held. Lesion absorption fell steadily with the twin weight (Table~\ref{tab:dose_response}); the
twin ensemble absorbed a third of what the released CT-supervised model absorbed and less than a third of what the matched control absorbed, on
an audit set no model had seen during development (Table~\ref{tab:audit_contrasts}); and it did so at non-inferior bone-region fidelity. The effect
replicated with two other encoders (Table~\ref{tab:encoders}) and on paediatric radiographs from a population no model was trained on (Section~\ref{sec:bdcxr}).

\paragraph{What the audit adds.} Fidelity metrics and lesion absorption rank suppressors differently (Figure~\ref{fig:pareto}). The v1 model has
the highest bone-region PSNR of the CT-supervised models and the highest absorption; ours and the reference model are within 0.06\,dB of each
other and a factor of three apart in absorption. BS-LDM, trained on dual-energy pairs from the same hospital, is about 7\,dB ahead in fidelity on
that hospital's images, yet in the development audit it absorbed 0.42 of an inserted lesion (at site positions not paired with the other
methods, so only descriptive). A suppressor that matches the dual-energy target closely can still remove lesions: fidelity to that target does
not protect what the target never showed it. Absorption must nevertheless be read beside suppression strength. The reference model removes more
rib-edge energy than ours (48\,\% against 30\,\%) and more non-rib energy (22\,\% against 8\,\%), so part of its higher absorption may reflect
stronger suppression in general; ours absorbs a third as much while removing about 60\,\% of the rib-edge energy.

\paragraph{What the regulariser costs.} On real paired radiographs the twin costs no measurable bone-region fidelity up to $w=16$ (21.22 to
21.28\,dB across $w=0\ldots16$); at $w=32$ and 64 the PSNR drops by about 0.3\,dB, so $w=16$ is the usable
operating point, chosen on development data before the freeze. The cost appears where the truth is exact: on the Qure-projection test split the
bone error of the twin ensemble is 41\,\% higher than that of the control (MAE 0.0041 against 0.0029), and in the ablation the validation error rises
from 0.0033 to 0.0057. The regulariser therefore trades some synthetic bone accuracy for lesion invariance; on real images that trade was not
visible. The ablation shows what carries the effect: gradients through both twin branches and a consistency term confined to the pasted lesion.
Showing lesions during training without the twin term (CLIA) did not lower absorption.

\paragraph{Relation to other methods.} The reference model is a SwinV2-B at $1024$~px trained with a far larger sample budget; we compare its
released weights as a user would run them, and we cannot say how much of the absorption difference is due to our regulariser rather than to
training differences, except through the matched no-twin control, which absorbs slightly more than the reference model (0.178 against 0.149). Diffusion models trained on real
dual-energy pairs \citep{sun2026bsldm,sun2026debonedit} answer a different question: they need paired data from the target site, sample slowly
(about 2.5 minutes per image for BS-LDM on a T4 GPU) and were not designed for lesion invariance.

\paragraph{Advice for evaluation and use.} We suggest that a bone suppressor be audited for lesion absorption before it is used, with the audit
reported next to a suppression measure; that suppressed images be offered beside the original rather than instead of it, which is also the only
use in which they helped downstream models here; and that the ensemble spread be used as a warning flag rather than as a calibrated interval: it
rose 2.7-fold on paediatric radiographs, but intervals calibrated on synthetic data did not cover real errors, and recalibration needed real
pairs. Whether classifiers that combine the radiograph with the soft-tissue and bone images in separate streams gain from suppression is tested
in \clsstudy.

\paragraph{Future work.} Prediction at $1024$~px, combining CT-derived and real dual-energy supervision with the twin term, conformal
recalibration on real pairs from more sites, lesion types other than nodules (consolidation, as in the paediatric pneumonia set), paediatric training data, and a
reader study that measures detectability rather than signal.


\section{Limitations}
\label{sec:limitations}
% PLAN (about 0.5 page, one short paragraph per item). Items 1-9 are known now from DEVIATIONS.md, DECISIONS.md and SAP_v2.md; add what the final
% results reveal. State each limitation with its consequence for how the results must be read. Do not soften a limitation that changes a conclusion.
\begin{enumerate}
\item \emph{Synthetic supervision and domain gap.} Training uses rendered or translated projections; the bone image of a real radiograph is never
 observed. Real-image fidelity is measured against dual-energy soft-tissue images, which differ from the radiograph in tone and in high-frequency
 content, so raw and band-limited metrics are both reported and the calibrated metrics are not absolute.
\item \emph{One source of paired real data.} SZCH-X-Rays comes from one partner hospital and one acquisition set-up (dual-exposure DES on a GE Discovery XR656).
\item \emph{The audit is linear and additive.} Lesions are added in image units to a real radiograph; beam hardening, scatter and the tone mapping of
 the real image are not modelled, so $A$ measures the leak of an added signal and not clinical detectability. $A=0$ is reached by doing nothing, and
 $A$ is therefore read next to a suppression gain.
\item \emph{Shared insertion mechanism.} The regulariser and the audit both insert lesions. The confirmatory audit, the real-nodule contrast and the
 downstream subgroup analysis were added to limit this, but they do not remove it.
\item \emph{Development use of SZCH.} All 75 SZCH images were used in development (audit) and 24 for paired decisions; only the 51 held-out pairs are
 clean for fidelity claims, and the SZCH audit is development evidence only. The plan changed once after the SZCH pairs were looked at
 (SAP\_v2.md, DEVIATIONS.md).
\item \emph{Resolution and backbone.} Our models predict at 512 px with tiny encoders (ConvNeXt-V2-T for ours, SwinV2-T for the dose-response);
 the reference model predicts at 1024 px with SwinV2-B. The SSIM gap sits in fine texture: with a band limit of $\sigma=4$ it shrinks from 0.021 to
 0.002, and ours keeps about 9\,\% more high-frequency energy than the dual-energy target in non-bone regions (the reference model 1\,\% less). Whether
 prediction at 1024 px closes the remaining gap was not tested.
\item \emph{Operating point.} The twin weight trades absorption against fidelity and against bone accuracy on exact-truth data (Section~\ref{sec:discussion}).
 Only $w=16$ was evaluated with the full protocol; the dose-response shows that $w=32$ and 64 lower bone-region PSNR by about 0.3\,dB.
\item \emph{Uncertainty.} Conformal calibration treats pixels as exchangeable, so coverage is marginal over pixels, and it is calibrated on synthetic
 data; coverage on real images is reported as measured under shift. The recalibration on real pairs was chosen after the shift was seen and
 used 24 pairs from one hospital. Ensemble members share training data.
\item \emph{Scope.} Training and paired evaluation use adult frontal radiographs; paediatric images enter only the external validation, where no
 paired truth exists; no lateral images and no reader study. Research use only and not a medical device. The weights are non-commercial because
 the training data are.
\item \emph{Overlap maps.} On JSRT the bone overlap of each audit site came from the reference model's bone map, not from the consensus map the
 plan names, so the JSRT interaction is defined by a comparator; on VinDr-RibCXR it comes from radiologist rib masks and shows no interaction.
\item \emph{Downstream and external validation.} Downstream models were trained once per seed with fixed settings, not tuned per input; the
 subgroup analysis by bone overlap is not estimable on TBX11K, where every test box overlaps bone. BDCXR-3257 is a single-centre set with
 image-level labels and no patient identifiers, so repeated images of one child can only be caught by the near-duplicate grouping; its audit
 inserts adult nodules into children's radiographs.
\item \emph{Post-freeze corrections.} Several analysis errors found in an independent check after the first locked run were corrected
 (a resampling error that narrowed interaction intervals, incomplete downstream statistics, and detection metrics that depended on the order of
 tied scores by up to 0.00055); some evaluations ran on cloud GPUs whose floating-point results differ from the local ones by at most
 0.0024\,dB, with the audit sites pinned to the local lung masks. No pre-specified verdict changed; every change is listed in the deviations log.
\item \emph{Registration.} The analysis plans were hash-frozen in the project's version control before testing but not deposited in a public
 registry.
\end{enumerate}


\section{Threats to validity}
\paragraph{Internal validity.} The operating point, the training source and the encoder were chosen on development data: the 24 SZCH
development pairs, the audit on all 75 SZCH images and synthetic validation data. The confirmatory sets (VinDr-RibCXR, JSRT, the 51 held-out
SZCH pairs, the downstream test sets and BDCXR-3257) were evaluated only after the plans were frozen, and the screening tools that had printed
held-out SZCH numbers next to development numbers were changed and the episode disclosed. Seeds were fixed per arm, and every final contrast
pairs models on identical images and sites; where cloud and local GPUs were mixed, the audit sites were pinned to the same lung masks and the
remaining floating-point difference was measured (at most 0.0024\,dB). The selection of ConvNeXt-V2-T over RAD-DINO was made by a pre-declared
tie-break at two seeds and was not reopened when the third seeds moved the gap just outside the band.

\paragraph{External validity.} Paired real truth comes from one hospital and one dual-energy set-up. Training data are adult CT projections,
and the external validation covers one paediatric hospital. The audit uses nodules only; consolidations, masses, effusions and lines were not
inserted. Lateral, portable and supine radiographs were not studied.

\paragraph{Construct validity.} Absorption measures how much of an added, linear lesion signal moves into the bone estimate. It is not
clinical detectability, and a model that suppresses less absorbs less for that reason alone; it is therefore read with the rib benchmark and
bone-region fidelity. Real-nodule contrast and downstream performance were added because they do not use insertion, and they agree with the
audit against the matched control but not uniformly against the reference model.

\paragraph{Statistical conclusion validity.} Holm correction is applied to the $p$-values within each hypothesis family; confidence intervals
are unadjusted. Bootstrap intervals resample whole images, patients or near-duplicate groups as the plans require; per-class DeLong tests on
NIH ignore the clustering of images within patients, which the macro-level patient bootstrap does not. An independent check after the first
locked run found and fixed several analysis errors, among them a resampling error that had narrowed the interaction intervals and detection
metrics that depended on the order of tied scores; no pre-specified verdict changed, and every correction is listed in
Appendix~\ref{app:deviations}.

\section{Future work}
Five directions follow from the results. Prediction at $1024$~px would show whether the remaining SSIM gap to the reference model is a matter of
resolution. Combining CT-derived supervision with real dual-energy pairs, where available, could raise fidelity without giving up the twin
term, and real pairs from the target site would allow conformal calibration that holds on real images. The audit should be extended to
lesion types other than nodules, such as consolidation, which the paediatric pneumonia set points to. Paediatric training data would test
whether the higher absorption and uncertainty on children can be reduced. Finally, a reader study would measure what the audit cannot:
whether radiologists find more lesions on images from a suppressor with low absorption.


\chapter{Reproducibility and ethics}\label{ch:ethics}
% Reproducibility and ethics. Licence facts are copied from the pages listed in paper/REFS_TODO.md (read on 2026-11-03); unread terms are marked.
\section{Reproducibility, data and code availability}
\label{sec:repro}
% PLAN (about 0.5 page). Say exactly what is released, under which licence, and how every table and figure is regenerated.
%   R1 Pre-specified plans and freeze: SAP.md, SAP_v2.md, freeze.json, freeze_v2.json (hashes of plan, split lists, weights, code commit); dates.
%   R2 Generated outputs: all tables and figures come from saved result files; commands in RUNPLAN.md section 6 (tools/final_local.sh,
%      tools/final_stats.py, tools/fig_*.py).
%   R3 Randomness: seeds, paired sites drawn once with a fixed seed, bootstrap seeds.
%   R4 Compute: Kaggle T4 for training, a single consumer GPU for the local evaluation; wall-clock budget.
%   R5 Data access: which data are public, which are licensed (VinDr-RibCXR: metrics only, never redistributed), which links.
%   R6 Release: weights (three checkpoints) and code, licence, model card, demo; nothing derived from licensed images.
%   R7 Reporting checklist: CLAIM 2024 (tejani2024claim), Appendix~\ref{app:claim}.
The analysis plans (SAP.md and SAP\_v2.md), the split files, the weights and the code commit were hashed into \path{freeze.json} and
\path{freeze_v2.json} before any held-out or locked evaluation: \path{freeze.json} on 27 October 2026 (code commit \texttt{8e9e3be}) and
\path{freeze_v2.json} on 1 November 2026 (commit \texttt{dd52b6a}); the BDCXR-3257 addendum was committed on 3 November 2026 before any model
saw those images. The plans were kept in the project's version control, not deposited in a public registry; they and the hash files are
supplementary material. Every table and figure is produced by a script from saved result files; no number in the text is typed by hand. From the saved
results, \texttt{python} \path{tools/final_stats.py} writes \path{results/final/stats_held.json} and the main tables, \texttt{python} \path{tools/paper_extras.py} \texttt{--audit-map} the ablation table and the dose-response, audit-map, downstream and reliability figures, \texttt{python} \path{tools/fig_pareto.py} \texttt{final}
and \path{tools/fig_real_qualitative.py} the remaining figures, and \texttt{python} \path{tools/bdcxr_stats.py} the BDCXR-3257 tables; the
evaluations themselves are run by \path{tools/final_local.sh} (locally or in cloud kernels). \todo{repository URL and release commit}

Training used Kaggle T4 GPUs: a run of \num{300000} samples took 6 to 7 GPU-hours on one T4 (two runs share a two-GPU session), a screening run
of \num{100000} samples about 2 GPU-hours; the per-run times are in the training logs. The locked-set evaluation ran on a consumer GPU and on
cloud T4 GPUs with identical code; the audit sites of every model were pinned to the same lung masks, and model outputs differ between the two
GPU types by at most 0.0024\,dB.
The Qure.ai projections, CT-RATE, SZCH-X-Rays, JSRT, Node21, TBX11K and NIH ChestX-ray14 are obtained from their providers under their own terms.
VinDr-RibCXR is licensed and was processed only on the licence holder's own computer and cloud account; VinBigData was processed only in
cloud accounts that accepted the competition rules; BDCXR-3257 was used under a data use agreement. Only metrics left these environments, and no
image, mask, map or model derived from VinDr-RibCXR, VinBigData or BDCXR-3257 images is released.
The three checkpoints of the ensemble, the inference code and a model card are released for research use under CC BY-NC-SA 4.0, the licence of the
training projections. \todo{link, revision hash and checksum after release approval}
The reporting follows the CLAIM 2024 checklist \citep{tejani2024claim} (Appendix~\ref{app:claim}).


\section{Licences and terms of use}
Each data set was used under its own terms; the licence pages listed in the reference notes of the repository were read on 3 November 2026.
\begin{description}
\item[LIDC-IDRI] (The Cancer Imaging Archive) CC BY 3.0; the renders and the lesion library are derived from it, so any release that contains
 lesion patches carries the attribution and the data DOI.
\item[Qure.ai CT projections] CC BY-NC-SA 4.0 on the Hugging Face Hub. Weights trained on them inherit the non-commercial, share-alike terms.
\item[CT-RATE] Gated on the Hugging Face Hub, CC BY-NC-SA 4.0, no redistribution; it was not used directly, only through the projections.
\item[SZCH-X-Rays] CC BY-NC 4.0; its images may be shown with attribution, which is why every qualitative figure of this thesis uses them.
\item[VinDr-RibCXR] A signed data use agreement: non-commercial research only, no redistribution without permission, the original images and
 annotations are not to be modified, and the data set page is cited. The images were processed only on the licence holder's computer and
 private cloud account; lesions were added only in memory for the audit and no modified image was stored or shared; only metrics left these
 environments.
\item[VinBigData] The Kaggle competition built on the VinDr-CXR images; its rules allow non-commercial research and academic use. It was
 processed only in cloud accounts that had accepted the rules.
\item[JSRT] Free of charge after registration with the Japanese Society of Radiological Technology, but not in the public domain: the
 archived request form forbids distributing images and commercial use. No JSRT image is shown or redistributed.
\item[Node21] CC BY-NC-ND 4.0: modified Node21 images may not be redistributed, and none is.
\item[TBX11K] Data under CC BY 4.0 (code under CC BY-NC-SA 4.0).
\item[NIH ChestX-ray14] No named licence; use is unrestricted provided the download site is linked, the NIH Clinical Center acknowledged and the
 data set paper cited \citep{wang2017chestxray}.
\item[BDCXR-3257] Published on figshare under CC BY 4.0 \citep{durjoy2026bdcxr}; this thesis follows the stricter data use agreement between
 the hospital and the author's research group, which allows academic use and publication of derived results but not redistribution or public
 release of the radiographs.
\end{description}
Pretrained encoders carry their own terms, read on their model pages: the ConvNeXt V2 weights used for the final model are CC BY-NC 4.0;
RAD-DINO is MIT-licensed and described by its authors as shared for research only and not for clinical practice; the DINOv3 weights fall under
the DINOv3 License, which requires derivatives to carry the same agreement and publications to acknowledge DINOv3. A release of the final
ensemble has to respect both the CC BY-NC-SA 4.0 terms of the training projections and the CC BY-NC 4.0 terms of the ConvNeXt V2 weights.

\section{Data protection and privacy}
All data are retrospective and were de-identified by their providers. No attempt was made to re-identify any person, and no data set was
linked to another at the level of individuals. Licensed and restricted data (VinDr-RibCXR, VinBigData, BDCXR-3257) were processed only by
their licence holders or within accounts permitted to hold them, and only aggregate metrics, predictions and statistics left those
environments. File names and per-image results of these sets are not released. No credential was ever written into a file of the repository,
a notebook or a log; access tokens were kept in local files outside the repository and are to be rotated before any release.
\todo{state the data protection law and institutional policy that apply to the author's institution}

\section{Ethical considerations}
This is a retrospective study of existing images and did not involve new patients. BDCXR-3257 was used under the ethical approval
REC-FSIT-2025/No:~11928 of Daffodil International University and the data use agreement with the providing hospital; the public and licensed
data sets were used under their terms. \todo{institutional ethics statement for the thesis as a whole: approval or exemption}

The main ethical risk of the subject itself is the one the thesis measures: a bone suppressor that removes lesions can hide disease from a
reader who trusts the suppressed image. The models of this thesis are for research only and are not medical devices; they were not evaluated
for clinical use, by readers, or on populations beyond those described. Their failures are reported with the same weight as their successes,
including the cases where the reference model did better and the conditions under which the uncertainty estimates do not hold. A suppressed image
should be shown beside the original, never instead of it.

\section{Release plan}
Nothing has been published yet; release requires the author's explicit decision. The planned release consists of the three checkpoints of the
final ensemble, the inference code, a model card that states the training data, the intended use, the measured limitations and the licence,
and optionally a demonstration that processes uploaded images in memory and stores nothing. Training, evaluation and table scripts, the
analysis plans, the decision and deviation logs and the freeze files can be released because they contain no data. Nothing derived from
VinDr-RibCXR, VinBigData or BDCXR-3257 images, and no real radiograph other than SZCH-X-Rays examples with their licence line, will be released.
Before publication the checkpoints are inspected for embedded paths or names, their SHA-256 hashes are written into the model card, and all
access tokens used during the project are rotated. The checklist is kept in \path{release/README.md}.


\chapter{Conclusion}\label{ch:conclusion}
% Conclusion: no new results. Each answer cites the table that carries it.
\section{Summary of findings}
\paragraph{RQ1: is exact CT supervision enough?} No. On rendered LIDC-IDRI data with exact truth, lesion insertion lowered absorption a little
at no cost in bone error, and the pre-specified hypothesis H1 passed (Table~\ref{tab:v1_h1}). On real radiographs, however, the v1 models
absorbed more than twice what the reference model absorbed, and on the locked VinDr-RibCXR audit the v1 model still absorbs 0.289 of an added
lesion. Exact supervision of the bone image does not by itself keep lesions out of it.

\paragraph{RQ2: does the twin regulariser help without losing fidelity?} Yes, for bone-region fidelity. Absorption falls with the twin weight
(H1$'$, Table~\ref{tab:dose_response}); the final ensemble absorbs 0.051 of an added lesion on VinDr-RibCXR against 0.149 for the released
reference model and 0.178 for the matched control (H3$'$, Table~\ref{tab:audit_contrasts}), while its bone-region PSNR is non-inferior to both
(H2$'$, Table~\ref{tab:h2}). Structural similarity to the dual-energy image missed its margin, mainly through fine texture, and on exact-truth
data the regulariser raises the bone error. The ablation shows that the effect needs a consistency term confined to the lesion and gradients
through both twins (Table~\ref{tab:ablation_twin}), and it appears with every encoder tried (Table~\ref{tab:encoders}).

\paragraph{RQ3: does the audit measure something real?} In part. Against the matched control, real-nodule contrast rose consistently with the
regulariser on both nodule data sets, which agrees with the audit. Against the reference model the answer depends on the data set: higher
contrast on JSRT, lower on Node21, where the reference model suppresses more strongly (Table~\ref{tab:nodule_cnr}). The audit and the evidence
without insertion agree on the effect of the regulariser; they do not agree on a single ranking of different models.

\paragraph{RQ4: does it matter downstream, and is the output reliable?} The soft-tissue image used alone did not improve classification or
detection, and adding it to the original image gave small gains (Tables~\ref{tab:downstream_cls} and~\ref{tab:downstream_det}). The
ensemble spread ranks the bone error well on rendered data, and conformal intervals reach their nominal coverage there, but neither holds on
real radiographs (Table~\ref{tab:reliability}); an exploratory recalibration on 24 real pairs restored the coverage of the intervals
(Table~\ref{tab:bone_recal}). On paediatric radiographs the suppressor was non-inferior for pneumonia classification, kept its absorption
ordering, and its uncertainty rose 2.7-fold (Chapter~\ref{ch:external}).

\paragraph{RQ5: do classifiers gain from the suppressor's images?} Not overall. A multi-stream classifier that read the radiograph,
soft-tissue and bone images did not improve macro-AUROC on any of three data sets, and it did not beat a control of the same capacity or
channel stacking (H6 and H7, Table~\ref{tab:cls_primary}). The largest effect on both adult data sets was a small gain for nodules and
masses, significant on NIH ChestX-ray14 only and offset there by a loss for diffuse findings (H8, Table~\ref{tab:cls_groups}); attribution
maps of the soft-tissue stream located nodules and masses more often than those of the radiograph stream (H10,
Table~\ref{tab:cls_pointing}). The suppressor's own uncertainty did not help a classifier decide when to abstain
(Table~\ref{tab:cls_selective}).

\paragraph{RQ6: does CT teach a Swin X-ray classifier?} Not with this design. Distilling a 2.5D Swin CT classifier into the X-ray model
left the classification of real radiographs unchanged: on the five shared findings the distilled classifier differed from one started from
ImageNet by $-0.0001$ $[-0.0016, 0.0014]$ in mean AUROC (H11), and neither pretraining on the projections nor the teacher on top of it made
a difference (H12, Table~\ref{tab:xm_hyp}); on the projections and in zero-shot transfer the distilled models were slightly worse
(Table~\ref{tab:xm_transfer}). The CT and X-ray heat maps of the same patient did not agree (H14, Table~\ref{tab:xm_crossmodal}). The Swin
classifier was slightly less accurate than the DenseNet-121 of Chapter~\ref{ch:classifier} (H13), almost entirely because of hernia; per
class it lost clearly on hernia and slightly on infiltration and matched it on the other 12 diseases (Table~\ref{tab:xm_per_class}). It was far less sensitive to image noise than DenseNet-121, and under rib emphasis the distilled classifier
flipped slightly fewer finding-free images than the one started from ImageNet (H15, $p=0.044$ uncorrected), as did the classifier pretrained
on the projections without the teacher (Table~\ref{tab:xm_halluc}). Grad-CAM++ of the Swin classifiers located expert boxes less than half
as often as that of DenseNet-121, while Integrated Gradients and decoder attention did about as well (Table~\ref{tab:xm_pointing}).

\section{Contributions}
The thesis contributes a counterfactual lesion audit that makes lesion loss by bone suppressors measurable on real radiographs, a twin
lesion-invariance regulariser that lowers it, a recipe that trains such suppressors on public data only, two pre-specified evaluations
on locked data with downstream tasks, an external validation and a full record of deviations and negative results, and a pre-specified test,
against a control of the same capacity, of whether classifiers gain from bone-suppressed inputs. It also contributes an explainable,
hallucination-aware Swin framework for chest CT and X-ray, with CT-to-X-ray distillation, a test of heat-map agreement between the
modalities and counterfactual hallucination tests, evaluated under a fourth pre-specified plan; its negative results bound what such a
transfer can add with public data of this size.

\section{Recommendations}
Whoever evaluates a bone suppressor should report how much of an inserted lesion it absorbs, on real radiographs and at sites shared by all
compared models, next to a measure of how much bone it removes, and not rely on PSNR and SSIM alone. Whoever uses one should show the
suppressed image beside the original rather than instead of it, should treat the ensemble spread as a warning that the image is unusual rather
than as a calibrated error bar, and should not use a suppressor on a population or image type it was not evaluated on. Whoever builds a
classifier should not expect bone-suppressed inputs to raise its overall accuracy, and should test any benefit per finding type against a
control of the same capacity. Whoever explains a Swin classifier of radiographs should prefer Integrated Gradients or decoder attention to
Grad-CAM++ and should score the maps against expert boxes; whoever claims that a classifier is robust should test it with perturbations that
add no disease, because accuracy did not reveal the nearly threefold difference in noise sensitivity between two encoders of almost equal
accuracy.

\section{Outlook}
The audit applies to any image-to-image model that is meant to remove one structure and keep another. The next steps are a reader study, other
lesion types, higher resolution, real dual-energy data for calibration, and paediatric training data. For the cross-modality framework, a
stable and stronger CT teacher comes first: the 2.5D recipe collapsed in one of seven seeds and was weak for lung lesions, so a full 3D
encoder and many more volumes are needed before distillation can be judged on its merits. Heat-map agreement should then be measured with
Integrated Gradients or attention maps rather than Grad-CAM++, and disagreement between the CT and X-ray models could be tested as a signal
for when to abstain.


\appendix
\chapter{Deviations from the plan}\label{app:deviations}
% Follows DEVIATIONS.md entry by entry (dates and substance copied); each item says what changed, why, and how it changes the reading.
The entries below follow \path{DEVIATIONS.md} of the repository in date order. Each says what changed, why, and how it affects the reading of
the results.
\begin{description}
\item[2026-10-26, physics.] The renderer uses a water-equivalent density with per-component ICRU-44 attenuation, an energy-integrating
 detector and a W-anode spectrum; scatter, noise and blur are applied at training time rather than rendered. \emph{Effect:} the v1 renders are a
 simplified physical model; they concern the v1 study only.
\item[2026-10-26, model size.] SwinV2-Tiny at $512$~px instead of the reference SwinV2-Base at $1024$~px, for compute reasons; the nested
 decoder covers the encoder stages with a plain tail to full resolution. \emph{Effect:} comparisons with the reference model mix the
 regulariser with resolution and capacity; the matched no-twin control isolates the regulariser.
\item[2026-10-26, LIDC sources.] Part of the series came from a Kaggle copy, the rest from LUNA16 (checked voxel-identical to DICOM) or the TCIA
 API; the frame is a body crop with a bounded aspect ratio. \emph{Effect:} none on the results; recorded for reproducibility.
\item[2026-10-26, MIMIC-CXR and Node21.] MIMIC-CXR was dropped (not needed, credentialed access). Node21 was not available as a ready copy;
 it was later fetched from its Zenodo record inside a cloud kernel. \emph{Effect:} none.
\item[2026-10-27, usable CTs.] Short scans needed many padded frame rows; the limit was raised to 50 rows and padded rows are set to air.
 \emph{Effect:} more CTs usable for v1; no locked set involved.
\item[2026-10-27, learning rate.] The pilot passed its gate and a check on synthetic validation CTs kept $2\times10^{-4}$ for every arm; no real
 image was used. \emph{Effect:} none.
\item[2026-10-27, SZCH looked at.] The pre-specified comparison with the reference model was looked at on all 75 SZCH pairs after the v1
 ablation. SZCH became a development set (24 development, 51 held-out images) and its audit development evidence only; the confirmatory audit
 moved to VinDr-RibCXR and JSRT. \emph{Effect:} the main deviation of the project; it is why the second plan exists and why only the 51 held-out
 pairs support fidelity claims.
\item[2026-10-27, secondary metrics.] Band-limited SSIM, the fidelity of the removed signal and a high-frequency energy ratio were added after
 the SSIM gap was seen. \emph{Effect:} they explain the gap; they do not replace the pre-specified metrics.
\item[2026-10-27, Qure projections.] The public Qure.ai projections became the v2 training source after their exact additivity was verified;
 split by patient. \emph{Effect:} the weights inherit CC BY-NC-SA 4.0.
\item[2026-10-29, second plan.] SAP\_v2.md was written after the v2 recipe had been developed on development data; held-out SZCH PSNR and SSIM of
 screening runs had been printed by a screening tool. \emph{Effect:} disclosed; no decision used them, and the tool was changed.
\item[2026-10-29, encoder screen and ablation arms.] Added as exploratory, with an encoder learning rate of $0.25\times$ for vision
 transformers fixed in advance; ConvNeXt-V2-T diverged at the full rate and was rerun at $0.25\times$. \emph{Effect:} the rerun extends a rule
 written for another encoder type.
\item[2026-10-29, accidental start.] A dry-run listing started the held-out uncertainty evaluation by mistake; it was stopped after two minutes
 before writing any output. \emph{Effect:} none; the driver now guards that step.
\item[2026-10-31, BS-LDM.] The released BS-LDM weights were run as an extra baseline outside the analysis plan, with a seeded sampler and our SZCH
 preprocessing. \emph{Effect:} descriptive only; its audit sites were drawn before the masks were pinned and are not paired with the others.
\item[2026-10-31, replacement and tie-break.] Rule R2 first promoted RAD-DINO, then the pre-declared tie-break between the two encoders that met
 R2 chose ConvNeXt-V2-T at two seeds each; the matched no-twin control is the same encoder. With three seeds the gap moved 0.0005 outside the
 tie band; the choice was not reopened. \emph{Effect:} both encoders are reported on the locked sets; the selection used no locked number.
\item[2026-11-01, stored weights.] Seven of 27 evaluated checkpoints carry overflowed (infinite) batch-normalisation variances from mixed-precision
 training in one decoder node. \emph{Effect:} the results describe the stored models as they are; cloud and local outputs agree to 0.0024\,dB.
\item[2026-11-03, cloud evaluation and mask pinning.] The last locked-set outputs were computed on cloud GPUs. Because lung masks differ between
 GPUs, an optional mask cache was added to the audit and rib-benchmark code so that every model is scored on the sites of the final tables; the
 sites were checked to be identical file by file. \emph{Effect:} a post-freeze code change that leaves the computation unchanged when the cache is
 not used.
\item[2026-11-03, synthetic-bank path.] A wrong path in the evaluation driver would have loaded no CT for the synthetic evaluation; fixed before
 it ran. \emph{Effect:} none on results.
\item[2026-11-03, BDCXR-3257 duplicates.] The planned hash rule chained most of the paediatric images into one group; replaced, before any model
 ran, by a correlation rule. The 2\,\% lung-mask rule of the audit is applied to the eroded mask (no image affected). \emph{Effect:} the folds
 are balanced; repeated images of one child cannot be detected without identifiers.
\item[2026-11-03, independent recheck.] Two independent reviews of the statistics and the code found and fixed: an interaction bootstrap that
 dropped repeated draws (intervals about 25\,\% too narrow; one exploratory reading changed), an incomplete implementation of the downstream
 analysis (now with the planned bootstrap intervals), detection metrics that depended on the order of tied scores (up to 0.00055), missing Holm
 correction for the real-nodule contrasts, missing intervals for the reliability results (rerun identical), and the overlap source of the JSRT
 audit (the reference model's map instead of the consensus map). The multiplicity rule was stated explicitly. \emph{Effect:} no
 pre-specified verdict changed; all corrected numbers are those of this thesis.
\item[2026-11-03, v1 H1 record.] A recheck entry had described the v1 H1 test as missing; it had been computed on 27 October 2026 on the test CTs
 with all three lesion-aware arms. The analysis of the 300k finals is a replication. \emph{Effect:} both are reported (Table~\ref{tab:v1_h1}).
\item[3--7 November, v3.] \todo{summarise the v3 entries of DEVIATIONS.md (3 to 7 November) here}
\item[2026-11-08, CT micro-batches.] Gradient accumulation was added to the CT stage after the freeze so that four volumes of 24 slices fit a
 T4 GPU; each step still averages the same four volumes and the network has no batch statistics, so the step is unchanged. \emph{Effect:}
 none on the computation; the frozen hash of the CT code no longer matches the file.
\item[2026-11-08, accounts.] At the author's request one cloud account took the place of another; it prepared its own NIH copy, which was
 compared file by file with the common copy before use. \emph{Effect:} none.
\item[2026-11-09, CT teacher outputs and seed 2.] The design keeps everything derived from CT-RATE in the cloud except metrics and heat-map
 frames, but its own account plan needs the teacher outputs (logits and pooled features of the subset volumes, no images) on other accounts,
 which is possible only through a download; the standard download also brought the teachers' weights and predictions to the author's
 computer. Teacher seed 2 collapsed (constant logits; no code or data fault found). Seeds 5 and 6 were trained with the unchanged recipe
 under a rule fixed before any P1 model of seed 2 (Section~\ref{sec:xm_models}), and seed 5 took slot 2. \emph{Effect:} a disclosed
 conflict with the design's data rule (nothing was released); the CT recipe is seed-unstable (Section~\ref{sec:xm_discussion}).
\item[2026-11-09, last kernels.] The weights of the distilled projection models moved through the author's computer to the only account
 holding CT-RATE, for the cross-modal maps and the VinBigData zero-shot transfer (the same conflict as above). The DenseNet-121 reference runs
 were copied to an account with the common NIH copy, because one account's soft image differed and the rib-emphasis test reads it. The
 pointing-game launcher gained a run filter that leaves the v3 behaviour unchanged. A zero-shot run that failed because it started before its
 model was written was rerun unchanged. VinBigData fine-tuning, planned only if the quota allowed, was not run (about 50 GPU hours).
 \emph{Effect:} none on the hypotheses; one descriptive analysis is missing.
\item[2026-11-10, collection and descriptive summaries.] A rerun that sat in the cloud queue for nine hours was resubmitted unchanged. CT
 teacher seed 5's predictions, first downloaded without them, were fetched like those of the other seeds. The pointing-game summary, the
 grouping of the per-class DeLong tests and the conformal analysis were fixed after the hypothesis results were known
 (Section~\ref{sec:xm_stats}). \emph{Effect:} these descriptive results are post hoc in form; no hypothesis test was affected.
\end{description}


\chapter{Negative and null results}\label{app:negative}
% Negative and null results, in neutral words. Numbers are copied from the generated tables, DECISIONS.md (development data) and
% results/final/stats_held.json.
Every run or question that did not give the expected benefit, or any benefit, is listed here.
\begin{description}
\item[N1] \emph{Lesion-weighted loss in v1.} It lowered absorption on synthetic validation CTs but sat on the bone-error margin, and combined
 with insertion it missed the margin by about 6\,\% (Table~\ref{tab:v1_ablation}); it was therefore not selected, although on the test CTs it
 was non-inferior (Table~\ref{tab:v1_h1}).
\item[N2] \emph{Lesion insertion on the Qure source.} Adding counterfactual insertion to the Qure-projection source did not lower real-image
 absorption in the screening (0.218 against 0.216) or in the ablation at the full budget (0.191 against 0.190; Table~\ref{tab:ablation_twin}).
\item[N3] \emph{Mixed sources.} Mixing LIDC renders with the Qure projections gave higher absorption than the Qure source alone (0.286, and 0.269
 with insertion, against 0.216 in round 1 of the screening).
\item[N4] \emph{Backbone size.} SwinV2-B instead of SwinV2-T changed absorption by about one seed spread (0.189 against 0.216 without the twin
 term, 0.063 against 0.073--0.103 with it, both at the screening budget) and did not remove the fidelity cost of large twin weights; it was not pursued.
\item[N5] \emph{Training length.} Three times more samples changed absorption and fidelity within the seed spread (for example 0.071 against
 a mean of 0.086 at $w=16$), but made the seeds agree more closely.
\item[N6] \emph{Large twin weights.} At $w=64$ two seeds disagreed (absorption 0.020 and 0.061, development PSNR 21.27 and 20.80\,dB); on the
 locked sets $w=32$ and 64 lower bone-region PSNR by about 0.3\,dB (Table~\ref{tab:dose_response}).
\item[N7] \emph{SSIM gap to the reference model.} Not closed: SSIM$_{cal}$ misses its margin by 0.013 at the lower bound (Table~\ref{tab:h2}).
 With identical smoothing of all outputs the gap halved on development data, and with a band limit of $\sigma=4$ it nearly vanishes on the held-out
 pairs (0.946 against 0.948); the rest is attributed to fine texture and possibly to the lower output resolution, which was not tested.
\item[N8] \emph{Encoders not promoted.} DINOv3 ViT-S/16 did not improve on SwinV2-T at the screening budget; RAD-DINO missed rule R1 on PSNR by
 0.05\,dB (it later met R2 and is reported as a secondary encoder); ConvNeXt-V2-T diverged at the full encoder learning rate
 (Table~\ref{tab:app_screen}).
\item[N9] \emph{Failed or mixed confirmatory results.}
 SSIM non-inferiority against the reference model (H2$'$) failed.
 Real-nodule contrast on Node21 rose less with ours than with the reference model ($-0.025$ $[-0.033,-0.018]$; H4$'$a).
 The soft-tissue image used alone did not improve downstream classification or detection (H4$'$b): NIH macro-AUROC $-0.001$
 $[-0.004,0.003]$, VinBigData CPM $-0.009$ $[-0.016,-0.002]$; no difference in differences by bone overlap excluded zero, and on TBX11K it could
 not be estimated.
 Conformal intervals calibrated on synthetic data covered 0.139 of real bone pixels at a nominal 0.90 (H5$'$).
 On BDCXR-3257 the soft image was non-inferior but not superior for pneumonia classification.
\item[N10] \emph{CT-to-X-ray distillation (v4).} Distilling a CT teacher into the X-ray model did not improve the NIH classifier: on the
 shared findings F$_\mathrm{p1}$ minus F$_\mathrm{in}$ was $-0.0001$ $[-0.0016,0.0014]$ (H11), and neither H12 contrast differed
 (Table~\ref{tab:xm_hyp}). On the projections and in zero-shot transfer the distilled models were slightly worse
 (Table~\ref{tab:xm_transfer}).
\item[N11] \emph{Cross-modal heat-map agreement (v4).} The CT and X-ray maps of the same scan correlated no better than those of different
 scans (H14: $+0.0076$ $[-0.0068,0.0223]$), and all mean correlations were negative (Table~\ref{tab:xm_crossmodal}).
\item[N12] \emph{Swin against DenseNet-121 (v4).} The Swin classifier was not better than DenseNet-121 (H13: $-0.0049$
 $[-0.0085,-0.0015]$), mostly because of hernia (Table~\ref{tab:xm_per_class}), and its Grad-CAM++ maps hit the NIH boxes less than half as
 often (Table~\ref{tab:xm_pointing}).
\item[N13] \emph{CT teacher stability (v4).} One of seven CT teacher runs collapsed and another learned only in its last epochs
 (Table~\ref{tab:xm_ct_seeds}).
\end{description}


\chapter{Additional tables}\label{app:extra}
\input{sections/appendix_tables}

\chapter{Classifier study: deviations and additional figures}\label{app:classifier}
\input{sections/cls_appendix}

\chapter{CLAIM 2024 checklist}\label{app:claim}
\input{sections/claim_checklist}

\bibliographystyle{plainnat}
\bibliography{../paper/refs}

\end{document}
